\section{Schémas formels}
\label{section:I.10-fr}

\subsection{Schémas formels affines.}
\label{subsection:I.10.1-fr}

\begin{env}[10.1.1]
\label{I.10.1.1-fr}
Soit $A$ un anneau topologique \emph{admissible}
(\hyperref[0.7.1.2-fr]{0, 7.1.2})~; pour tout idéal de définition
$\mathfrak{J}$ de $A$, $\operatorname{Spec}(A/\mathfrak{J})$ s'identifie
au sous-espace fermé $V(\mathfrak{J})$ de $\operatorname{Spec}(A)$
(\hyperref[I.1.1.11-fr]{1.1.11}), ensemble des idéaux premiers
\emph{ouverts} de $A$~; cet espace topologique ne dépend pas de
\oldpage[I]{181}
l'idéal de définition $\mathfrak{J}$ considéré~; notons-le
$\mathfrak{X}$. Soit $(\mathfrak{J}_\lambda)$ un système fondamental de
voisinages de 0 dans $A$, formé d'idéaux de définition, et pour tout
$\lambda$, soit $\mathscr{O}_\lambda$ le faisceau structural de
$\operatorname{Spec}(A/\mathfrak{J}_\lambda)$~; ce faisceau est induit sur
$\mathfrak{X}$ par
$\widetilde{A}/\widetilde{\mathfrak{J}}_\lambda$ (lequel est nul hors de
$\mathfrak{X}$).

Pour $\mathfrak{J}_\mu\subset\mathfrak{J}_\lambda$, l'homomorphisme
canonique $A/\mathfrak{J}_\mu\to A/\mathfrak{J}_\lambda$ définit donc un
homomorphisme
$u_{\lambda\mu}:\mathscr{O}_\mu\to\mathscr{O}_\lambda$ de faisceaux
d'anneaux (\hyperref[I.1.6.1-fr]{1.6.1}), et
$(\mathscr{O}_\lambda)$ est un \emph{système projectif de faisceaux
d'anneaux} pour ces homomorphismes. Comme la topologie de $\mathfrak{X}$
admet une base formée d'ouverts quasi-compacts, on peut associer à tout
$\mathscr{O}_\lambda$ un \emph{faisceau d'anneaux topologiques
pseudo-discrets} (\hyperref[0.3.8.1-fr]{0, 3.8.1}) qui a
$\mathscr{O}_\lambda$ comme faisceau d'anneaux (sans topologie) sous-jacent,
et que nous noterons encore $\mathscr{O}_\lambda$~; et les
$\mathscr{O}_\lambda$ forment encore un \emph{système projectif de faisceaux
d'anneaux topologiques} (\hyperref[0.3.8.2-fr]{0, 3.8.2}). Nous désignerons
par $\mathscr{O}_{\mathfrak{X}}$ le \emph{faisceau d'anneaux topologiques}
sur $\mathfrak{X}$, limite projective du système
$(\mathscr{O}_\lambda)$~; pour tout ouvert \emph{quasi-compact} $U$ de
$\mathfrak{X}$, $\Gamma(U,\mathscr{O}_{\mathfrak{X}})$ est donc l'anneau
topologique limite projective du système d'anneaux \emph{discrets}
$\Gamma(U,\mathscr{O}_\lambda)$
(\hyperref[0.3.2.6-fr]{0, 3.2.6}).
\end{env}

\begin{definition}[10.1.2]
\label{I.10.1.2-fr}
Étant donné un anneau topologique admissible $A$, on appelle spectre formel
de $A$ et on note $\operatorname{Spf}(A)$ le sous-espace fermé
$\mathfrak{X}$ de $\operatorname{Spec}(A)$ formé des idéaux premiers
ouverts de $A$. On dit qu'un espace topologiquement annelé est un schéma
formel affine s'il est isomorphe à un spectre formel
$\operatorname{Spf}(A)=\mathfrak{X}$ muni du faisceau d'anneaux
topologiques $\mathscr{O}_{\mathfrak{X}}$ limite projective des faisceaux
d'anneaux pseudo-discrets
$(\widetilde{A}/\widetilde{\mathfrak{J}}_\lambda)|\mathfrak{X}$, où
$\mathfrak{J}_\lambda$ parcourt l'ensemble filtrant des idéaux de
définition de $A$.
\end{definition}

Quand nous parlerons d'un \emph{spectre formel}
$\mathfrak{X}=\operatorname{Spf}(A)$ comme d'un schéma formel affine, il
s'agira toujours de l'espace topologiquement annelé
$(\mathfrak{X},\mathscr{O}_{\mathfrak{X}})$ où
$\mathscr{O}_{\mathfrak{X}}$ est défini comme ci-dessus.

On notera que tout \emph{schéma affine} $X=\operatorname{Spec}(A)$ peut être
considéré comme un schéma formel affine d'une seule manière, en considérant
$A$ comme un anneau topologique discret~: les anneaux topologiques
$\Gamma(U,\mathscr{O}_X)$ sont alors discrets lorsque $U$ est
quasi-compact (mais non en général lorsque $U$ est un ouvert quelconque de
$X$).

\begin{proposition}[10.1.3]
\label{I.10.1.3-fr}
Si $\mathfrak{X}=\operatorname{Spf}(A)$, où $A$ est un anneau admissible,
$\Gamma(\mathfrak{X},\mathscr{O}_{\mathfrak{X}})$ est topologiquement
isomorphe à $A$.
\end{proposition}

\begin{proof}
En effet, comme $\mathfrak{X}$ est fermé dans $\operatorname{Spec}(A)$, il
est quasi-compact, et par suite
$\Gamma(\mathfrak{X},\mathscr{O}_{\mathfrak{X}})$ est topologiquement
isomorphe à la limite projective des anneaux discrets
$\Gamma(\mathfrak{X},\mathscr{O}_\lambda)$~; mais
$\Gamma(\mathfrak{X},\mathscr{O}_\lambda)$ est isomorphe à
$A/\mathfrak{J}_\lambda$ (\hyperref[I.1.3.7-fr]{1.3.7})~; comme $A$ est
séparé et complet, il est topologiquement isomorphe à
$\varprojlim A/\mathfrak{J}_\lambda$
(\hyperref[0.7.2.1-fr]{0, 7.2.1}), d'où la proposition.
\end{proof}

\begin{proposition}[10.1.4]
\label{I.10.1.4-fr}
Soient $A$ un anneau admissible, $\mathfrak{X}=\operatorname{Spf}(A)$, et
pour tout $f\in A$, soit
$\mathfrak{D}(f)=D(f)\cap\mathfrak{X}$~; l'espace topologiquement annelé
$(\mathfrak{D}(f),\mathscr{O}_{\mathfrak{X}}|\mathfrak{D}(f))$ est
isomorphe au spectre formel affine
$\operatorname{Spf}(A_{\{f\}})$
(\hyperref[0.7.6.15-fr]{0, 7.6.15}).
\end{proposition}

\begin{proof}
Pour tout idéal de définition $\mathfrak{J}$ de $A$, l'anneau discret
$S_f^{-1}A/S_f^{-1}\mathfrak{J}$ s'identifie canoniquement à
$A_{\{f\}}/\mathfrak{J}_{\{f\}}$
(\hyperref[0.7.6.9-fr]{0, 7.6.9}), donc
(\hyperref[I.1.2.5-fr]{1.2.5} et \hyperref[I.1.2.6-fr]{1.2.6}), l'espace
topologique $\operatorname{Spf}(A_{\{f\}})$ s'identifie canoniquement à
$\mathfrak{D}(f)$. En outre, pour tout ouvert quasi-compact $U$ de
$\mathfrak{X}$ contenu dans $\mathfrak{D}(f)$,
$\Gamma(U,\mathscr{O}_\lambda)$ s'identifie au module des sections du
faisceau structural de
$\operatorname{Spec}(S_f^{-1}A/S_f^{-1}\mathfrak{J}_\lambda)$ au-dessus de
$U$ (\hyperref[I.1.3.6-fr]{1.3.6}), donc, si on pose
$\mathfrak{Y}=\operatorname{Spf}(A_{\{f\}})$,
$\Gamma(U,\mathscr{O}_{\mathfrak{X}})$ s'identifie au module des sections
$\Gamma(U,\mathscr{O}_{\mathfrak{Y}})$, ce qui démontre la proposition.
\end{proof}

\oldpage[I]{182}
\begin{env}[10.1.5]
\label{I.10.1.5-fr}
En tant que faisceau d'anneaux \emph{sans topologie}, le faisceau
structural $\mathscr{O}_{\mathfrak{X}}$ de
$\operatorname{Spf}(A)$ admet pour tout $x\in\mathfrak{X}$ une fibre qui,
en vertu de (\hyperref[I.10.1.4-fr]{10.1.4}), s'identifie à la limite
inductive
$\varinjlim_{f\notin\mathfrak{j}_x}A_{\{f\}}$. Par suite
(\hyperref[0.7.6.17-fr]{0, 7.6.17} et
\hyperref[0.7.6.18-fr]{7.6.18})~:
\end{env}

\begin{proposition}[10.1.6]
\label{I.10.1.6-fr}
Pour tout $x\in\mathfrak{X}=\operatorname{Spf}(A)$, la fibre
$\mathscr{O}_x$ est un anneau local dont le corps résiduel est isomorphe à
$k(x)=A_x/\mathfrak{j}_xA_x$. Si en outre $A$ est adique et noethérien,
$\mathscr{O}_x$ est un anneau noethérien.
\end{proposition}

Comme $k(x)$ n'est pas réduit à 0, on conclut de ce résultat que le
\emph{support} du faisceau d'anneaux $\mathscr{O}_{\mathfrak{X}}$ est
\emph{égal à $\mathfrak{X}$}.

\subsection{Morphismes de schémas formels affines.}
\label{subsection:I.10.2-fr}

\begin{env}[10.2.1]
\label{I.10.2.1-fr}
Soient $A$, $B$ deux anneaux admissibles, et soit
$\varphi:B\to A$ un homomorphisme \emph{continu}. L'application continue
$ {}^a\varphi:\operatorname{Spec}(A)\to\operatorname{Spec}(B)$
(\hyperref[I.1.2.1-fr]{1.2.1}) applique alors
$\mathfrak{X}=\operatorname{Spf}(A)$ dans
$\mathfrak{Y}=\operatorname{Spf}(B)$, car l'image réciproque par
$\varphi$ d'un idéal premier ouvert de $A$ est un idéal premier ouvert de
$B$. D'autre part, pour tout $g\in B$, $\varphi$ définit un homomorphisme
continu
\[
  \Gamma(\mathfrak{D}(g),\mathscr{O}_{\mathfrak{Y}})
  \longrightarrow
  \Gamma(\mathfrak{D}(\varphi(g)),\mathscr{O}_{\mathfrak{X}})
\]
en vertu de (\hyperref[I.10.1.4-fr]{10.1.4}),
(\hyperref[I.10.1.3-fr]{10.1.3}) et
(\hyperref[0.7.6.7-fr]{0, 7.6.7})~; comme ces homomorphismes satisfont aux
conditions de compatibilité pour les restrictions correspondant au
passage de $g$ à un multiple de $g$, et que
$\mathfrak{D}(\varphi(g))={}^a\varphi^{-1}(\mathfrak{D}(g))$, ils
définissent un homomorphisme \emph{continu} de faisceaux d'anneaux
topologiques
$\mathscr{O}_{\mathfrak{Y}}\to
{}^a\varphi_*(\mathscr{O}_{\mathfrak{X}})$
(\hyperref[0.3.2.5-fr]{0, 3.2.5}), que nous noterons encore
$\widetilde{\varphi}$~; on a ainsi obtenu un morphisme
$\Phi=({}^a\varphi,\widetilde{\varphi})$ d'espaces topologiquement annelés
$\mathfrak{X}\to\mathfrak{Y}$. On notera qu'en tant qu'homomorphisme de
faisceaux d'anneaux sans topologie, $\widetilde{\varphi}$ définit un
homomorphisme
$\widetilde{\varphi}_x^\sharp:
\mathscr{O}_{{}^a\varphi(x)}\to\mathscr{O}_x$ sur les fibres, pour tout
$x\in\mathfrak{X}$.
\end{env}

\begin{proposition}[10.2.2]
\label{I.10.2.2-fr}
Soient $A$, $B$ deux anneaux topologiques admissibles, et soient
$\mathfrak{X}=\operatorname{Spf}(A)$,
$\mathfrak{Y}=\operatorname{Spf}(B)$. Pour qu'un morphisme
$u=(\psi,\theta):\mathfrak{X}\to\mathfrak{Y}$ d'espaces topologiquement
annelés soit de la forme $({}^a\varphi,\widetilde{\varphi})$, où
$\varphi$ est un homomorphisme continu d'anneaux $B\to A$, il faut et il
suffit que pour tout $x\in\mathfrak{X}$,
$\theta_x^\sharp$ soit un homomorphisme local
$\mathscr{O}_{\psi(x)}\to\mathscr{O}_x$.
\end{proposition}

\begin{proof}
La condition est nécessaire~: en effet, soit
$\mathfrak{p}=\mathfrak{j}_x\in\operatorname{Spf}(A)$, et soit
$\mathfrak{q}=\varphi^{-1}(\mathfrak{j}_x)$~; si
$g\notin\mathfrak{q}$, on a donc $\varphi(g)\notin\mathfrak{p}$, et il est
immédiat que l'homomorphisme
$B_{\{g\}}\to A_{\{\varphi(g)\}}$ déduit de $\varphi$
(\hyperref[0.7.6.7-fr]{0, 7.6.7}) transforme
$\mathfrak{q}_{\{g\}}$ en une partie de
$\mathfrak{p}_{\{\varphi(g)\}}$~; passant à la limite inductive, on voit
donc (compte tenu de (\hyperref[I.10.1.5-fr]{10.1.5}) et de
(\hyperref[0.7.6.17-fr]{0, 7.6.17})) que
$\widetilde{\varphi}_x^\sharp$ est un homomorphisme local.

Inversement, soit $(\psi,\theta)$ un morphisme vérifiant la condition de
l'énoncé~; en vertu de (\hyperref[I.10.1.3-fr]{10.1.3}), $\theta$ définit
un homomorphisme continu d'anneaux
\[
  \varphi=\Gamma(\theta):
  B=\Gamma(\mathfrak{Y},\mathscr{O}_{\mathfrak{Y}})
  \longrightarrow
  \Gamma(\mathfrak{X},\mathscr{O}_{\mathfrak{X}})=A.
\]
En vertu de l'hypothèse sur $\theta$, pour que la section $\varphi(g)$ de
$\mathscr{O}_{\mathfrak{X}}$ au-dessus de $\mathfrak{X}$ ait un germe
inversible au point $x$, il faut et il suffit que $g$ ait un germe
inversible au point $\psi(x)$. Mais en vertu de
(\hyperref[0.7.6.17-fr]{0, 7.6.17}), les sections de
$\mathscr{O}_{\mathfrak{X}}$ (resp. $\mathscr{O}_{\mathfrak{Y}}$)
au-dessus de $\mathfrak{X}$ (resp. $\mathfrak{Y}$) dont le germe n'est pas
inversible au point $x$ (resp. $\psi(x)$) sont exactement les éléments de
$\mathfrak{j}_x$
\oldpage[I]{183}
(resp. $\mathfrak{j}_{\psi(x)}$)~; la remarque précédente montre donc que
${}^a\varphi=\psi$. Enfin, pour tout $g\in B$, le diagramme
\[
  \xymatrix{
    B=\Gamma(\mathfrak{Y},\mathscr{O}_{\mathfrak{Y}})
      \ar[r]^\varphi\ar[d] &
    \Gamma(\mathfrak{X},\mathscr{O}_{\mathfrak{X}})=A\ar[d]\\
    B_{\{g\}}=\Gamma(\mathfrak{D}(g),\mathscr{O}_{\mathfrak{Y}})
      \ar[r]_{\Gamma(\theta_{\mathfrak{D}(g)})} &
    \Gamma(\mathfrak{D}(\varphi(g)),\mathscr{O}_{\mathfrak{X}})
      =A_{\{\varphi(g)\}}
  }
\]
est commutatif~; d'après la propriété universelle des anneaux complets de
fractions (\hyperref[0.7.6.6-fr]{0, 7.6.6}),
$\theta_{\mathfrak{D}(g)}$ est égale à
$\widetilde{\varphi}_{\mathfrak{D}(g)}$ pour tout $g\in B$, donc
(\hyperref[0.3.2.5-fr]{0, 3.2.5}) on a
$\theta=\widetilde{\varphi}$.
\end{proof}

Nous dirons qu'un morphisme $(\psi,\theta)$ d'espaces topologiquement
annelés vérifiant la condition de (\hyperref[I.10.2.2-fr]{10.2.2}) est un
\emph{morphisme de schémas formels affines}. On peut dire que les foncteurs
$\operatorname{Spf}(A)$ en $A$ et
$\Gamma(\mathfrak{X},\mathscr{O}_{\mathfrak{X}})$ en $\mathfrak{X}$
définissent une \emph{équivalence} de la catégorie des anneaux admissibles
et de la catégorie duale de la catégorie des schémas formels affines
(T, I, 1.2).

\begin{env}[10.2.3]
\label{I.10.2.3-fr}
Comme cas particulier de (\hyperref[I.10.2.2-fr]{10.2.2}), notons que, pour
$f\in A$, l'injection canonique du schéma formel affine induit par
$\mathfrak{X}$ sur $\mathfrak{D}(f)$ correspond à l'homomorphisme continu
canonique $A\to A_{\{f\}}$. Sous les hypothèses de
(\hyperref[I.10.2.2-fr]{10.2.2}), soient $h$ un élément de $B$, $g$ un
élément de $A$, multiple de $\varphi(h)$~; on a alors
$\psi(\mathfrak{D}(g))\subset\mathfrak{D}(h)$~; la restriction de $u$ à
$\mathfrak{D}(g)$, considérée comme morphisme de $\mathfrak{D}(g)$ dans
$\mathfrak{D}(h)$, est l'unique morphisme $v$ rendant commutatif le
diagramme
\[
  \xymatrix{
    \mathfrak{D}(g)\ar[r]^v\ar[d] &
    \mathfrak{D}(h)\ar[d]\\
    \mathfrak{X}\ar[r]_u &
    \mathfrak{Y}
  }
\]
Ce morphisme correspond à l'unique homomorphisme continu
$\varphi':B_{\{h\}}\to A_{\{g\}}$
(\hyperref[0.7.6.7-fr]{0, 7.6.7}) rendant commutatif le diagramme
\[
  \xymatrix{
    A\ar[d] &
    B\ar[l]_\varphi\ar[d]\\
    A_{\{g\}} &
    B_{\{h\}}\ar[l]_{\varphi'}
  }
\]
\end{env}

\subsection{Idéaux de définition d'un schéma formel affine.}
\label{subsection:I.10.3-fr}

\begin{env}[10.3.1]
\label{I.10.3.1-fr}
Soient $A$ un anneau admissible, $\mathfrak{J}$ un idéal ouvert de $A$,
$\mathfrak{X}$ le schéma formel affine $\operatorname{Spf}(A)$. Soit
$(\mathfrak{J}_\lambda)$ l'ensemble des idéaux de définition de $A$
contenus dans $\mathfrak{J}$~; alors
$\widetilde{\mathfrak{J}}/\widetilde{\mathfrak{J}}_\lambda$ est un
faisceau d'idéaux de
$\widetilde{A}/\widetilde{\mathfrak{J}}_\lambda$. Désignons par
$\mathfrak{J}^\Delta$ la limite projective des faisceaux induits sur
$\mathfrak{X}$ par
$\widetilde{\mathfrak{J}}/\widetilde{\mathfrak{J}}_\lambda$, qui
s'identifie à un \emph{faisceau d'idéaux} de
$\mathscr{O}_{\mathfrak{X}}$ (\hyperref[0.3.2.6-fr]{0, 3.2.6}). Pour tout
$f\in A$, $\Gamma(\mathfrak{D}(f),\mathfrak{J}^\Delta)$ est la limite
projective de
$S_f^{-1}\mathfrak{J}/S_f^{-1}\mathfrak{J}_\lambda$, autrement dit
s'identifie à l'idéal ouvert $\mathfrak{J}_{\{f\}}$ de l'anneau
$A_{\{f\}}$ (\hyperref[0.7.6.9-fr]{0, 7.6.9}), et en particulier
$\Gamma(\mathfrak{X},\mathfrak{J}^\Delta)=\mathfrak{J}$~; on en conclut
(les $\mathfrak{D}(f)$ formant une base de la topologie de $\mathfrak{X}$)
que l'on a
\[
  \mathfrak{J}^\Delta|\mathfrak{D}(f)
  =(\mathfrak{J}_{\{f\}})^\Delta.
  \tag{10.3.1.1}\label{I.10.3.1.1-fr}
\]
\end{env}

\oldpage[I]{184}
\begin{env}[10.3.2]
\label{I.10.3.2-fr}
Avec les notations de (\hyperref[I.10.3.1-fr]{10.3.1}), pour tout $f\in A$,
l'application canonique de
$A_{\{f\}}=\Gamma(\mathfrak{D}(f),\mathscr{O}_{\mathfrak{X}})$ dans
$\Gamma\bigl(\mathfrak{D}(f),
(\widetilde{A}/\widetilde{\mathfrak{J}})|\mathfrak{X}\bigr)
=S_f^{-1}A/S_f^{-1}\mathfrak{J}$
est \emph{surjective} et a pour noyau
$\Gamma(\mathfrak{D}(f),\mathfrak{J}^\Delta)=\mathfrak{J}_{\{f\}}$
(\hyperref[0.7.6.9-fr]{0, 7.6.9})~; ces applications définissent donc un
homomorphisme \emph{surjectif} continu, dit \emph{canonique}, du faisceau
d'anneaux topologiques $\mathscr{O}_{\mathfrak{X}}$ sur le faisceau
d'anneaux discrets
$(\widetilde{A}/\widetilde{\mathfrak{J}})|\mathfrak{X}$, dont le noyau est
$\mathfrak{J}^\Delta$~; cet homomorphisme n'est autre d'ailleurs que
$\widetilde{\varphi}$ (\hyperref[I.10.2.1-fr]{10.2.1}), où $\varphi$ est
l'homomorphisme continu $A\to A/\mathfrak{J}$~; le morphisme
$({}^a\varphi,\widetilde{\varphi}):
\operatorname{Spec}(A/\mathfrak{J})\to\mathfrak{X}$ de schémas formels
affines (où ${}^a\varphi$ est d'ailleurs l'homéomorphisme identique de
$\mathfrak{X}$ sur lui-même) est encore dit \emph{canonique}. On a donc,
d'après ce qui précède, un \emph{isomorphisme canonique}
\[
  \mathscr{O}_{\mathfrak{X}}/\mathfrak{J}^\Delta
  \xrightarrow{\sim}
  (\widetilde{A}/\widetilde{\mathfrak{J}})|\mathfrak{X}.
  \tag{10.3.2.1}\label{I.10.3.2.1-fr}
\]

Il est clair (en vertu de
$\Gamma(\mathfrak{X},\mathfrak{J}^\Delta)=\mathfrak{J}$) que
l'application $\mathfrak{J}\to\mathfrak{J}^\Delta$ est
\emph{strictement croissante}~; d'après ce qui précède, pour
$\mathfrak{J}\subset\mathfrak{J}'$, le faisceau
${\mathfrak{J}'}^\Delta/\mathfrak{J}^\Delta$ est canoniquement isomorphe à
$\widetilde{\mathfrak{J}'}/\widetilde{\mathfrak{J}}
=\widetilde{\mathfrak{J}'/\mathfrak{J}}$.
\end{env}

\begin{env}[10.3.3]
\label{I.10.3.3-fr}
Les hypothèses et notations étant toujours celles de
(\hyperref[I.10.3.1-fr]{10.3.1}), nous dirons qu'un faisceau d'idéaux
$\mathscr{J}$ de $\mathscr{O}_{\mathfrak{X}}$ est un
\emph{faisceau d'idéaux de définition de $\mathfrak{X}$} (ou un
\emph{Idéal de définition de $\mathfrak{X}$}) si, pour tout
$x\in\mathfrak{X}$, il existe un voisinage ouvert de $x$ de la forme
$\mathfrak{D}(f)$, où $f\in A$, tel que
$\mathscr{J}|\mathfrak{D}(f)$ soit de la forme $\mathfrak{H}^\Delta$, où
$\mathfrak{H}$ est un idéal de définition de $A_{\{f\}}$.
\end{env}

\begin{proposition}[10.3.4]
\label{I.10.3.4-fr}
Pour tout $f\in A$, tout Idéal de définition de $\mathfrak{X}$ induit un
Idéal de définition de $\mathfrak{D}(f)$.
\end{proposition}

\begin{proof}
Cela résulte de (\hyperref[I.10.3.1.1-fr]{10.3.1.1}).
\end{proof}

\begin{proposition}[10.3.5]
\label{I.10.3.5-fr}
Si $A$ est un anneau admissible, tout Idéal de définition de
$\mathfrak{X}=\operatorname{Spf}(A)$ est de la forme
$\mathfrak{J}^\Delta$, où $\mathfrak{J}$ est un idéal de définition de
$A$, uniquement déterminé.
\end{proposition}

\begin{proof}
En effet, soit $\mathscr{J}$ un Idéal de définition de $\mathfrak{X}$~;
par hypothèse, et puisque $\mathfrak{X}$ est quasi-compact, il y a un nombre
fini d'éléments $f_i\in A$ tels que les $\mathfrak{D}(f_i)$ recouvrent
$\mathfrak{X}$ et que
$\mathscr{J}|\mathfrak{D}(f_i)=\mathfrak{H}_i^\Delta$, où
$\mathfrak{H}_i$ est un idéal de définition de $A_{\{f_i\}}$. Pour tout
$i$, il existe donc un idéal ouvert $\mathfrak{K}_i$ de $A$ tel que
$(\mathfrak{K}_i)_{\{f_i\}}=\mathfrak{H}_i$
(\hyperref[0.7.6.9-fr]{0, 7.6.9})~; soit $\mathfrak{K}$ un idéal de
définition de $A$ contenu dans tous les $\mathfrak{K}_i$. L'image canonique
de $\mathscr{J}/\mathfrak{K}^\Delta$ dans le faisceau structural
$\widetilde{A/\mathfrak{K}}$ de
$\operatorname{Spec}(A/\mathfrak{K})$
(\hyperref[I.10.3.2-fr]{10.3.2}) est donc telle que sa restriction à
$\mathfrak{D}(f_i)$ soit égale à celle de
$\widetilde{\mathfrak{K}_i/\mathfrak{K}}$~; on en conclut que cette image
canonique est un faisceau quasi-cohérent sur
$\operatorname{Spec}(A/\mathfrak{K})$, donc de la forme
$\widetilde{\mathfrak{J}/\mathfrak{K}}$, où $\mathfrak{J}$ est un idéal de
$A$ contenant $\mathfrak{K}$ (\hyperref[I.1.4.1-fr]{1.4.1}) d'où
$\mathscr{J}=\mathfrak{J}^\Delta$
(\hyperref[I.10.3.2-fr]{10.3.2})~; en outre, comme pour tout $i$ il existe
un entier $n_i$ tel que
$\mathfrak{H}_i^{n_i}\subset\mathfrak{K}_{\{f_i\}}$, on aura, en désignant
par $n$ le plus grand des $n_i$,
$(\mathscr{J}/\mathfrak{K}^\Delta)^n=0$, et par suite
(\hyperref[I.10.3.2-fr]{10.3.2})
$(\widetilde{\mathfrak{J}/\mathfrak{K}})^n=0$, d'où finalement
$(\mathfrak{J}/\mathfrak{K})^n=0$
(\hyperref[I.1.3.13-fr]{1.3.13}), ce qui prouve que $\mathfrak{J}$ est un
idéal de définition de $A$ (\hyperref[0.7.1.4-fr]{0, 7.1.4}).
\end{proof}

\begin{proposition}[10.3.6]
\label{I.10.3.6-fr}
Soient $A$ un anneau adique, $\mathfrak{J}$ un idéal de définition de $A$
tel que $\mathfrak{J}/\mathfrak{J}^2$ soit un
$A/\mathfrak{J}$-module de type fini. Pour tout entier $n>0$, on a alors
$(\mathfrak{J}^\Delta)^n=(\mathfrak{J}^n)^\Delta$.
\end{proposition}

\begin{proof}
En effet, pour tout $f\in A$, on a (puisque $\mathfrak{J}^n$ est un idéal
ouvert)
\[
  \bigl(\Gamma(\mathfrak{D}(f),\mathfrak{J}^\Delta)\bigr)^n
  =(\mathfrak{J}_{\{f\}})^n
  =(\mathfrak{J}^n)_{\{f\}}
  =\Gamma\bigl(\mathfrak{D}(f^n),(\mathfrak{J}^n)^\Delta\bigr)
\]
\oldpage[I]{185}
en vertu de (\hyperref[I.10.3.1.1-fr]{10.3.1.1}) et de
(\hyperref[0.7.6.12-fr]{0, 7.6.12}). Comme
$(\mathfrak{J}^\Delta)^n$ est associé au préfaisceau
$U\mapsto(\Gamma(U,\mathfrak{J}^\Delta))^n$
(\hyperref[0.4.1.6-fr]{0, 4.1.6}), le corollaire en résulte, puisque les
$\mathfrak{D}(f)$ forment une base de la topologie de $\mathfrak{X}$.
\end{proof}

\begin{env}[10.3.7]
\label{I.10.3.7-fr}
On dit qu'une famille $(\mathscr{J}_\lambda)$ d'Idéaux de définition de
$\mathfrak{X}$ est un \emph{système fondamental d'Idéaux de définition}
si tout Idéal de définition de $\mathfrak{X}$ contient un des
$\mathscr{J}_\lambda$~; comme
$\mathscr{J}_\lambda=\mathfrak{J}_\lambda^\Delta$, il revient au même de
dire que les $\mathfrak{J}_\lambda$ forment un \emph{système fondamental
de voisinages de 0} dans $A$. Soit $(f_\alpha)$ une famille d'éléments de
$A$ tels que les $\mathfrak{D}(f_\alpha)$ recouvrent $\mathfrak{X}$. Si
$(\mathscr{J}_\lambda)$ est une famille filtrante décroissante d'idéaux de
$\mathscr{O}_{\mathfrak{X}}$ telle que pour tout $\alpha$, la famille
$(\mathscr{J}_\lambda|\mathfrak{D}(f_\alpha))$ soit un système fondamental
d'Idéaux de définition de $\mathfrak{D}(f_\alpha)$, alors
$(\mathscr{J}_\lambda)$ est un système fondamental d'Idéaux de définition
de $\mathfrak{X}$. En effet, pour tout Idéal de définition $\mathscr{J}$
de $\mathfrak{X}$, il y a un recouvrement fini de $\mathfrak{X}$ par des
$\mathfrak{D}(f_i)$ tel que, pour tout $i$,
$\mathscr{J}_{\lambda_i}|\mathfrak{D}(f_i)$ soit un Idéal de définition
de $\mathfrak{D}(f_i)$ contenu dans
$\mathscr{J}|\mathfrak{D}(f_i)$. Si $\mu$ est un indice tel que
$\mathscr{J}_\mu\subset\mathscr{J}_{\lambda_i}$ pour tout $i$, il résulte
de (\hyperref[I.10.3.3-fr]{10.3.3}) que $\mathscr{J}_\mu$ est un Idéal de
définition de $\mathfrak{X}$, évidemment contenu dans $\mathscr{J}$, d'où
notre assertion.
\end{env}

\subsection{Préschémas formels et morphismes de préschémas formels.}
\label{subsection:I.10.4-fr}

\begin{env}[10.4.1]
\label{I.10.4.1-fr}
Étant donné un espace topologiquement annelé $\mathfrak{X}$, on dit qu'un
ouvert $U\subset\mathfrak{X}$ est un \emph{ouvert formel affine} (resp. un
\emph{ouvert formel affine adique}, resp. un \emph{ouvert formel affine
noethérien}) si l'espace topologiquement annelé induit par $\mathfrak{X}$
sur $U$ est un schéma formel affine (resp. un tel schéma dont l'anneau est
adique, resp. adique et noethérien).
\end{env}

\begin{definition}[10.4.2]
\label{I.10.4.2-fr}
On appelle préschéma formel un espace topologiquement annelé $\mathfrak{X}$
dont tout point admet un voisinage ouvert formel affine. On dit que le
préschéma formel $\mathfrak{X}$ est adique (resp. localement noethérien) si
tout point de $\mathfrak{X}$ admet un voisinage ouvert formel affine adique
(resp. noethérien). On dit que $\mathfrak{X}$ est noethérien s'il est
localement noethérien et si son espace sous-jacent est quasi-compact (donc
noethérien).
\end{definition}

\begin{proposition}[10.4.3]
\label{I.10.4.3-fr}
Si $\mathfrak{X}$ est un préschéma formel (resp. localement noethérien),
les ensembles ouverts formels affines (resp. affines noethériens) forment
une base de la topologie de $\mathfrak{X}$.
\end{proposition}

\begin{proof}
Cela résulte de (\hyperref[I.10.4.2-fr]{10.4.2}) et
(\hyperref[I.10.1.4-fr]{10.1.4}) en tenant compte de ce que si $A$ est un
anneau adique noethérien, il en est de même de $A_{\{f\}}$ pour tout
$f\in A$ (\hyperref[0.7.6.11-fr]{0, 7.6.11}).
\end{proof}

\begin{corollary}[10.4.4]
\label{I.10.4.4-fr}
Si $\mathfrak{X}$ est un préschéma formel (resp. un préschéma formel
localement noethérien, resp. noethérien), l'espace topologiquement annelé
induit sur tout ouvert de $\mathfrak{X}$ est encore un préschéma formel
(resp. un préschéma formel localement noethérien, resp. noethérien).
\end{corollary}

\begin{definition}[10.4.5]
\label{I.10.4.5-fr}
Étant donnés deux préschémas formels $\mathfrak{X}$, $\mathfrak{Y}$, on
appelle morphisme (de préschémas formels) de $\mathfrak{X}$ dans
$\mathfrak{Y}$ tout morphisme $(\psi,\theta)$ d'espaces topologiquement
annelés tel que, pour tout $x\in\mathfrak{X}$, $\theta_x^\sharp$ soit un
homomorphisme local
$\mathscr{O}_{\psi(x)}\to\mathscr{O}_x$.
\end{definition}

Il est immédiat que le composé de deux morphismes de préschémas formels est
encore un tel morphisme~; les préschémas formels forment donc une
\emph{catégorie}, et on notera
$\operatorname{Hom}(\mathfrak{X},\mathfrak{Y})$ l'ensemble des morphismes
d'un préschéma formel $\mathfrak{X}$ dans un préschéma formel
$\mathfrak{Y}$.

\oldpage[I]{186}
Si $U$ est une partie ouverte de $\mathfrak{X}$, l'injection canonique
dans $\mathfrak{X}$ du préschéma formel induit par $\mathfrak{X}$ sur $U$
est un morphisme de préschémas formels (et même un
\emph{monomorphisme} d'espaces topologiquement annelés
(\hyperref[0.4.1.1-fr]{0, 4.1.1})).

\begin{proposition}[10.4.6]
\label{I.10.4.6-fr}
Soient $\mathfrak{X}$ un préschéma formel,
$\mathfrak{S}=\operatorname{Spf}(A)$ un schéma formel affine. Il existe
une correspondance biunivoque canonique entre les morphismes du préschéma
formel $\mathfrak{X}$ dans le préschéma formel $\mathfrak{S}$ et les
homomorphismes continus de l'anneau $A$ dans l'anneau topologique
$\Gamma(\mathfrak{X},\mathscr{O}_{\mathfrak{X}})$.
\end{proposition}

La démonstration est la même que celle de
(\hyperref[I.2.2.4-fr]{2.2.4}), en remplaçant «~homomorphisme~» par
«~homomorphisme continu~», «~ouvert affine~» par «~ouvert formel
affine~», et en utilisant (\hyperref[I.10.2.2-fr]{10.2.2}) au lieu de
(\hyperref[I.1.7.3-fr]{1.7.3})~; nous en laissons les détails au lecteur.

\begin{env}[10.4.7]
\label{I.10.4.7-fr}
Étant donné un préschéma formel $\mathfrak{S}$, on dit que la donnée d'un
préschéma formel $\mathfrak{X}$ et d'un morphisme
$\varphi:\mathfrak{X}\to\mathfrak{S}$ définit un préschéma formel
$\mathfrak{X}$ \emph{au-dessus de $\mathfrak{S}$} ou un
\emph{$\mathfrak{S}$-préschéma formel}, $\varphi$ étant appelé le
\emph{morphisme structural} du $\mathfrak{S}$-préschéma
$\mathfrak{X}$. Si $\mathfrak{S}=\operatorname{Spf}(A)$, où $A$ est un
anneau admissible, on dit aussi que le $\mathfrak{S}$-préschéma formel
$\mathfrak{X}$ est un \emph{$A$-préschéma formel} ou un préschéma formel
\emph{au-dessus de $A$}. Un préschéma formel quelconque peut toujours
être considéré comme un préschéma formel au-dessus de $\mathbf{Z}$ (muni
de la topologie discrète).

Si $\mathfrak{X}$, $\mathfrak{Y}$ sont deux $\mathfrak{S}$-préschémas
formels, on dit qu'un morphisme
$u:\mathfrak{X}\to\mathfrak{Y}$ est un
\emph{$\mathfrak{S}$-morphisme} si le diagramme
\[
  \xymatrix{
    \mathfrak{X}\ar[rr]^u\ar[rd] & &
    \mathfrak{Y}\ar[ld]\\
    & \mathfrak{S}
  }
\]
(où les flèches obliques sont les morphismes structuraux) est
commutatif. Avec cette définition, les $\mathfrak{S}$-préschémas formels
(pour $\mathfrak{S}$ fixé) forment une \emph{catégorie}. On désigne par
$\operatorname{Hom}_{\mathfrak{S}}(\mathfrak{X},\mathfrak{Y})$
l'ensemble des $\mathfrak{S}$-morphismes du $\mathfrak{S}$-préschéma
formel $\mathfrak{X}$ dans le $\mathfrak{S}$-préschéma formel
$\mathfrak{Y}$. Lorsque $\mathfrak{S}=\operatorname{Spf}(A)$, on dit
aussi \emph{$A$-morphisme} au lieu de
\emph{$\mathfrak{S}$-morphisme}.
\end{env}

\begin{env}[10.4.8]
\label{I.10.4.8-fr}
Comme tout schéma affine peut être considéré comme un schéma formel affine
(\hyperref[I.10.1.2-fr]{10.1.2}), tout préschéma (usuel) peut être
considéré comme un préschéma formel. Il résulte en outre de
(\hyperref[I.10.4.5-fr]{10.4.5}) que pour les préschémas \emph{usuels},
les morphismes (resp. $S$-morphismes) de préschémas \emph{formels}
coïncident avec les morphismes (resp. $S$-morphismes) définis au
\S~\hyperref[section:I.2-fr]{2}.
\end{env}

\subsection{Idéaux de définition des préschémas formels.}
\label{subsection:I.10.5-fr}

\begin{env}[10.5.1]
\label{I.10.5.1-fr}
Soit $\mathfrak{X}$ un préschéma formel~; on dit qu'un
$\mathscr{O}_{\mathfrak{X}}$-Idéal $\mathscr{J}$ est un
\emph{faisceau d'idéaux de définition} (ou un
\emph{Idéal de définition}) de $\mathfrak{X}$ si tout
$x\in\mathfrak{X}$ possède un voisinage ouvert formel affine $U$ tel que
$\mathscr{J}|U$ soit un Idéal de définition du schéma formel affine
induit par $\mathfrak{X}$ sur $U$
(\hyperref[I.10.3.3-fr]{10.3.3})~; en vertu de
(\hyperref[I.10.3.1.1-fr]{10.3.1.1}) et
(\hyperref[I.10.4.3-fr]{10.4.3}), pour tout ouvert
$V\subset\mathfrak{X}$, $\mathscr{J}|V$ est alors un Idéal de définition
du préschéma formel induit par $\mathfrak{X}$ sur $V$.

On dit qu'une famille $(\mathscr{J}_\lambda)$ d'Idéaux de définition de
$\mathfrak{X}$ est un \emph{système fondamental}
\oldpage[I]{187}
\emph{d'Idéaux de définition} s'il existe un recouvrement
$(U_\alpha)$ de $\mathfrak{X}$ par des ouverts formels affines tel que,
pour tout $\alpha$, la famille des
$\mathscr{J}_\lambda|U_\alpha$ soit un système fondamental d'Idéaux de
définition (\hyperref[I.10.3.6-fr]{10.3.6}) du schéma formel affine
induit par $\mathfrak{X}$ sur $U_\alpha$. Il résulte de la remarque finale
de (\hyperref[I.10.3.7-fr]{10.3.7}) que lorsque $\mathfrak{X}$ est un
schéma formel affine, cette définition coïncide avec la définition donnée
dans (\hyperref[I.10.3.7-fr]{10.3.7}). Pour tout ouvert $V$ de
$\mathfrak{X}$, les restrictions $\mathscr{J}_\lambda|V$ forment alors un
système fondamental d'Idéaux de définition du préschéma formel induit sur
$V$, en vertu de (\hyperref[I.10.3.1.1-fr]{10.3.1.1}). Si
$\mathfrak{X}$ est un préschéma formel \emph{localement noethérien}, et
$\mathscr{J}$ un Idéal de définition de $\mathfrak{X}$, il résulte de
(\hyperref[I.10.3.6-fr]{10.3.6}) que les puissances $\mathscr{J}^n$
forment un système fondamental d'Idéaux de définition de $\mathfrak{X}$.
\end{env}

\begin{env}[10.5.2]
\label{I.10.5.2-fr}
Soient $\mathfrak{X}$ un préschéma formel, $\mathscr{J}$ un Idéal de
définition de $\mathfrak{X}$. Alors l'espace annelé
$(\mathfrak{X},\mathscr{O}_{\mathfrak{X}}/\mathscr{J})$ est un
\emph{préschéma} (usuel), qui est affine (resp. localement noethérien,
resp. noethérien) lorsque $\mathfrak{X}$ est un schéma formel affine
(resp. un préschéma formel localement noethérien, resp. noethérien)~; on
est en effet aussitôt ramené au cas affine, et alors la proposition a déjà
été démontrée dans (\hyperref[I.10.3.2-fr]{10.3.2}). En outre, si
$\theta:\mathscr{O}_{\mathfrak{X}}\to
\mathscr{O}_{\mathfrak{X}}/\mathscr{J}$ est l'homomorphisme canonique,
$u=(1_{\mathfrak{X}},\theta)$ est un \emph{morphisme} (dit
\emph{canonique}) de préschémas formels
$(\mathfrak{X},\mathscr{O}_{\mathfrak{X}}/\mathscr{J})\to
(\mathfrak{X},\mathscr{O}_{\mathfrak{X}})$, car ici encore, cela a été vu
dans le cas affine (\hyperref[I.10.3.2-fr]{10.3.2}), auquel on se ramène
aussitôt.
\end{env}

\begin{proposition}[10.5.3]
\label{I.10.5.3-fr}
Soient $\mathfrak{X}$ un préschéma formel, $(\mathscr{J}_\lambda)$ un
système fondamental d'Idéaux de définition de $\mathfrak{X}$. Alors le
faisceau d'anneaux topologiques $\mathscr{O}_{\mathfrak{X}}$ est limite
projective des faisceaux d'anneaux pseudo-discrets
(\hyperref[0.3.8.1-fr]{0, 3.8.1})
$\mathscr{O}_{\mathfrak{X}}/\mathscr{J}_\lambda$.
\end{proposition}

\begin{proof}
Comme la topologie de $\mathfrak{X}$ admet une base d'ouverts formels
affines quasi-compacts (\hyperref[I.10.4.3-fr]{10.4.3}), on est ramené au
cas affine, où la proposition est conséquence de
(\hyperref[I.10.3.5-fr]{10.3.5}),
(\hyperref[I.10.3.2-fr]{10.3.2}) et de la définition
(\hyperref[I.10.1.1-fr]{10.1.1}).
\end{proof}

Il n'est pas certain que tout préschéma formel admette des Idéaux de
définition. Toutefois~:

\begin{proposition}[10.5.4]
\label{I.10.5.4-fr}
Soit $\mathfrak{X}$ un préschéma formel localement noethérien. Il existe
un plus grand Idéal de définition $\mathscr{T}$ de $\mathfrak{X}$~;
c'est le seul Idéal de définition $\mathscr{J}$ tel que le préschéma
$(\mathfrak{X},\mathscr{O}_{\mathfrak{X}}/\mathscr{J})$ soit réduit. Si
$\mathscr{J}$ est un Idéal de définition de $\mathfrak{X}$,
$\mathscr{T}$ est l'image réciproque par
$\mathscr{O}_{\mathfrak{X}}\to
\mathscr{O}_{\mathfrak{X}}/\mathscr{J}$ du Nilradical de
$\mathscr{O}_{\mathfrak{X}}/\mathscr{J}$.
\end{proposition}

\begin{proof}
Supposons d'abord que $\mathfrak{X}=\operatorname{Spf}(A)$, où $A$ est
un anneau adique noethérien. L'existence et les propriétés de
$\mathscr{T}$ résultent immédiatement de
(\hyperref[I.10.3.4-fr]{10.3.4}) et
(\hyperref[I.5.1.1-fr]{5.1.1}), compte tenu de l'existence et des
propriétés du plus grand idéal de définition de $A$
(\hyperref[0.7.1.6-fr]{0, 7.1.6} et
\hyperref[0.7.1.7-fr]{7.1.7}).

Pour prouver l'existence et les propriétés de $\mathscr{T}$ dans le cas
général, il suffit de montrer que si $U\supset V$ sont deux ouverts
formels affines noethériens de $\mathfrak{X}$, le plus grand Idéal de
définition $\mathscr{T}_U$ de $U$ induit le plus grand Idéal de définition
$\mathscr{T}_V$ de $V$~; mais comme
$(V,(\mathscr{O}_{\mathfrak{X}}|V)/(\mathscr{T}_U|V))$ est réduit, cela
résulte de ce qui précède.
\end{proof}

On désigne par $\mathfrak{X}_{\mathrm{red}}$ le préschéma (usuel) réduit
$(\mathfrak{X},\mathscr{O}_{\mathfrak{X}}/\mathscr{T})$.

\begin{corollary}[10.5.5]
\label{I.10.5.5-fr}
Soient $\mathfrak{X}$ un préschéma formel localement noethérien,
$\mathscr{T}$ le plus grand Idéal de définition de $\mathfrak{X}$~; pour
tout ouvert $V$ de $\mathfrak{X}$, $\mathscr{T}|V$ est le plus grand
Idéal de définition du préschéma formel induit par $\mathfrak{X}$ sur
$V$.
\end{corollary}

\oldpage[I]{188}

\begin{proposition}[10.5.6]
\label{I.10.5.6-fr}
Soient $\mathfrak{X}$, $\mathfrak{Y}$ deux préschémas formels,
$\mathscr{J}$ (resp. $\mathscr{K}$) un Idéal de définition de
$\mathfrak{X}$ (resp. $\mathfrak{Y}$),
$f:\mathfrak{X}\to\mathfrak{Y}$ un morphisme de préschémas formels.
\begin{enumerate}
  \item[\rm{(i)}] Si
    $f^*(\mathscr{K})\mathscr{O}_{\mathfrak{X}}\subset\mathscr{J}$,
    il existe un unique morphisme
    $f':(\mathfrak{X},\mathscr{O}_{\mathfrak{X}}/\mathscr{J})\to
    (\mathfrak{Y},\mathscr{O}_{\mathfrak{Y}}/\mathscr{K})$ de
    préschémas usuels rendant commutatif le diagramme
    \[
      \label{I.10.5.6.1-fr}
      \xymatrix{
        (\mathfrak{X},\mathscr{O}_{\mathfrak{X}})\ar[r]^f &
        (\mathfrak{Y},\mathscr{O}_{\mathfrak{Y}})\\
        (\mathfrak{X},\mathscr{O}_{\mathfrak{X}}/\mathscr{J})
          \ar[r]_{f'}\ar[u] &
        (\mathfrak{Y},\mathscr{O}_{\mathfrak{Y}}/\mathscr{K})\ar[u]
      }
      \tag{10.5.6.1}
    \]
    où les flèches verticales sont les morphismes canoniques.
  \item[\rm{(ii)}] Supposons que
    $\mathfrak{X}=\operatorname{Spf}(A)$,
    $\mathfrak{Y}=\operatorname{Spf}(B)$ soient des schémas formels
    affines, $\mathscr{J}=\mathfrak{J}^{\Delta}$,
    $\mathscr{K}=\mathfrak{K}^{\Delta}$, où $\mathfrak{J}$ (resp.
    $\mathfrak{K}$) est un idéal de définition de $A$ (resp. $B$), et
    $f=({}^a\varphi,\widetilde{\varphi})$, où
    $\varphi:B\to A$ est un homomorphisme continu~; pour que
    $f^*(\mathscr{K})\mathscr{O}_{\mathfrak{X}}\subset\mathscr{J}$,
    il faut et il suffit que
    $\varphi(\mathfrak{K})\subset\mathfrak{J}$, et $f'$ est alors le
    morphisme $({}^a\varphi',\widetilde{\varphi'})$, où
    $\varphi':B/\mathfrak{K}\to A/\mathfrak{J}$ est l'homomorphisme
    déduit de $\varphi$ par passage aux quotients.
\end{enumerate}
\end{proposition}

\begin{proof}
\medskip\noindent
\begin{enumerate}
  \item[\rm{(i)}] Si $f=(\psi,\theta)$, l'hypothèse entraîne que
    l'image par
    $\theta^\sharp:\psi^*(\mathscr{O}_{\mathfrak{Y}})\to
    \mathscr{O}_{\mathfrak{X}}$ du faisceau d'idéaux
    $\psi^*(\mathscr{K})$ de $\psi^*(\mathscr{O}_{\mathfrak{Y}})$ est
    contenue dans $\mathscr{J}$
    (\hyperref[0.4.3.5-fr]{0, 4.3.5}). Par passage aux quotients, on
    déduit donc de $\theta^\sharp$ un homomorphisme de faisceau d'anneaux
    \[
      \omega:\psi^*(\mathscr{O}_{\mathfrak{Y}}/\mathscr{K})
      =\psi^*(\mathscr{O}_{\mathfrak{Y}})/\psi^*(\mathscr{K})
      \to\mathscr{O}_{\mathfrak{X}}/\mathscr{J}~;
    \]
    en outre, comme pour tout $x\in\mathfrak{X}$,
    $\theta_x^\sharp$ est un homomorphisme \emph{local}, il en est de
    même de $\omega_x$. Le morphisme d'espaces annelés
    $(\psi,\omega^b)$ est donc
    (\hyperref[I.2.2.1-fr]{2.2.1}) l'unique morphisme $f'$ d'espaces
    annelés répondant à la question.
  \item[\rm{(ii)}] La correspondance canonique fonctorielle entre
    morphismes de préschémas formels affines et homomorphismes continus
    d'anneaux (\hyperref[I.10.2.2-fr]{10.2.2}) montre que dans le cas
    considéré, la relation
    $f^*(\mathscr{K})\mathscr{O}_{\mathfrak{X}}\subset\mathscr{J}$
    entraîne que l'on a
    $f'=({}^a\varphi',\widetilde{\varphi'})$, où
    $\varphi':B/\mathfrak{K}\to A/\mathfrak{J}$ est l'unique
    homomorphisme rendant commutatif le diagramme
    \[
      \label{I.10.5.6.2-fr}
      \xymatrix{
        B\ar[r]^\varphi\ar[d] & A\ar[d]\\
        B/\mathfrak{K}\ar[r]_{\varphi'} & A/\mathfrak{J}
      }
      \tag{10.5.6.2}
    \]
    L'existence de $\varphi'$ implique donc que
    $\varphi(\mathfrak{K})\subset\mathfrak{J}$. Inversement, si cette
    condition est vérifiée, en désignant par $\varphi'$ l'unique
    homomorphisme rendant commutatif le diagramme
    (\hyperref[I.10.5.6.2-fr]{10.5.6.2}) et posant
    $f'=({}^a\varphi',\widetilde{\varphi'})$, il est clair que le
    diagramme (\hyperref[I.10.5.6.1-fr]{10.5.6.1}) est commutatif~; la
    considération des homomorphismes
    ${}^a\varphi^*(\mathscr{O}_{\mathfrak{Y}})\to
    \mathscr{O}_{\mathfrak{X}}$ et
    ${}^a{\varphi'}^*(\mathscr{O}_{\mathfrak{Y}}/\mathscr{K})\to
    \mathscr{O}_{\mathfrak{X}}/\mathscr{J}$ correspondant à $f$ et
    $f'$ respectivement, montre alors que cela entraîne la relation
    $f^*(\mathscr{K})\mathscr{O}_{\mathfrak{X}}\subset\mathscr{J}$.
\end{enumerate}
\end{proof}

Il est clair que la correspondance $f\mapsto f'$ définie ci-dessus est
\emph{fonctorielle}.

\subsection{Préschémas formels comme limites inductives de préschémas.}
\label{subsection:I.10.6-fr}

\begin{env}[10.6.1]
\label{I.10.6.1-fr}
Soient $\mathfrak{X}$ un préschéma formel,
$(\mathscr{J}_\lambda)$ un système fondamental d'Idéaux de définition de
$\mathfrak{X}$~; pour tout $\lambda$, soit $f_\lambda$ le morphisme
canonique
$(\mathfrak{X},\mathscr{O}_{\mathfrak{X}}/\mathscr{J}_\lambda)
\to\mathfrak{X}$ (\hyperref[I.10.5.2-fr]{10.5.2})~; pour
$\mathscr{J}_\mu\subset\mathscr{J}_\lambda$, l'homomorphisme canonique
$\mathscr{O}_{\mathfrak{X}}/\mathscr{J}_\mu\to
\mathscr{O}_{\mathfrak{X}}/\mathscr{J}_\lambda$ définit un morphisme

\oldpage[I]{189}

canonique
$f_{\mu\lambda}:(\mathfrak{X},
\mathscr{O}_{\mathfrak{X}}/\mathscr{J}_\lambda)\to
(\mathfrak{X},\mathscr{O}_{\mathfrak{X}}/\mathscr{J}_\mu)$ de préschémas
(usuels) tel que l'on ait $f_\lambda=f_\mu\circ f_{\mu\lambda}$. Les
préschémas
$X_\lambda=(\mathfrak{X},
\mathscr{O}_{\mathfrak{X}}/\mathscr{J}_\lambda)$ et les morphismes
$f_{\mu\lambda}$ constituent donc (en vertu de
\hyperref[I.10.4.8-fr]{10.4.8}) un \emph{système inductif} dans la
catégorie des préschémas formels.
\end{env}

\begin{proposition}[10.6.2]
\label{I.10.6.2-fr}
Avec les notations de (\hyperref[I.10.6.1-fr]{10.6.1}), le préschéma
formel $\mathfrak{X}$ et les morphismes $f_\lambda$ constituent une
limite inductive (T, I, 1.8) du système $(X_\lambda,f_{\mu\lambda})$
dans la catégorie des préschémas formels.
\end{proposition}

\begin{proof}
Soit $\mathfrak{Y}$ un préschéma formel, et pour chaque indice $\lambda$,
soit
\[
  g_\lambda=(\psi_\lambda,\theta_\lambda):X_\lambda\to\mathfrak{Y}
\]
un morphisme, tel que l'on ait
$g_\lambda=g_\mu\circ f_{\mu\lambda}$ pour
$\mathscr{J}_\mu\subset\mathscr{J}_\lambda$. Cette dernière condition
et la définition des $X_\lambda$ entraînent d'abord que tous les
$\psi_\lambda$ sont identiques à une même application continue
$\psi:\mathfrak{X}\to\mathfrak{Y}$ des espaces sous-jacents~; en outre,
les homomorphismes
$\theta_\lambda^\sharp:\psi^*(\mathscr{O}_{\mathfrak{Y}})\to
\mathscr{O}_{X_\lambda}=
\mathscr{O}_{\mathfrak{X}}/\mathscr{J}_\lambda$ forment un
\emph{système projectif} d'homomorphismes de faisceaux d'anneaux. Par
passage à la limite projective, on en déduit donc un homomorphisme
\[
  \omega:\psi^*(\mathscr{O}_{\mathfrak{Y}})\to
  \varprojlim\mathscr{O}_{\mathfrak{X}}/\mathscr{J}_\lambda
  =\mathscr{O}_{\mathfrak{X}},
\]
et il est clair que le morphisme $g=(\psi,\omega^b)$ d'\emph{espaces
annelés} est le \emph{seul} rendant commutatif les diagrammes
\[
  \label{I.10.6.2.1-fr}
  \xymatrix{
    X_\lambda\ar[rr]^{g_\lambda}\ar[rd]_{f_\lambda} &&
    \mathfrak{Y}\\
    & \mathfrak{X}\ar[ru]_g
  }
  \tag{10.6.2.1}
\]
Il reste donc à prouver que $g$ est un morphisme de \emph{préschémas
formels}~; la question étant locale sur $\mathfrak{X}$ et
$\mathfrak{Y}$, on peut donc supposer
$\mathfrak{X}=\operatorname{Spf}(A)$,
$\mathfrak{Y}=\operatorname{Spf}(B)$, $A$ et $B$ étant des anneaux
admissibles, avec
$\mathscr{J}_\lambda=\mathfrak{J}_\lambda^{\Delta}$, où
$(\mathfrak{J}_\lambda)$ est un système fondamental d'idéaux de
définition de $A$ (\hyperref[I.10.3.5-fr]{10.3.5})~; comme
$A=\varprojlim A/\mathfrak{J}_\lambda$, l'existence d'un morphisme de
schémas formels affines $g$ rendant commutatifs les diagrammes
(\hyperref[I.10.6.2.1-fr]{10.6.2.1}) résulte alors de la correspondance
biunivoque (\hyperref[I.10.2.2-fr]{10.2.2}) entre morphismes de schémas
formels affines et homomorphismes continus d'anneaux, et de la définition
de la limite projective. Mais l'unicité de $g$ en tant que morphisme
d'espaces annelés montre qu'il coïncide avec le morphisme noté de même au
début de la démonstration.
\end{proof}

La proposition suivante établit, sous certaines conditions
supplémentaires, l'existence de la limite inductive d'un système inductif
donné de préschémas (usuels) dans la catégorie des préschémas formels~:

\begin{proposition}[10.6.3]
\label{I.10.6.3-fr}
Soient $\mathfrak{X}$ un espace topologique,
$(\mathscr{O}_i,u_{ji})$ un système projectif de faisceaux d'anneaux sur
$\mathfrak{X}$, ayant $\mathbf{N}$ pour ensemble d'indices. Soit
$\mathscr{J}_i$ le noyau de
$u_{0i}:\mathscr{O}_i\to\mathscr{O}_0$. On suppose que~:
\begin{enumerate}
  \item[a)] L'espace annelé $(\mathfrak{X},\mathscr{O}_i)$ est un
    préschéma $X_i$.
  \item[b)] Pour tout $x\in X$ et tout $i$, il existe un voisinage
    ouvert $U_i$ de $x$ dans $\mathfrak{X}$ tel que la restriction
    $\mathscr{J}_i|U_i$ soit nilpotente.
  \item[c)] Les homomorphismes $u_{ji}$ sont surjectifs.
\end{enumerate}

\oldpage[I]{190}

Soit $\mathscr{O}_{\mathfrak{X}}$ le faisceau d'anneaux topologiques
limite projective des faisceaux d'anneaux pseudo-discrets
$\mathscr{O}_i$, et soit
$u_i:\mathscr{O}_{\mathfrak{X}}\to\mathscr{O}_i$ l'homomorphisme
canonique. Alors l'espace topologiquement annelé
$(\mathfrak{X},\mathscr{O}_{\mathfrak{X}})$ est un préschéma formel~;
les homomorphismes $u_i$ sont surjectifs~; leurs noyaux
$\mathscr{J}^{(i)}$ forment un système fondamental d'Idéaux de définition
de $\mathfrak{X}$, et $\mathscr{J}^{(0)}$ est la limite projective des
faisceaux d'idéaux $\mathscr{J}_i$.
\end{proposition}

\begin{proof}
Notons d'abord que sur chaque fibre, $u_{ji}$ est un homomorphisme
surjectif et \emph{a fortiori} un homomorphisme local~; donc
$v_{ij}=(1_{\mathfrak{X}},u_{ji})$ est un morphisme de préschémas
$X_j\to X_i$ ($i\geq j$) (\hyperref[I.2.2.1-fr]{2.2.1}). Supposons
d'abord que chaque $X_i$ soit un schéma affine d'anneau $A_i$. Il existe
un homomorphisme \emph{d'anneaux}
$\varphi_{ji}:A_i\to A_j$ tel que
$u_{ji}=\widetilde{\varphi_{ji}}$
(\hyperref[I.1.7.3-fr]{1.7.3})~; par suite
(\hyperref[I.1.6.3-fr]{1.6.3}), le faisceau $\mathscr{O}_j$ est un
$\mathscr{O}_i$-Module quasi-cohérent sur $X_i$ (pour la loi externe
définie par $u_{ji}$), associé à $A_j$ considéré comme $A_i$-module à
l'aide de $\varphi_{ji}$. Pour tout $f\in A_i$, soit
$f'=\varphi_{ji}(f)$~; par hypothèse, les ouverts $D(f)$ et $D(f')$ sont
identiques dans $\mathfrak{X}$, et l'homomorphisme de
$\Gamma(D(f),\mathscr{O}_i)=(A_i)_f$ dans
$\Gamma(D(f),\mathscr{O}_j)=(A_j)_{f'}$ correspondant à $u_{ji}$ n'est
autre que $(\varphi_{ji})_f$ (\hyperref[I.1.6.1-fr]{1.6.1}). Mais
lorsqu'on considère $A_j$ comme $A_i$-module, $(A_j)_{f'}$ est le
$(A_i)_f$-module $(A_j)_f$, donc on a aussi
$u_{ji}=\widetilde{\varphi_{ji}}$ lorsque $\varphi_{ji}$ est cette fois
considéré comme homomorphisme de $A_i$-modules. Alors, comme $u_{ji}$ est
surjectif, on en conclut que $\varphi_{ji}$ l'est aussi
(\hyperref[I.1.3.9-fr]{1.3.9}) et si $\mathfrak{J}_{ji}$ est le noyau de
$\varphi_{ji}$, le noyau de $u_{ji}$ est un $\mathscr{O}_i$-Module
quasi-cohérent égal à $\widetilde{\mathfrak{J}_{ji}}$. En particulier,
on a $\mathscr{J}_i=\widetilde{\mathfrak{J}_i}$, où $\mathfrak{J}_i$ est
le noyau de $\varphi_{0i}:A_i\to A_0$. L'hypothèse b) entraîne que
$\mathscr{J}_i$ est \emph{nilpotent}~: en effet, comme $\mathfrak{X}$
est quasi-compact, on peut recouvrir $\mathfrak{X}$ par un nombre fini
d'ouverts $U_k$ tels que $(\mathscr{J}_i|U_k)^{n_k}=0$ et en prenant pour
$n$ le plus grand des $n_k$, on a $\mathscr{J}_i^n=0$. On en conclut que
$\mathfrak{J}_i$ est nilpotent (\hyperref[I.1.3.13-fr]{1.3.13}). Alors
l'anneau $A=\varprojlim A_i$ est admissible
(\hyperref[0.7.2.2-fr]{0, 7.2.2}), l'homomorphisme canonique
$\varphi_i:A\to A_i$ est surjectif et son noyau
$\mathfrak{J}^{(i)}$ est égal à la limite projective des
$\mathfrak{J}_{ik}$ pour $k\geq i$~; les $\mathfrak{J}^{(i)}$ forment
un système fondamental de voisinages de 0 dans $A$. Les assertions de
(\hyperref[I.10.6.3-fr]{10.6.3}) résultent dans ce cas de
(\hyperref[I.10.1.1-fr]{10.1.1}) et
(\hyperref[I.10.3.2-fr]{10.3.2}),
$(\mathfrak{X},\mathscr{O}_{\mathfrak{X}})$ n'étant autre que
$\operatorname{Spf}(A)$.

Toujours dans ce même cas particulier, notons que si $f=(f_i)$ est un
élément de la limite projective $A=\varprojlim A_i$, tous les ouverts
$D(f_i)$ (ouvert affine dans $X_i$) s'identifient à l'ouvert
$\mathfrak{D}(f)$ de $\mathfrak{X}$, le préschéma induit par $X_i$ sur
$\mathfrak{D}(f)$ s'identifiant donc au schéma affine
$\operatorname{Spec}((A_i)_{f_i})$.

Dans le cas général, remarquons d'abord que pour tout ouvert quasi-compact
$U$ de $\mathfrak{X}$, chacun des $\mathscr{J}_i|U$ est nilpotent, comme
le montre le raisonnement fait ci-dessus. Nous allons voir que pour tout
$x\in\mathfrak{X}$, il y a un voisinage ouvert $U$ de $x$ dans
$\mathfrak{X}$ qui est un \emph{ouvert affine pour tous les $X_i$}. En
effet, prenons $U$ ouvert affine pour $X_0$, et observons que
$\mathscr{O}_{X_0}=\mathscr{O}_{X_i}/\mathscr{J}_i$. Comme
$\mathscr{J}_i|U$ est nilpotent, en vertu de ce qui précède, $U$ est
ouvert affine aussi pour chaque $X_i$ en vertu de
(\hyperref[I.5.1.9-fr]{5.1.9}). Cela étant, pour tout $U$ satisfaisant
aux conditions précédentes, l'étude du cas affine faite plus haut montre
que $(U,\mathscr{O}_{\mathfrak{X}}|U)$ est un préschéma formel dont les
$\mathscr{J}^{(i)}|U$ forment un système fondamental d'idéaux de
définition et $\mathscr{J}^{(0)}|U$ est limite projective des
$\mathscr{J}_i|U$~; d'où la conclusion.
\end{proof}

\begin{corollary}[10.6.4]
\label{I.10.6.4-fr}
Supposons que pour $i\geq j$, le noyau de $u_{ji}$ soit
$\mathscr{J}_i^{j+1}$ et que $\mathscr{J}_1/\mathscr{J}_1^2$

\oldpage[I]{191}

soit de type fini sur
$\mathscr{O}_0=\mathscr{O}_1/\mathscr{J}_1$. Alors $\mathfrak{X}$ est
un préschéma formel adique, et si $\mathscr{J}^{(n)}$ est le noyau de
$\mathscr{O}_{\mathfrak{X}}\to\mathscr{O}_n$, on a
$\mathscr{J}^{(n)}=\mathscr{J}^{n+1}$ et
$\mathscr{J}/\mathscr{J}^2$ est isomorphe à $\mathscr{J}_1$. Si en
outre $X_0$ est localement noethérien (resp. noethérien),
$\mathfrak{X}$ est localement noethérien (resp. noethérien).
\end{corollary}

\begin{proof}
Comme les espaces sous-jacents à $\mathfrak{X}$ et $X_0$ sont les mêmes,
la question est locale et l'on peut supposer tous les $X_i$ affines~;
compte tenu des relations
$\mathscr{J}_{ij}=\widetilde{\mathfrak{J}_{ji}}$ (avec les notations
de (\hyperref[I.10.6.3-fr]{10.6.3})), on est aussitôt ramené aux
assertions correspondantes de
(\hyperref[0.7.2.7-fr]{0, 7.2.7} et
\hyperref[0.7.2.8-fr]{7.2.8}), en notant que
$\mathfrak{J}_1/\mathfrak{J}_1^2$ est alors un $A_0$-module de type
fini (\hyperref[I.1.3.9-fr]{1.3.9}).
\end{proof}

En particulier, \emph{tout préschéma formel localement noethérien
$\mathfrak{X}$} est limite inductive d'une suite $(X_n)$ de préschémas
(usuels) localement noethériens vérifiant les conditions de
(\hyperref[I.10.6.3-fr]{10.6.3}) et
(\hyperref[I.10.6.4-fr]{10.6.4})~: il suffit de considérer un Idéal de
définition $\mathscr{J}$ de $\mathfrak{X}$
(\hyperref[I.10.5.4-fr]{10.5.4}) et de prendre
$X_n=(\mathfrak{X},
\mathscr{O}_{\mathfrak{X}}/\mathscr{J}^{n+1})$
((\hyperref[I.10.5.1-fr]{10.5.1}) et
(\hyperref[I.10.6.2-fr]{10.6.2})).

\begin{corollary}[10.6.5]
\label{I.10.6.5-fr}
Soit $A$ un anneau admissible. Pour que le schéma formel affine
$\mathfrak{X}=\operatorname{Spf}(A)$ soit noethérien, il faut et il
suffit que $A$ soit adique et noethérien.
\end{corollary}

\begin{proof}
La condition est évidemment suffisante. Inversement, supposons que
$\mathfrak{X}$ soit noethérien, et soient $\mathfrak{J}$ un idéal de
définition de $A$, $\mathscr{J}=\mathfrak{J}^\Delta$ l'Idéal de
définition correspondant de $\mathfrak{X}$. Les préschémas (usuels)
$X_n=(\mathfrak{X},
\mathscr{O}_{\mathfrak{X}}/\mathscr{J}^{n+1})$ sont alors affines et
noethériens, donc les anneaux
$A_n=A/\mathfrak{J}^{n+1}$ sont noethériens
(\hyperref[I.6.1.3-fr]{6.1.3}), d'où on conclut que
$\mathfrak{J}/\mathfrak{J}^2$ est un $A/\mathfrak{J}$-module de type
fini. Comme les $\mathscr{J}^n$ forment un système fondamental d'Idéaux
de définition de $\mathfrak{X}$
(\hyperref[I.10.5.1-fr]{10.5.1}), on a
$\mathscr{O}_{\mathfrak{X}}
=\varprojlim(\mathscr{O}_{\mathfrak{X}}/\mathscr{J}^n)$
(\hyperref[I.10.5.3-fr]{10.5.3})~; on en conclut
(\hyperref[I.10.1.3-fr]{10.1.3}) que $A$ est topologiquement isomorphe à
$\varprojlim A/\mathfrak{J}^n$, donc est adique et noethérien
(\hyperref[0.7.2.8-fr]{0, 7.2.8}).
\end{proof}

\begin{remark}[10.6.6]
\label{I.10.6.6-fr}
Avec les notations de (\hyperref[I.10.6.3-fr]{10.6.3}), soit
$\mathscr{F}_i$ un $\mathscr{O}_i$-Module, et supposons donné, pour
$i\geq j$, un $v_{ij}$-morphisme
$\theta_{ji}:\mathscr{F}_i\to\mathscr{F}_j$, de sorte que
$\theta_{kj}\circ\theta_{ji}=\theta_{ki}$ pour $k\leq j\leq i$.
Comme l'application continue sous-jacente à $v_{ij}$ est l'identité,
$\theta_{ji}$ est un homomorphisme de faisceaux de groupes abéliens sur
l'espace $\mathfrak{X}$~; en outre, si $\mathscr{F}$ est la limite
projective du système projectif $(\mathscr{F}_i)$ de faisceaux de groupes
abéliens, le fait que les $\theta_{ji}$ sont des $v_{ij}$-morphismes
permet de définir sur $\mathscr{F}$ une structure de
$\mathscr{O}_{\mathfrak{X}}$-Module par passage à la limite projective~;
muni de cette structure, nous dirons que $\mathscr{F}$ est la
\emph{limite projective} (pour les $\theta_{ji}$) du système de
$\mathscr{O}_i$-Modules $(\mathscr{F}_i)$. Dans le cas particulier où
$v_{ij}^*(\mathscr{F}_i)=\mathscr{F}_j$ et où $\theta_{ji}$ est
l'\emph{identité}, nous dirons pour abréger, que $\mathscr{F}$ est la
limite projective d'un système $(\mathscr{F}_i)$ tel que
$v_{ij}^*(\mathscr{F}_i)=\mathscr{F}_j$ pour $j\leq i$ (sans
mentionner les $\theta_{ji}$).
\end{remark}

\begin{env}[10.6.7]
\label{I.10.6.7-fr}
Soient $\mathfrak{X}$, $\mathfrak{Y}$ deux préschémas formels,
$\mathscr{J}$ (resp. $\mathscr{K}$) un Idéal de définition de
$\mathfrak{X}$ (resp. $\mathfrak{Y}$),
$f:\mathfrak{X}\to\mathfrak{Y}$ un morphisme tel que
$f^*(\mathscr{K})\mathscr{O}_{\mathfrak{X}}\subset\mathscr{J}$. On a
alors pour tout entier $n>0$,
$f^*(\mathscr{K}^n)\mathscr{O}_{\mathfrak{X}}
=(f^*(\mathscr{K})\mathscr{O}_{\mathfrak{X}})^n
\subset\mathscr{J}^n$~; on peut donc
(\hyperref[I.10.5.6-fr]{10.5.6}) déduire de $f$ un morphisme de
préschémas (usuels) $f_n:X_n\to Y_n$, en posant
$X_n=(\mathfrak{X},
\mathscr{O}_{\mathfrak{X}}/\mathscr{J}^{n+1})$,
$Y_n=(\mathfrak{Y},
\mathscr{O}_{\mathfrak{Y}}/\mathscr{K}^{n+1})$, et il résulte
aussitôt des définitions que les diagrammes
\[
  \label{I.10.6.7.1-fr}
  \xymatrix{
    X_m\ar[r]^{f_m}\ar[d] &
    Y_m\ar[d]\\
    X_n\ar[r]_{f_n} &
    Y_n
  }
  \tag{10.6.7.1}
\]

\oldpage[I]{192}

sont commutatifs pour $m\leq n$~; autrement dit, $(f_n)$ est un
\emph{système inductif} de morphismes.
\end{env}

\begin{env}[10.6.8]
\label{I.10.6.8-fr}
Inversement, soit $(X_n)$ (resp. $(Y_n)$) un système inductif de
préschémas (usuels) satisfaisant aux conditions b) et c) de
(\hyperref[I.10.6.3-fr]{10.6.3}), et soit $\mathfrak{X}$ (resp.
$\mathfrak{Y}$) sa limite inductive. Par définition des limites
inductives, toute suite $(f_n)$ de morphismes $X_n\to Y_n$ formant un
système inductif admet une \emph{limite inductive}
$f:\mathfrak{X}\to\mathfrak{Y}$, qui est l'unique morphisme de
préschémas formels rendant commutatifs les diagrammes
\[
  \xymatrix{
    X_n\ar[r]^{f_n}\ar[d] &
    Y_n\ar[d]\\
    \mathfrak{X}\ar[r]_f &
    \mathfrak{Y}
  }
\]
\end{env}

\begin{proposition}[10.6.9]
\label{I.10.6.9-fr}
Soient $\mathfrak{X}$, $\mathfrak{Y}$ deux préschémas formels
localement noethériens, $\mathscr{J}$ (resp. $\mathscr{K}$) un Idéal
de définition de $\mathfrak{X}$ (resp. $\mathfrak{Y}$)~;
l'application $f\mapsto(f_n)$ définie dans
(\hyperref[I.10.6.7-fr]{10.6.7}) est une bijection de l'ensemble des
morphismes $f:\mathfrak{X}\to\mathfrak{Y}$ tels que
$f^*(\mathscr{K})\mathscr{O}_{\mathfrak{X}}\subset\mathscr{J}$, sur
l'ensemble des suites $(f_n)$ de morphismes rendant commutatifs les
diagrammes (\hyperref[I.10.6.7.1-fr]{10.6.7.1}).
\end{proposition}

\begin{proof}
Si $f$ est la limite inductive d'une telle suite, il faut montrer que
$f^*(\mathscr{K})\mathscr{O}_{\mathfrak{X}}\subset\mathscr{J}$. La
question étant locale sur $\mathfrak{X}$ et $\mathfrak{Y}$, on peut se
borner au cas où $\mathfrak{X}=\operatorname{Spf}(A)$,
$\mathfrak{Y}=\operatorname{Spf}(B)$ sont affines, $A$ et $B$ étant
adiques noethériens,
$\mathscr{J}=\mathfrak{J}^\Delta$,
$\mathscr{K}=\mathfrak{K}^\Delta$, où $\mathfrak{J}$ (resp.
$\mathfrak{K}$) est un idéal de définition de $A$ (resp. $B$). On a
alors $X_n=\operatorname{Spec}(A_n)$,
$Y_n=\operatorname{Spec}(B_n)$, avec
$A_n=A/\mathfrak{J}^{n+1}$ et $B_n=B/\mathfrak{K}^{n+1}$, en vertu de
(\hyperref[I.10.3.6-fr]{10.3.6}) et
(\hyperref[I.10.3.2-fr]{10.3.2})~;
$f_n=({}^a\varphi_n,\widetilde{\varphi_n})$, où les homomorphismes
$\varphi_n:B_n\to A_n$ forment un système projectif, donc
$f=({}^a\varphi,\widetilde{\varphi})$, où
$\varphi=\varprojlim\varphi_n$. La commutativité du diagramme
(\hyperref[I.10.6.7.1-fr]{10.6.7.1}) pour $m=0$ donne alors la
condition
$\varphi_n(\mathfrak{K}/\mathfrak{K}^{n+1})
\subset\mathfrak{J}/\mathfrak{J}^{n+1}$ pour tout $n$, donc, en
passant à la limite projective, $\varphi(\mathfrak{K})\subset
\mathfrak{J}$, et cela entraîne
$f^*(\mathscr{K})\mathscr{O}_{\mathfrak{X}}\subset\mathscr{J}$
(\hyperref[I.10.5.6-fr]{10.5.6, (ii)}).
\end{proof}

\begin{corollary}[10.6.10]
\label{I.10.6.10-fr}
Soient $\mathfrak{X}$, $\mathfrak{Y}$ deux préschémas formels
localement noethériens, $\mathscr{T}$ le plus grand Idéal de définition
de $\mathfrak{X}$ (\hyperref[I.10.5.4-fr]{10.5.4}).
\begin{enumerate}
  \item[\rm{(i)}] Pour tout Idéal de définition $\mathscr{K}$ de
  $\mathfrak{Y}$ et tout morphisme
  $f:\mathfrak{X}\to\mathfrak{Y}$, on a
  $f^*(\mathscr{K})\mathscr{O}_{\mathfrak{X}}\subset\mathscr{T}$.
  \item[\rm{(ii)}] Il y a correspondance biunivoque canonique entre
  $\operatorname{Hom}(\mathfrak{X},\mathfrak{Y})$ et l'ensemble des
  suites $(f_n)$ de morphismes rendant commutatifs les diagrammes
  (\hyperref[I.10.6.7.1-fr]{10.6.7.1}), où
  $X_n=(\mathfrak{X},
  \mathscr{O}_{\mathfrak{X}}/\mathscr{T}^{n+1})$,
  $Y_n=(\mathfrak{Y},
  \mathscr{O}_{\mathfrak{Y}}/\mathscr{K}^{n+1})$.
\end{enumerate}
\end{corollary}

\begin{proof}
(ii) résulte aussitôt de (i) et de
(\hyperref[I.10.6.9-fr]{10.6.9}). Pour démontrer (i), on peut se borner
au cas où $\mathfrak{X}=\operatorname{Spf}(A)$,
$\mathfrak{Y}=\operatorname{Spf}(B)$, $A$ et $B$ étant noethériens,
$\mathscr{T}=\mathfrak{T}^\Delta$,
$\mathscr{K}=\mathfrak{K}^\Delta$, où $\mathfrak{T}$ est le plus grand
idéal de définition de $A$ et $\mathfrak{K}$ un idéal de définition de
$B$. Soit $f=({}^a\varphi,\widetilde{\varphi})$, où
$\varphi:B\to A$ est un homomorphisme continu~; comme les éléments de
$\mathfrak{K}$ sont topologiquement nilpotents
(\hyperref[0.7.1.4-fr]{0, 7.1.4, (ii)}), il en est de même de ceux de
$\varphi(\mathfrak{K})$, donc
$\varphi(\mathfrak{K})\subset\mathfrak{T}$ puisque $\mathfrak{T}$ est
l'ensemble des éléments topologiquement nilpotents de $A$
(\hyperref[0.7.1.6-fr]{0, 7.1.6})~; d'où la conclusion en vertu de
(\hyperref[I.10.5.6-fr]{10.5.6, (ii)}).
\end{proof}

\begin{corollary}[10.6.11]
\label{I.10.6.11-fr}
Soient $\mathfrak{S}$, $\mathfrak{X}$, $\mathfrak{Y}$ trois
préschémas formels localement noethériens,
$f:\mathfrak{X}\to\mathfrak{S}$,
$g:\mathfrak{Y}\to\mathfrak{S}$ des morphismes faisant de
$\mathfrak{X}$ et $\mathfrak{Y}$ des $\mathfrak{S}$-préschémas
formels. Soit $\mathscr{J}$ (resp. $\mathscr{K}$, $\mathscr{L}$) un
Idéal de définition de $\mathfrak{S}$ (resp. $\mathfrak{X}$,
$\mathfrak{Y}$), et supposons que
$f^*(\mathscr{J})\mathscr{O}_{\mathfrak{X}}\subset\mathscr{K}$,
$g^*(\mathscr{J})\mathscr{O}_{\mathfrak{Y}}=\mathscr{L}$~; posons
$S_n=(\mathfrak{S},
\mathscr{O}_{\mathfrak{S}}/\mathscr{J}^{n+1})$,
$X_n=(\mathfrak{X},
\mathscr{O}_{\mathfrak{X}}/\mathscr{K}^{n+1})$,
$Y_n=(\mathfrak{Y},
\mathscr{O}_{\mathfrak{Y}}/\mathscr{L}^{n+1})$. Il y a alors
correspondance

\oldpage[I]{193}

biunivoque canonique entre
$\operatorname{Hom}_{\mathfrak{S}}(\mathfrak{X},\mathfrak{Y})$ et
l'ensemble des suites $(u_n)$ de $S_n$-morphismes
$u_n:X_n\to Y_n$ rendant commutatifs les diagrammes
(\hyperref[I.10.6.7.1-fr]{10.6.7.1}).
\end{corollary}

\begin{proof}
Pour tout $\mathfrak{S}$-morphisme
$u:\mathfrak{X}\to\mathfrak{Y}$, on a par définition $f=g\circ u$, donc
\[
  u^*(\mathscr{L})\mathscr{O}_{\mathfrak{X}}
  =u^*(g^*(\mathscr{J})\mathscr{O}_{\mathfrak{Y}})
    \mathscr{O}_{\mathfrak{X}}
  =f^*(\mathscr{J})\mathscr{O}_{\mathfrak{X}}
  \subset\mathscr{K}
\]
le corollaire découle donc de
(\hyperref[I.10.6.9-fr]{10.6.9}).
\end{proof}

On notera que, pour $m\leq n$, la donnée d'un morphisme
$f_n:X_n\to Y_n$ détermine un morphisme et un seul
$f_m:X_m\to Y_m$ rendant commutatif le diagramme
(\hyperref[I.10.6.7.1-fr]{10.6.7.1}), comme on le voit aussitôt en se
ramenant au cas affine~; on a ainsi défini une application
$\varphi_{mn}:\operatorname{Hom}_{S_n}(X_n,Y_n)
\to\operatorname{Hom}_{S_m}(X_m,Y_m)$ et les
$\operatorname{Hom}_{S_n}(X_n,Y_n)$ forment pour les $\varphi_{mn}$ un
\emph{système projectif d'ensembles}~;
(\hyperref[I.10.6.11-fr]{10.6.11}) s'énonce encore en disant qu'il existe
une bijection canonique
\[
  \operatorname{Hom}_{\mathfrak{S}}(\mathfrak{X},\mathfrak{Y})
  \xrightarrow{\sim}
  \varprojlim_n\operatorname{Hom}_{S_n}(X_n,Y_n).
\]

\subsection{Produit de préschémas formels.}
\label{subsection:I.10.7-fr}

\begin{env}[10.7.1]
\label{I.10.7.1-fr}
Soit $\mathfrak{S}$ un préschéma formel~; les $\mathfrak{S}$-préschémas
formels formant une catégorie, on peut définir la notion de
\emph{produit} de $\mathfrak{S}$-préschémas formels.
\end{env}

\begin{proposition}[10.7.2]
\label{I.10.7.2-fr}
Soient $\mathfrak{X}=\operatorname{Spf}(B)$,
$\mathfrak{Y}=\operatorname{Spf}(C)$ deux schémas formels affines
au-dessus d'un schéma formel affine
$\mathfrak{S}=\operatorname{Spf}(A)$. Soient
$\mathfrak{Z}=\operatorname{Spf}(B\widehat{\otimes}_A C)$,
$p_1$, $p_2$ les $\mathfrak{S}$-morphismes correspondant
(\hyperref[I.10.2.2-fr]{10.2.2}) aux $A$-homomorphismes canoniques
(continus) $\rho$ et $\sigma$ de $B$ et $C$ dans
$B\widehat{\otimes}_A C$~; alors $(\mathfrak{Z},p_1,p_2)$ est un produit
des $\mathfrak{S}$-schémas formels affines $\mathfrak{X}$ et
$\mathfrak{Y}$.
\end{proposition}

\begin{proof}
En vertu de (\hyperref[I.10.4.6-fr]{10.4.6}), tout revient à vérifier que
si, à tout $A$-homomorphisme continu
$\varphi:B\widehat{\otimes}_A C\to D$, où $D$ est un anneau admissible qui
est une $A$-algèbre topologique, on associe le couple
$(\varphi\circ\rho,\varphi\circ\sigma)$, on définit une bijection
\[
  \operatorname{Hom}_A(B\widehat{\otimes}_A C,D)
  \xrightarrow{\sim}
  \operatorname{Hom}_A(B,D)\times\operatorname{Hom}_A(C,D)
\]
ce qui n'est autre que la propriété universelle du produit tensoriel
complété (\hyperref[0.7.7.6-fr]{0, 7.7.6}).
\end{proof}

\begin{proposition}[10.7.3]
\label{I.10.7.3-fr}
Étant donnés deux $\mathfrak{S}$-préschémas formels
$\mathfrak{X}$, $\mathfrak{Y}$, le produit
$\mathfrak{X}\times_{\mathfrak{S}}\mathfrak{Y}$ existe.
\end{proposition}

\begin{proof}
La démonstration est identique à celle de
(\hyperref[I.3.2.6-fr]{3.2.6}), en y remplaçant les schémas affines
(resp. les ouverts affines) par les schémas formels affines (resp. les
ouverts formels affines), et la prop.
(\hyperref[I.3.2.2-fr]{3.2.2}) par
(\hyperref[I.10.7.2-fr]{10.7.2}).
\end{proof}

Toutes les propriétés formelles du produit de préschémas
(\hyperref[I.3.2.7-fr]{3.2.7} et
\hyperref[I.3.2.8-fr]{3.2.8},
\hyperref[I.3.3.1-fr]{3.3.1} à
\hyperref[I.3.3.12-fr]{3.3.12}) sont valables sans aucune modification
pour le produit de préschémas formels.

\begin{env}[10.7.4]
\label{I.10.7.4-fr}
Soient $\mathfrak{S}$, $\mathfrak{X}$, $\mathfrak{Y}$ trois préschémas
formels et soient $f:\mathfrak{X}\to\mathfrak{S}$,
$g:\mathfrak{Y}\to\mathfrak{S}$ deux morphismes. Supposons qu'il existe
dans $\mathfrak{S}$, $\mathfrak{X}$, $\mathfrak{Y}$ respectivement,
trois systèmes fondamentaux d'Idéaux de définition
$(\mathscr{J}_\lambda)$, $(\mathscr{K}_\lambda)$,
$(\mathscr{L}_\lambda)$ respectivement, ayant même ensemble d'indices
$I$, tels que
$f^*(\mathscr{J}_\lambda)\mathscr{O}_{\mathfrak{X}}
\subset\mathscr{K}_\lambda$ et
$g^*(\mathscr{J}_\lambda)\mathscr{O}_{\mathfrak{Y}}
\subset\mathscr{L}_\lambda$ pour tout $\lambda$. Posons
$S_\lambda=(\mathfrak{S},
\mathscr{O}_{\mathfrak{S}}/\mathscr{J}_\lambda)$,
$X_\lambda=(\mathfrak{X},
\mathscr{O}_{\mathfrak{X}}/\mathscr{K}_\lambda)$,
$Y_\lambda=(\mathfrak{Y},
\mathscr{O}_{\mathfrak{Y}}/\mathscr{L}_\lambda)$~; pour
$\mathscr{J}_\mu\subset\mathscr{J}_\lambda$,
$\mathscr{K}_\mu\subset\mathscr{K}_\lambda$,
$\mathscr{L}_\mu\subset\mathscr{L}_\lambda$, notons que
$S_\lambda$ (resp. $X_\lambda$, $Y_\lambda$) est un sous-préschéma fermé
de $S_\mu$ (resp. $X_\mu$, $Y_\mu$) ayant même

\oldpage[I]{194}

espace sous-jacent (\hyperref[I.10.6.1-fr]{10.6.1}). Comme
$S_\lambda\to S_\mu$ est un monomorphisme de préschémas, on voit donc
d'abord que les produits
$X_\lambda\times_{S_\lambda}Y_\lambda$ et
$X_\lambda\times_{S_\mu}Y_\lambda$ sont identiques
(\hyperref[I.3.2.4-fr]{3.2.4}), puis que
$X_\lambda\times_{S_\mu}Y_\lambda$ s'identifie à un sous-préschéma fermé
de $X_\mu\times_{S_\mu}Y_\mu$ ayant même espace sous-jacent
(\hyperref[I.4.3.1-fr]{4.3.1}). Cela étant, le produit
$\mathfrak{X}\times_{\mathfrak{S}}\mathfrak{Y}$ est la
\emph{limite inductive} des préschémas usuels
$X_\lambda\times_{S_\lambda}Y_\lambda$~: en effet, on voit comme dans
(\hyperref[I.10.6.2-fr]{10.6.2}) qu'on peut se ramener au cas où
$\mathfrak{S}$, $\mathfrak{X}$ et $\mathfrak{Y}$ sont des schémas formels
affines. Compte tenu de (\hyperref[I.10.5.6-fr]{10.5.6, (ii)}) et de
l'hypothèse sur les systèmes fondamentaux d'Idéaux de définition de
$\mathfrak{S}$, $\mathfrak{X}$ et $\mathfrak{Y}$, on voit aussitôt que
notre assertion résulte de la définition du produit tensoriel complété de
deux algèbres (\hyperref[0.7.7.1-fr]{0, 7.7.1}).

En outre, soit $\mathfrak{Z}$ un $\mathfrak{S}$-préschéma formel,
$(\mathscr{M}_\lambda)$ un système fondamental d'Idéaux de définition de
$\mathfrak{Z}$ ayant $I$ comme ensemble d'indices,
$u:\mathfrak{Z}\to\mathfrak{X}$,
$v:\mathfrak{Z}\to\mathfrak{Y}$ deux $\mathfrak{S}$-morphismes tels que
$u^*(\mathscr{K}_\lambda)\mathscr{O}_{\mathfrak{Z}}
\subset\mathscr{M}_\lambda$ et
$v^*(\mathscr{L}_\lambda)\mathscr{O}_{\mathfrak{Z}}
\subset\mathscr{M}_\lambda$. Si l'on pose
$Z_\lambda=(\mathfrak{Z},
\mathscr{O}_{\mathfrak{Z}}/\mathscr{M}_\lambda)$, et si
$u_\lambda:Z_\lambda\to X_\lambda$ et
$v_\lambda:Z_\lambda\to Y_\lambda$ sont les $S_\lambda$-morphismes
correspondant à $u$ et $v$ (\hyperref[I.10.5.6-fr]{10.5.6}), on vérifie
aussitôt que $(u,v)_{\mathfrak{S}}$ est la limite inductive des
$S_\lambda$-morphismes $(u_\lambda,v_\lambda)_{S_\lambda}$.

Les considérations de ce numéro s'appliquent en particulier lorsque
$\mathfrak{S}$, $\mathfrak{X}$ et $\mathfrak{Y}$ sont localement
noethériens, en prenant comme systèmes fondamentaux d'Idéaux de définition
les systèmes formés des puissances d'un Idéal de définition
(\hyperref[I.10.5.1-fr]{10.5.1}). Mais on notera que
$\mathfrak{X}\times_{\mathfrak{S}}\mathfrak{Y}$ n'est pas nécessairement
localement noethérien (voir toutefois
(\hyperref[I.10.13.5-fr]{10.13.5})).
\end{env}

\subsection{Complété formel d'un préschéma le long d'une partie fermée.}
\label{subsection:I.10.8-fr}

\begin{env}[10.8.1]
\label{I.10.8.1-fr}
Soient $X$ un préschéma (usuel) localement noethérien, $X'$ une partie
fermée de l'espace sous-jacent à $X$~; désignons par $\Phi$ l'ensemble des
faisceaux cohérents d'idéaux $\mathscr{J}$ dans $\mathscr{O}_X$, tels que
le support de $\mathscr{O}_X/\mathscr{J}$ soit $X'$. L'ensemble $\Phi$
n'est pas vide (\hyperref[I.5.2.1-fr]{5.2.1},
\hyperref[I.4.1.4-fr]{4.1.4} et
\hyperref[I.6.1.1-fr]{6.1.1})~; nous l'ordonnerons par la relation
$\supset$.
\end{env}

\begin{lemma}[10.8.2]
\label{I.10.8.2-fr}
L'ensemble ordonné $\Phi$ est filtrant~; si $X$ est noethérien, pour tout
$\mathscr{J}_0\in\Phi$, l'ensemble des puissances
$\mathscr{J}_0^n$ ($n>0$) est cofinal à $\Phi$.
\end{lemma}

\begin{proof}
En effet, si $\mathscr{J}_1$ et $\mathscr{J}_2$ appartiennent à $\Phi$, et
si on pose $\mathscr{J}=\mathscr{J}_1\cap\mathscr{J}_2$,
$\mathscr{J}$ est cohérent puisque $\mathscr{O}_X$ est cohérent
(\hyperref[I.6.1.1-fr]{6.1.1} et
\hyperref[0.5.3.4-fr]{0, 5.3.4}), et l'on a
$\mathscr{J}_x=(\mathscr{J}_1)_x\cap(\mathscr{J}_2)_x$ pour tout
$x\in X$, donc $\mathscr{J}_x=\mathscr{O}_x$ pour $x\notin X'$ et
$\mathscr{J}_x\neq\mathscr{O}_x$ pour $x\in X'$, ce qui prouve que
$\mathscr{J}\in\Phi$. D'autre part, si $X$ est noethérien, et si
$\mathscr{J}_0$ et $\mathscr{J}$ appartiennent à $\Phi$, il existe un
entier $n>0$ tel que
$\mathscr{J}_0^n(\mathscr{O}_X/\mathscr{J})=0$
(\hyperref[I.9.3.4-fr]{9.3.4}), ce qui signifie que
$\mathscr{J}_0^n\subset\mathscr{J}$.
\end{proof}

\begin{env}[10.8.3]
\label{I.10.8.3-fr}
Soit maintenant $\mathscr{F}$ un $\mathscr{O}_X$-Module cohérent~; pour
tout $\mathscr{J}\in\Phi$,
$\mathscr{F}\otimes_{\mathscr{O}_X}(\mathscr{O}_X/\mathscr{J})$ est un
$\mathscr{O}_X$-Module cohérent
(\hyperref[I.9.1.1-fr]{9.1.1}), de support contenu dans $X'$, et que nous
identifierons le plus souvent à sa restriction à $X'$. Lorsque
$\mathscr{J}$ parcourt $\Phi$, ces faisceaux forment un
\emph{système projectif} de faisceaux de groupes abéliens.
\end{env}

\begin{definition}[10.8.4]
\label{I.10.8.4-fr}
Étant donnés une partie fermée $X'$ d'un préschéma localement noethérien
$X$ et un $\mathscr{O}_X$-Module cohérent $\mathscr{F}$, on appelle
\emph{complété de $\mathscr{F}$ le long de $X'$} et on désigne par
$\mathscr{F}_{/X'}$ ou par $\widehat{\mathscr{F}}$ (lorsqu'aucune confusion
n'est possible) la restriction à $X'$ du faisceau

\oldpage[I]{195}

$\varprojlim_{\Phi}
(\mathscr{F}\otimes_{\mathscr{O}_X}
(\mathscr{O}_X/\mathscr{J}))$~; on dit que ses sections au-dessus de $X'$ sont
les \emph{sections formelles de $\mathscr{F}$ le long de $X'$}.
\end{definition}

Il est immédiat que pour tout ouvert $U\subset X$, on a
$(\mathscr{F}|U)_{/(U\cap X')}=(\mathscr{F}_{/X'})|(U\cap X')$.

Par passage à la limite projective, il est clair que
$(\mathscr{O}_X)_{/X'}$ est un faisceau d'anneaux, et que
$\mathscr{F}_{/X'}$ peut être considéré comme un
$(\mathscr{O}_X)_{/X'}$-Module. En outre, comme il existe une base de la
topologie de $X'$ formée d'ouverts quasi-compacts, on peut considérer
$(\mathscr{O}_X)_{/X'}$ (resp. $\mathscr{F}_{/X'}$) comme un
\emph{faisceau d'anneaux topologiques} (resp. de
\emph{groupes topologiques}) limite projective des faisceaux d'anneaux
(resp. groupes) \emph{pseudo-discrets}
$\mathscr{O}_X/\mathscr{J}$ (resp.
$\mathscr{F}\otimes_{\mathscr{O}_X}
(\mathscr{O}_X/\mathscr{J})=\mathscr{F}/\mathscr{J}\mathscr{F}$), et, par
passage à la limite projective, $\mathscr{F}_{/X'}$ devient alors un
\emph{$(\mathscr{O}_X)_{/X'}$-Module topologique}
(\hyperref[0.3.8.1-fr]{0, 3.8.1} et
\hyperref[0.3.8.2-fr]{3.8.2})~; rappelons que pour tout ouvert
\emph{quasi-compact} $U\subset X$,
$\Gamma(U\cap X',(\mathscr{O}_X)_{/X'})$ (resp.
$\Gamma(U\cap X',\mathscr{F}_{/X'})$) est alors limite projective des
anneaux (resp. groupes) discrets $\Gamma(U,\mathscr{O}_X/\mathscr{J})$
(resp. $\Gamma(U,\mathscr{F}/\mathscr{J}\mathscr{F})$).

Si maintenant $u:\mathscr{F}\to\mathscr{G}$ est un homomorphisme de
$\mathscr{O}_X$-Modules, on en déduit canoniquement des homomorphismes
$u_{\mathscr{J}}:
\mathscr{F}\otimes_{\mathscr{O}_X}(\mathscr{O}_X/\mathscr{J})
\to
\mathscr{G}\otimes_{\mathscr{O}_X}(\mathscr{O}_X/\mathscr{J})$
pour tout $\mathscr{J}\in\Phi$, et ces homomorphismes forment un système
projectif. Par passage à la limite projective et restriction à $X'$, ils
donnent donc un $(\mathscr{O}_X)_{/X'}$-homomorphisme continu
$\mathscr{F}_{/X'}\to\mathscr{G}_{/X'}$, noté $u_{/X'}$ ou $\widehat u$,
et appelé le \emph{complété de l'homomorphisme $u$ le long de $X'$}. Il est
clair que si $v:\mathscr{G}\to\mathscr{H}$ est un second homomorphisme de
$\mathscr{O}_X$-Modules, on a
$(v\circ u)_{/X'}=(v_{/X'})\circ(u_{/X'})$ donc
$\mathscr{F}_{/X'}$ est un \emph{foncteur additif covariant} en
$\mathscr{F}$, de la catégorie des $\mathscr{O}_X$-Modules cohérents, à
valeurs dans la catégorie des $(\mathscr{O}_X)_{/X'}$-Modules topologiques.

\begin{proposition}[10.8.5]
\label{I.10.8.5-fr}
Le support de $(\mathscr{O}_X)_{/X'}$ est $X'$~; l'espace topologiquement
annelé $(X',(\mathscr{O}_X)_{/X'})$ est un préschéma formel localement
noethérien, et si $\mathscr{J}\in\Phi$, $\mathscr{J}_{/X'}$ est un Idéal de
définition de ce préschéma formel. Si $X=\operatorname{Spec}(A)$ est un
schéma affine d'anneau noethérien, $\mathscr{J}=\widetilde{\mathfrak{J}}$
où $\mathfrak{J}$ est un idéal de $A$, et $X'=V(\mathfrak{J})$,
$(X',(\mathscr{O}_X)_{/X'})$ s'identifie canoniquement à
$\operatorname{Spf}(\widehat A)$, où $\widehat A$ est le séparé complété de
$A$ pour la topologie $\mathfrak{J}$-préadique.
\end{proposition}

\begin{proof}
On peut évidemment se borner à prouver la dernière assertion. On sait
(\hyperref[0.7.3.3-fr]{0, 7.3.3}) que le séparé complété
$\widehat{\mathfrak{J}}$ de $\mathfrak{J}$ pour la topologie
$\mathfrak{J}$-préadique s'identifie à l'idéal
$\mathfrak{J}\widehat A$ de $\widehat A$, et que $\widehat A$ est un anneau
$\widehat{\mathfrak{J}}$-adique noethérien tel que
$\widehat A/\widehat{\mathfrak{J}}^n=A/\mathfrak{J}^n$
(\hyperref[0.7.2.6-fr]{0, 7.2.6}). Cette dernière relation montre que les
idéaux premiers ouverts de $\widehat A$ sont les idéaux
$\widehat{\mathfrak p}=\mathfrak p\widehat A$, où $\mathfrak p$ est un
idéal premier de $A$ contenant $\mathfrak{J}$, et que l'on a
$\widehat{\mathfrak p}\cap A=\mathfrak p$, d'où
$\operatorname{Spf}(\widehat A)=X'$. Comme
$\mathscr{O}_X/\mathscr{J}^n=(A/\mathfrak{J}^n)^{\sim}$, la proposition
découle aussitôt des définitions.
\end{proof}

On dit que le préschéma formel ainsi défini est le
\emph{complété de $X$ le long de $X'$} et on le note $X_{/X'}$ ou
$\widehat X$ si aucune confusion n'est à craindre. Lorsque l'on prend
$X'=X$, on peut prendre $\mathscr{J}=0$, et on a donc $X_{/X}=X$.

Il est clair que si $U$ est un sous-préschéma induit sur un ouvert de $X$,
$U_{/(U\cap X')}$ s'identifie canoniquement au sous-préschéma formel induit
par $X_{/X'}$ sur l'ouvert $U\cap X'$ de $X'$.

\begin{corollary}[10.8.6]
\label{I.10.8.6-fr}
Le préschéma (usuel) $\widehat X_{\mathrm{red}}$ est l'unique
sous-préschéma réduit de $X$ ayant pour espace sous-jacent $X'$
(\hyperref[I.5.2.1-fr]{5.2.1}). Pour que $\widehat X$ soit noethérien, il
faut et il suffit que $\widehat X_{\mathrm{red}}$ le soit, et il suffit que
$X$ le soit.
\end{corollary}

\oldpage[I]{196}

La détermination de $\widehat X_{\mathrm{red}}$ étant locale
(\hyperref[I.10.5.4-fr]{10.5.4}), on peut encore supposer que $X$ est un
schéma affine d'anneau noethérien~; avec les notations de
(\hyperref[I.10.8.5-fr]{10.8.5}), l'idéal $\mathfrak{T}$ des éléments
topologiquement nilpotents de $\widehat A$ est l'image réciproque par
l'application canonique
$\widehat A\to\widehat A/\widehat{\mathfrak{J}}=A/\mathfrak{J}$ du
nilradical de $A/\mathfrak{J}$
(\hyperref[0.7.1.3-fr]{0, 7.1.3}), donc $\widehat A/\mathfrak{T}$ est
isomorphe au quotient de $A/\mathfrak{J}$ par son nilradical. La première
assertion résulte donc de (\hyperref[I.10.5.4-fr]{10.5.4}) et
(\hyperref[I.5.1.1-fr]{5.1.1}). Si $\widehat X_{\mathrm{red}}$ est
noethérien, son espace sous-jacent $X'$ l'est aussi, donc les
$X'_n=\operatorname{Spec}(\mathscr{O}_X/\mathscr{J}^n)$ sont noethériens
(\hyperref[I.6.1.2-fr]{6.1.2}) et il en est de même de $\widehat X$
(\hyperref[I.10.6.4-fr]{10.6.4})~; la réciproque est immédiate, en vertu
de (\hyperref[I.6.1.2-fr]{6.1.2}).

\begin{env}[10.8.7]
\label{I.10.8.7-fr}
Les homomorphismes canoniques
$\mathscr{O}_X\to\mathscr{O}_X/\mathscr{J}$ (pour
$\mathscr{J}\in\Phi$) forment un système projectif et donnent donc, par
passage à la limite projective, un homomorphisme de faisceau d'anneaux
$\theta:\mathscr{O}_X\to\psi_*((\mathscr{O}_X)_{/X'})
=\varprojlim_\Phi(\mathscr{O}_X/\mathscr{J})$, en désignant par $\psi$
l'injection canonique $X'\to X$ des espaces sous-jacents. Nous désignerons
par $i$ (ou $i_X$) le morphisme (dit canonique)
\[
  (\psi,\theta):X_{/X'}\to X
\]
d'espaces annelés.

Par tensorisation, pour tout $\mathscr{O}_X$-Module cohérent
$\mathscr{F}$, les homomorphismes canoniques
$\mathscr{O}_X\to\mathscr{O}_X/\mathscr{J}$ donnent des homomorphismes
$\mathscr{F}\to
\mathscr{F}\otimes_{\mathscr{O}_X}(\mathscr{O}_X/\mathscr{J})$ de
$\mathscr{O}_X$-Modules qui forment encore un système projectif, et
donnent donc, par passage à la limite projective, un homomorphisme
canonique fonctoriel
$\gamma:\mathscr{F}\to\psi_*(\mathscr{F}_{/X'})$ de
$\mathscr{O}_X$-Modules.
\end{env}

\begin{proposition}[10.8.8]
\label{I.10.8.8-fr}
\medskip\noindent
\begin{enumerate}
  \item[\rm{(i)}] Le foncteur $\mathscr{F}_{/X'}$ (en $\mathscr{F}$) est
    exact.
  \item[\rm{(ii)}] L'homomorphisme fonctoriel
    $\gamma^\sharp:i^*(\mathscr{F})\to\mathscr{F}_{/X'}$ de
    $(\mathscr{O}_X)_{/X'}$-Modules est un isomorphisme.
\end{enumerate}
\end{proposition}

\medskip\noindent
\begin{enumerate}
  \item[\rm{(i)}] Il suffit de prouver que si
    $0\to\mathscr{F}'\to\mathscr{F}\to\mathscr{F}''\to0$ est une suite
    exacte de $\mathscr{O}_X$-Modules cohérents, et $U$ un ouvert affine
    de $X$, d'anneau noethérien $A$, la suite
    \[
      0\to\Gamma(U\cap X',\mathscr{F}'_{/X'})
      \to\Gamma(U\cap X',\mathscr{F}_{/X'})
      \to\Gamma(U\cap X',\mathscr{F}''_{/X'})\to0
    \]
    est exacte. On a alors $\mathscr{F}|U=\widetilde M$,
    $\mathscr{F}'|U=\widetilde{M'}$, $\mathscr{F}''|U=\widetilde{M''}$,
    où $M$, $M'$, $M''$ sont trois $A$-modules de type fini tels que la
    suite $0\to M'\to M\to M''\to0$ soit exacte
    (\hyperref[I.1.5.1-fr]{1.5.1} et
    \hyperref[I.1.3.11-fr]{1.3.11})~; soit $\mathscr{J}\in\Phi$ et soit
    $\mathfrak{J}$ un idéal de $A$ tel que
    $\mathscr{J}|U=\widetilde{\mathfrak{J}}$. On a alors
    \[
      \Gamma(U\cap X',
      \mathscr{F}\otimes_{\mathscr{O}_X}(\mathscr{O}_X/\mathscr{J}^n))
      =M\otimes_A(A/\mathfrak{J}^n)
    \]
    (\hyperref[I.1.3.12-fr]{1.3.12})~; donc, par définition de la limite
    projective, on a
    \[
      \Gamma(U\cap X',\mathscr{F}_{/X'})
      =\varprojlim_n(M\otimes_A(A/\mathfrak{J}^n))=\widehat M
    \]
    séparé complété de $M$ pour la topologie $\mathfrak{J}$-préadique, et
    de même
    \[
      \Gamma(U\cap X',\mathscr{F}'_{/X'})=\widehat{M'},\qquad
      \Gamma(U\cap X',\mathscr{F}''_{/X'})=\widehat{M''}~;
    \]
    notre assertion résulte alors de ce que lorsque $A$ est noethérien,
    le foncteur $\widehat M$ en $M$ est exact sur la catégorie des
    $A$-modules de type fini (\hyperref[0.7.3.3-fr]{0, 7.3.3}).

\oldpage[I]{197}
  \item[\rm{(ii)}] La question étant locale, on peut supposer que l'on a
    une suite exacte
    $\mathscr{O}_X^m\to\mathscr{O}_X^n\to\mathscr{F}\to0$
    (\hyperref[0.5.3.2-fr]{0, 5.3.2})~; comme $\gamma^\sharp$ est
    fonctoriel, et que les foncteurs $i^*(\mathscr{F})$ et
    $\mathscr{F}_{/X'}$ sont exacts à droite (d'après (i) et
    (\hyperref[0.4.3.1-fr]{0, 4.3.1})), on a le diagramme commutatif
    \[
      \xymatrix{
        i^*(\mathscr{O}_X^m)\ar[r]\ar[d]_{\gamma^\sharp} &
        i^*(\mathscr{O}_X^n)\ar[r]\ar[d]_{\gamma^\sharp} &
        i^*(\mathscr{F})\ar[r]\ar[d]_{\gamma^\sharp} &
        0\\
        (\mathscr{O}_X^m)_{/X'}\ar[r] &
        (\mathscr{O}_X^n)_{/X'}\ar[r] &
        \mathscr{F}_{/X'}\ar[r] &
        0
      }
      \tag{10.8.8.1}\label{I.10.8.8.1-fr}
    \]
    dont les lignes sont exactes. En outre, les deux foncteurs
    $i^*(\mathscr{F})$ et $\mathscr{F}_{/X'}$ commutent aux sommes
    directes finies (\hyperref[0.3.2.6-fr]{0, 3.2.6} et
    \hyperref[0.4.3.2-fr]{4.3.2}) et on est donc ramené à démontrer notre
    assertion pour $\mathscr{F}=\mathscr{O}_X$. On a alors
    $i^*(\mathscr{O}_X)=(\mathscr{O}_X)_{/X'}=
    \mathscr{O}_{\widehat X}$
    (\hyperref[0.4.3.4-fr]{0, 4.3.4}), et $\gamma^\sharp$ est un
    homomorphisme de $\mathscr{O}_{\widehat X}$-Modules~; il suffit donc
    de vérifier que $\gamma^\sharp$ transforme la section unité de
    $\mathscr{O}_{\widehat X}$ au-dessus d'un ouvert de $X'$ en elle-même,
    ce qui est immédiat et montre donc que dans ce cas $\gamma^\sharp$ est
    l'identité.
\end{enumerate}

\begin{corollary}[10.8.9]
\label{I.10.8.9-fr}
Le morphisme d'espaces annelés $i:X_{/X'}\to X$ est plat.
\end{corollary}

Cela résulte en effet de (\hyperref[0.6.7.3-fr]{0, 6.7.3}) et de
(\hyperref[I.10.8.8-fr]{10.8.8, (i)}).

\begin{corollary}[10.8.10]
\label{I.10.8.10-fr}
Si $\mathscr{F}$ et $\mathscr{G}$ sont des $\mathscr{O}_X$-Modules
cohérents, il existe des isomorphismes canoniques fonctoriels (en
$\mathscr{F}$ et $\mathscr{G}$)
\[
  (\mathscr{F}_{/X'})\otimes_{(\mathscr{O}_X)_{/X'}}
  (\mathscr{G}_{/X'})
  \xrightarrow{\sim}
  (\mathscr{F}\otimes_{\mathscr{O}_X}\mathscr{G})_{/X'}
  \tag{10.8.10.1}\label{I.10.8.10.1-fr}
\]
\[
  (\mathscr{H}\!om_{\mathscr{O}_X}(\mathscr{F},\mathscr{G}))_{/X'}
  \xrightarrow{\sim}
  \mathscr{H}\!om_{(\mathscr{O}_X)_{/X'}}
  (\mathscr{F}_{/X'},\mathscr{G}_{/X'})
  \tag{10.8.10.2}\label{I.10.8.10.2-fr}
\]
\end{corollary}

Cela résulte de l'identification canonique de $i^*(\mathscr{F})$ et de
$\mathscr{F}_{/X'}$~; l'existence du premier isomorphisme est alors un
résultat valable pour tous les morphismes d'espaces annelés
(\hyperref[0.4.3.3.1-fr]{0, 4.3.3.1}) et celle du second est un résultat
valable pour tous les morphismes plats
(\hyperref[0.6.7.6-fr]{0, 6.7.6}), donc découle de
(\hyperref[I.10.8.9-fr]{10.8.9}).

\begin{proposition}[10.8.11]
\label{I.10.8.11-fr}
Pour tout $\mathscr{O}_X$-Module cohérent $\mathscr{F}$, le noyau de
l'homomorphisme canonique
$\Gamma(X,\mathscr{F})\to\Gamma(X',\mathscr{F}_{/X'})$ déduit de
$\mathscr{F}\to\mathscr{F}_{/X'}$ est formé des sections nulles dans un
voisinage de $X'$.
\end{proposition}

Il résulte de la définition de $\mathscr{F}_{/X'}$ que l'image canonique
d'une telle section est nulle. Réciproquement, si
$s\in\Gamma(X,\mathscr{F})$ a une image nulle dans
$\Gamma(X',\mathscr{F}_{/X'})$, il suffit de voir que tout $x\in X'$ admet
un voisinage dans $X$ dans lequel $s$ est nulle, et on peut donc se ramener
au cas où $X=\operatorname{Spec}(A)$ est affine, $A$ noethérien,
$X'=V(\mathfrak{J})$, où $\mathfrak{J}$ est un idéal de $A$, et
$\mathscr{F}=\widetilde M$, où $M$ est un $A$-module de type fini. Alors
$\Gamma(X',\mathscr{F}_{/X'})$ est le séparé complété $\widehat M$ de $M$
pour la topologie $\mathfrak{J}$-préadique, et l'homomorphisme
$\Gamma(X,\mathscr{F})\to\Gamma(X',\mathscr{F}_{/X'})$ est
l'homomorphisme canonique $M\to\widehat M$. On sait
(\hyperref[0.7.3.7-fr]{0, 7.3.7}) que le noyau de cet homomorphisme est
l'ensemble des $z\in M$ annulés par un élément de $1+\mathfrak{J}$. On a
donc $(1+f)s=0$ pour un $f\in\mathfrak{J}$~; pour tout $x\in X'$ on en
déduit $(1_x+f_x)s_x=0$, et comme $1_x+f_x$ est inversible dans
$\mathscr{O}_x$ ($\mathfrak{J}_x\mathscr{O}_x$ étant contenu dans l'idéal
maximal de $\mathscr{O}_x$), on a $s_x=0$, ce qui démontre la proposition.

\begin{corollary}[10.8.12]
\label{I.10.8.12-fr}
Le support de $\mathscr{F}_{/X'}$ est égal à
$\operatorname{Supp}(\mathscr{F})\cap X'$.
\end{corollary}

Il est clair que $\mathscr{F}_{/X'}$ est un
$(\mathscr{O}_X)_{/X'}$-Module de type fini
(\hyperref[I.10.8.8-fr]{10.8.8, (ii)}) et
(\hyperref[0.5.2.4-fr]{0, 5.2.4}),

\oldpage[I]{198}

donc son support est fermé (\hyperref[0.5.2.2-fr]{0, 5.2.2}) et
évidemment contenu dans $\operatorname{Supp}(\mathscr{F})\cap X'$. Pour
montrer qu'il est égal à ce dernier ensemble, on est aussitôt ramené à
prouver que la relation $\Gamma(X',\mathscr{F}_{/X'})=0$ entraîne
$\operatorname{Supp}(\mathscr{F})\cap X'=\varnothing$~; or cela résulte de
(\hyperref[I.10.8.11-fr]{10.8.11}) et de
(\hyperref[I.1.4.1-fr]{1.4.1}).

\begin{corollary}[10.8.13]
\label{I.10.8.13-fr}
Soit $u:\mathscr{F}\to\mathscr{G}$ un homomorphisme de
$\mathscr{O}_X$-Modules cohérents. Pour que
$u_{/X'}:\mathscr{F}_{/X'}\to\mathscr{G}_{/X'}$ soit nul, il faut et il
suffit que $u$ soit nul dans un voisinage de $X'$.
\end{corollary}

En effet, d'après (\hyperref[I.10.8.8-fr]{10.8.8, (ii)}), $u_{/X'}$
s'identifie à $i^*(u)$, donc si on considère $u$ comme une section au-dessus
de $X$ du faisceau
$\mathscr{H}=\mathscr{H}\!om_{\mathscr{O}_X}(\mathscr{F},\mathscr{G})$,
$u_{/X'}$ est la section de
$i^*(\mathscr{H})=\mathscr{H}_{/X'}$ au-dessus de $X'$ qui lui correspond
canoniquement ((\hyperref[I.10.8.10.2-fr]{10.8.10.2}) et
(\hyperref[0.4.4.6-fr]{0, 4.4.6})). Il suffit donc d'appliquer
(\hyperref[I.10.8.11-fr]{10.8.11}) au $\mathscr{O}_X$-Module cohérent
$\mathscr{H}$.

\begin{corollary}[10.8.14]
\label{I.10.8.14-fr}
Soit $u:\mathscr{F}\to\mathscr{G}$ un homomorphisme de
$\mathscr{O}_X$-Modules cohérents. Pour que $u_{/X'}$ soit un monomorphisme
(resp. un épimorphisme), il faut et il suffit que $u$ soit un monomorphisme
(resp. un épimorphisme) dans un voisinage de $X'$.
\end{corollary}

Soient $\mathscr{P}$ et $\mathscr{N}$ le conoyau et le noyau de $u$, de
sorte qu'on a la suite exacte
\[
  0\to\mathscr{N}\xrightarrow{v}\mathscr{F}\xrightarrow{u}
  \mathscr{G}\xrightarrow{w}\mathscr{P}\to0,
\]
d'où (\hyperref[I.10.8.8-fr]{10.8.8, (i)}) la suite exacte
\[
  0\to\mathscr{N}_{/X'}\xrightarrow{v_{/X'}}
  \mathscr{F}_{/X'}\xrightarrow{u_{/X'}}
  \mathscr{G}_{/X'}\xrightarrow{w_{/X'}}
  \mathscr{P}_{/X'}\to0.
\]
Si $u_{/X'}$ est un monomorphisme (resp. un épimorphisme), on a
$v_{/X'}=0$ (resp. $w_{/X'}=0$), donc il y a un voisinage de $X'$ dans
lequel $v=0$ (resp. $w=0$) en vertu de
(\hyperref[I.10.8.13-fr]{10.8.13}).

\subsection{Prolongement d'un morphisme aux complétés.}
\label{subsection:I.10.9-fr}

\begin{env}[10.9.1]
\label{I.10.9.1-fr}
Soient $X$, $Y$ deux préschémas (usuels) localement noethériens,
$f:X\to Y$ un morphisme, $X'$ (resp. $Y'$) une partie fermée de l'espace
sous-jacent $X$ (resp. $Y$), telles que $f(X')\subset Y'$. Soit
$\mathscr{J}$ (resp. $\mathscr{K}$) un faisceau d'idéaux de
$\mathscr{O}_X$ (resp. $\mathscr{O}_Y$) tel que le support de
$\mathscr{O}_X/\mathscr{J}$ (resp. $\mathscr{O}_Y/\mathscr{K}$) soit
$X'$ (resp. $Y'$) et que
$f^*(\mathscr{K})\mathscr{O}_X\subset\mathscr{J}$~; on notera qu'il
existe toujours de tels faisceaux d'idéaux, car on peut par exemple prendre
pour $\mathscr{J}$ le plus grand faisceau d'idéaux de $\mathscr{O}_X$
définissant un sous-préschéma de $X$ ayant $X'$ pour espace sous-jacent
(\hyperref[I.5.2.1-fr]{5.2.1}), et l'hypothèse $f(X')\subset Y'$ entraîne
alors $f^*(\mathscr{K})\mathscr{O}_X\subset\mathscr{J}$
(\hyperref[I.5.2.4-fr]{5.2.4}). On a donc pour tout entier $n>0$,
$f^*(\mathscr{K}^n)\mathscr{O}_X\subset\mathscr{J}^n$
(\hyperref[0.4.3.5-fr]{0, 4.3.5})~; par suite
(\hyperref[I.4.4.6-fr]{4.4.6}), si on pose
\[
  X'_n=(X',\mathscr{O}_X/\mathscr{J}^{n+1}),\qquad
  Y'_n=(Y',\mathscr{O}_Y/\mathscr{K}^{n+1}),
\]
on déduit de $f$ un morphisme $f_n:X'_n\to Y'_n$, et il est immédiat que
les $f_n$ forment un système inductif. Nous désignerons sa limite inductive
(\hyperref[I.10.6.8-fr]{10.6.8}) par
$\widehat f:X_{/X'}\to Y_{/Y'}$, et nous dirons (par abus de langage) que
$\widehat f$ est le \emph{prolongement de $f$ aux complétés de $X$ et $Y$
le long de $X'$ et $Y'$}. Il est immédiat de vérifier que ce morphisme ne
dépend pas du choix des faisceaux d'idéaux $\mathscr{J}$, $\mathscr{K}$
vérifiant les conditions ci-dessus. Il suffit de le voir en effet lorsque
$X$ et $Y$ sont des schémas affines noethériens d'anneaux $A$, $B$~; alors
$\mathscr{J}=\widetilde{\mathfrak{J}}$,
$\mathscr{K}=\widetilde{\mathfrak{K}}$, où $\mathfrak{J}$ (resp.
$\mathfrak{K}$) est un idéal de $A$ (resp. $B$), $f$ correspond à un
homomorphisme d'anneaux $\varphi:B\to A$ tel que
$\varphi(\mathfrak{K})\subset\mathfrak{J}$
(\hyperref[I.4.4.6-fr]{4.4.6} et
\hyperref[I.1.7.4-fr]{1.7.4})~; \emph{$\widehat f$ est alors le morphisme
qui correspond (\hyperref[I.10.2.2-fr]{10.2.2}) à l'homomorphisme continu
$\widehat\varphi:\widehat B\to\widehat A$}, où $\widehat A$ (resp.
$\widehat B$) est le séparé complété de $A$ (resp. $B$) pour la topologie
$\mathfrak{J}$-préadique (resp. $\mathfrak{K}$-préadique)
(\hyperref[I.10.6.8-fr]{10.6.8})~; et on sait que si on remplace
$\mathscr{J}$ par un autre

\oldpage[I]{199}

faisceau d'idéaux
$\mathscr{J}'=\widetilde{\mathfrak{J}'}$ tel que le support de
$\mathscr{O}_X/\mathscr{J}'$ soit encore $X'$, les topologies
$\mathfrak{J}$-préadique et $\mathfrak{J}'$-préadique sur $A$ sont les
mêmes (\hyperref[I.10.8.2-fr]{10.8.2}).

On notera que, d'après cette définition, l'application continue
$X'\to Y'$ des espaces sous-jacents à $X_{/X'}$ et $Y_{/Y'}$, qui
correspond à $\widehat f$, n'est autre que la restriction à $X'$ de $f$.
\end{env}

\begin{env}[10.9.2]
\label{I.10.9.2-fr}
Il résulte aussitôt de la définition précédente que le diagramme de
morphismes d'espaces annelés
\[
  \label{I.10.9.2.1-fr}
  \xymatrix{
    \widehat X\ar[r]^{\widehat f}\ar[d]_{i_X} &
    \widehat Y\ar[d]^{i_Y}\\
    X\ar[r]_f &
    Y
  }
\]
est commutatif, les flèches verticales étant les morphismes canoniques
(\hyperref[I.10.8.7-fr]{10.8.7}).
\end{env}

\begin{env}[10.9.3]
\label{I.10.9.3-fr}
Soient $Z$ un troisième préschéma, $g:Y\to Z$ un morphisme, $Z'$ une
partie fermée de $Z$ telle que $g(Y')\subset Z'$. Si $\widehat g$ désigne
le complété le long de $Y'$ et $Z'$ du morphisme $g$, il résulte aussitôt
de (\hyperref[I.10.9.1-fr]{10.9.1}) que l'on a
$\widehat{(g\circ f)}=\widehat g\circ\widehat f$.
\end{env}

\begin{proposition}[10.9.4]
\label{I.10.9.4-fr}
Soient $X$, $Y$ deux $S$-préschémas localement noethériens, $Y$ étant de
type fini sur $S$. Soient $f$, $g$, deux $S$-morphismes de $X$ dans $Y$
tels que $f(X')\subset Y'$, $g(X')\subset Y'$. Pour que
$\widehat f=\widehat g$, il faut et il suffit que $f$ et $g$ coïncident
dans un voisinage de $X'$.
\end{proposition}

La condition est évidemment suffisante (sans hypothèse de finitude sur
$Y$). Pour voir qu'elle est nécessaire, remarquons d'abord que l'hypothèse
$\widehat f=\widehat g$ implique $f(x)=g(x)$ pour tout $x\in X'$. D'autre
part, la question étant locale, on peut supposer que $X$ et $Y$ sont des
voisinages ouverts affines respectifs de $x$ et de $y=f(x)=g(x)$,
d'anneaux noethériens, que $S$ est affine et que
$\Gamma(Y,\mathscr{O}_Y)$ est une $\Gamma(S,\mathscr{O}_S)$-algèbre de
type fini (\hyperref[I.6.3.3-fr]{6.3.3}). Alors $f$ et $g$ correspondent à
deux $\Gamma(S,\mathscr{O}_S)$-homomorphismes $\rho$, $\sigma$ de
$\Gamma(Y,\mathscr{O}_Y)$ dans $\Gamma(X,\mathscr{O}_X)$
(\hyperref[I.1.7.3-fr]{1.7.3}), et par hypothèse, les prolongements par
continuité de ces homomorphismes au séparé complété de
$\Gamma(Y,\mathscr{O}_Y)$ sont les mêmes. On conclut de
(\hyperref[I.10.8.11-fr]{10.8.11}) que pour toute section
$s\in\Gamma(Y,\mathscr{O}_Y)$, les sections $\rho(s)$ et $\sigma(s)$
coïncident dans un voisinage de $X'$ (dépendant de $s$)~; comme
$\Gamma(Y,\mathscr{O}_Y)$ est une algèbre de type fini sur
$\Gamma(S,\mathscr{O}_S)$, on en déduit aussitôt qu'il existe un
voisinage $V$ de $X'$ tel que $\rho(s)$ et $\sigma(s)$ coïncident dans
$V$ pour \emph{toute} section $s\in\Gamma(Y,\mathscr{O}_Y)$. Si
$h\in\Gamma(Y,\mathscr{O}_X)$ est tel que $D(h)$ soit un voisinage de $x$
contenu dans $V$, on conclut de ce qui précède et de
(\hyperref[I.1.4.1-fr]{1.4.1, d)}) que $f$ et $g$ coïncident dans $D(h)$.

\begin{proposition}[10.9.5]
\label{I.10.9.5-fr}
Sous les hypothèses de (\hyperref[I.10.9.1-fr]{10.9.1}), pour tout
$\mathscr{O}_Y$-Module cohérent $\mathscr{G}$, il existe un isomorphisme
canonique fonctoriel de $(\mathscr{O}_X)_{/X'}$-Modules
\[
  (f^*(\mathscr{G}))_{/X'}\xrightarrow{\sim}
  \widehat f^*(\mathscr{G}_{/Y'})
\]
\end{proposition}

Si on identifie canoniquement $(f^*(\mathscr{G}))_{/X'}$ à
$i_X^*(f^*(\mathscr{G}))$ et $\widehat f^*(\mathscr{G}_{/Y'})$ à
$\widehat f^*(i_Y^*(\mathscr{G}))$
(\hyperref[I.10.8.8-fr]{10.8.8}), la proposition résulte aussitôt de la
commutativité du diagramme de
(\hyperref[I.10.9.2-fr]{10.9.2}).

\begin{env}[10.9.6]
\label{I.10.9.6-fr}
Soient maintenant $\mathscr{F}$ un $\mathscr{O}_X$-Module cohérent,
$\mathscr{G}$ un $\mathscr{O}_Y$-Module cohérent. Si
$u:\mathscr{G}\to\mathscr{F}$ est un $f$-morphisme de $\mathscr{G}$ dans
$\mathscr{F}$, il lui correspond un $\mathscr{O}_X$-homomorphisme
$u^\sharp:f^*(\mathscr{G})\to\mathscr{F}$, donc par complétion un
$(\mathscr{O}_X)_{/X'}$-homomorphisme continu
$(u^\sharp)_{/X'}:(f^*(\mathscr{G}))_{/X'}\to\mathscr{F}_{/X'}$, et en
vertu de (\hyperref[I.10.9.5-fr]{10.9.5}) il existe un
$\widehat f$-morphisme $v:\mathscr{G}_{/Y'}\to\mathscr{F}_{/X'}$,
\oldpage[I]{200}

et un seul, tel que $v^\sharp=(u^\sharp)_{/X'}$. Si on considère les
triplets $(\mathscr{F},X,X')$ ($\mathscr{F}$ étant un
$\mathscr{O}_X$-Module cohérent et $X'$ une partie fermée de $X$) comme
une \emph{catégorie}, les morphismes
$(\mathscr{F},X,X')\to(\mathscr{G},Y,Y')$ consistant en un morphisme de
préschémas $f:X\to Y$ tel que $f(X')\subset Y'$ et en un
$f$-morphisme $u:\mathscr{G}\to\mathscr{F}$, on peut donc dire que
$(X_{/X'},\mathscr{F}_{/X'})$ est un \emph{foncteur} en
$(\mathscr{F},X,X')$, prenant ses valeurs dans la catégorie des couples
$(\mathfrak{Z},\mathscr{H})$ formés d'un préschéma formel localement
noethérien $\mathfrak{Z}$ et d'un
$\mathscr{O}_{\mathfrak{Z}}$-Module $\mathscr{H}$, les morphismes de
cette dernière catégorie consistant en les couples formés d'un morphisme
$g$ de préschémas formels et d'un $g$-morphisme.
\end{env}

\begin{proposition}[10.9.7]
\label{I.10.9.7-fr}
Soient $S$, $X$, $Y$ trois préschémas localement noethériens,
$g:X\to S$, $h:Y\to S$ deux morphismes, $S'$ une partie fermée de $S$,
$X'$ (resp. $Y'$) une partie fermée de $X$ (resp. $Y$) telle que
$g(X')\subset S'$ (resp. $h(Y')\subset S'$)~; soit
$Z=X\times_S Y$~; supposons $Z$ localement noethérien, et soit
$Z'=p^{-1}(X')\cap q^{-1}(Y')$, où $p$ et $q$ sont les projections de
$X\times_S Y$. Dans ces conditions, le complété $Z_{/Z'}$ s'identifie au
produit des $S_{/S'}$-préschémas formels
$(X_{/X'})\times_{S_{/S'}}(Y_{/Y'})$, les morphismes structuraux
s'identifiant à $\widehat g$ et $\widehat h$, et les projections à
$\widehat p$ et $\widehat q$.
\end{proposition}

Il est immédiat que la question est locale pour $S$, $X$ et $Y$, et on
est donc ramené au cas où $S=\operatorname{Spec}(A)$,
$X=\operatorname{Spec}(B)$, $Y=\operatorname{Spec}(C)$,
$S'=V(\mathfrak{J})$, $X'=V(\mathfrak{K})$, $Y'=V(\mathfrak{L})$, où
$\mathfrak{J}$, $\mathfrak{K}$, $\mathfrak{L}$ sont trois idéaux tels
que $\varphi(\mathfrak{J})\subset\mathfrak{K}$ et
$\psi(\mathfrak{J})\subset\mathfrak{L}$, en désignant par $\varphi$ et
$\psi$ les homomorphismes $A\to B$ et $A\to C$ qui correspondent à $g$
et $h$. Alors on sait que $Z=\operatorname{Spec}(B\otimes_A C)$ et que
$Z'=V(\mathfrak{M})$, où $\mathfrak{M}$ est l'idéal
$\operatorname{Im}(\mathfrak{K}\otimes_A C)
+\operatorname{Im}(B\otimes_A\mathfrak{L})$. La conclusion résulte
(\hyperref[I.10.7.2-fr]{10.7.2}) de ce que le produit tensoriel complété
$(\widehat B\otimes_{\widehat A}\widehat C)^\wedge$ (où $\widehat A$,
$\widehat B$, $\widehat C$ sont respectivement les séparés complétés de
$A$, $B$, $C$ pour les topologies $\mathfrak{J}$-, $\mathfrak{K}$- et
$\mathfrak{L}$-préadiques) est le séparé complété du produit tensoriel
$B\otimes_A C$ pour la topologie $\mathfrak{M}$-préadique
(\hyperref[0.7.7.2-fr]{0, 7.7.2}).

On notera en outre que si $T$ est un $S$-préschéma localement
noethérien, $u:T\to X$, $v:T\to Y$ deux $S$-morphismes, $T'$ une partie
fermée de $T$ telle que $u(T')\subset X'$, $v(T')\subset Y'$, alors le
prolongement aux complétés $((u,v)_S)^\wedge$ s'identifie à
$(\widehat u,\widehat v)_{S_{/S'}}$.

\begin{corollary}[10.9.8]
\label{I.10.9.8-fr}
Soient $X$, $Y$ deux $S$-préschémas localement noethériens tels que
$X\times_S Y$ soit localement noethérien~; soient $S'$ une partie fermée
de $S$, $X'$ (resp. $Y'$) une partie fermée de $X$ (resp. $Y$) dont
l'image dans $S$ est contenue dans $S'$. Pour tout $S$-morphisme
$f:X\to Y$ tel que $f(X')\subset Y'$, le morphisme graphe
$\Gamma_{\widehat f}$ s'identifie au prolongement $(\Gamma_f)^\wedge$ du
morphisme graphe de $f$.
\end{corollary}

\begin{corollary}[10.9.9]
\label{I.10.9.9-fr}
Soient $X$, $Y$ deux préschémas localement noethériens, $f:X\to Y$ un
morphisme, $Y'$ une partie fermée de $Y$, $X'=f^{-1}(Y')$. Alors le
préschéma $X_{/X'}$ s'identifie par le diagramme commutatif
\[
  \label{I.10.9.9.1-fr}
  \xymatrix{
    X\ar[d]_f &
    X_{/X'}\ar[l]\ar[d]^{\widehat f}\\
    Y &
    Y_{/Y'}\ar[l]
  }
\]
au produit $X\times_Y(Y_{/Y'})$ de préschémas formels.
\end{corollary}

Il suffit d'appliquer (\hyperref[I.10.9.7-fr]{10.9.7}) en remplaçant
$S$ et $S'$ par $Y$, $X$ et $X'$ par $X$.
\oldpage[I]{201}

\begin{remark}[10.9.10]
\label{I.10.9.10-fr}
Si $X$ est la somme $X_1\amalg X_2$
(\hyperref[subsection:I.3.1-fr]{3.1}), $X'$ la réunion
$X'_1\cup X'_2$, où $X'_i$ est une partie fermée de $X_i$ ($i=1,2$),
on voit aussitôt que l'on a
$X_{/X'}={X_1}_{/X'_1}\amalg {X_2}_{/X'_2}$.
\end{remark}

\subsection{Application aux faisceaux cohérents sur les schémas formels
affines.}
\label{subsection:I.10.10-fr}

\begin{env}[10.10.1]
\label{I.10.10.1-fr}
Dans tout ce paragraphe, $A$ désignera un \emph{anneau adique
noethérien}, $\mathfrak{J}$ un idéal de définition de $A$. Soit
$X=\operatorname{Spec}(A)$, $\mathfrak{X}=\operatorname{Spf}(A)$, qui
s'identifie à la partie fermée $V(\mathfrak{J})$ de $X$
(\hyperref[I.10.1.2-fr]{10.1.2}). En outre, la définition
(\hyperref[I.10.1.2-fr]{10.1.2}) et la définition
(\hyperref[I.10.8.4-fr]{10.8.4}) montrent que le
\emph{schéma formel affine $\mathfrak{X}$} est identique au
\emph{complété $X_{/\mathfrak{X}}$} du schéma affine $X$ le long de la
partie fermée $\mathfrak{X}$ de son espace sous-jacent. A tout
$\mathscr{O}_X$-Module cohérent $\mathscr{F}$ correspond donc un
$\mathscr{O}_{\mathfrak{X}}$-Module de type fini
$\mathscr{F}_{/\mathfrak{X}}$, qui est d'ailleurs un faisceau de modules
topologiques sur le faisceau d'anneaux topologiques
$\mathscr{O}_{\mathfrak{X}}$. Mais tout $\mathscr{O}_X$-Module cohérent
$\mathscr{F}$ est de la forme $\widetilde M$, où $M$ est un $A$-module
de type fini (\hyperref[I.1.5.1-fr]{1.5.1})~; nous poserons
$(\widetilde M)_{/\mathfrak{X}}=M^\Delta$. En outre, si $u:M\to N$ est
un $A$-homomorphisme de $A$-modules de type fini, il lui correspond un
homomorphisme $\widetilde u:\widetilde M\to\widetilde N$, et par suite
aussi un homomorphisme continu
$\widetilde u_{/\mathfrak{X}}:(\widetilde M)_{/\mathfrak{X}}\to
(\widetilde N)_{/\mathfrak{X}}$, que nous noterons $u^\Delta$. Il est
immédiat que
$(v\circ u)^\Delta=v^\Delta\circ u^\Delta$~; on a ainsi défini un
\emph{foncteur additif covariant $M^\Delta$} de la catégorie des
$A$-modules de type fini dans celle des
$\mathscr{O}_{\mathfrak{X}}$-Modules de type fini. Lorsque $A$ est un
anneau \emph{discret}, on a $M^\Delta=\widetilde M$.
\end{env}

\begin{proposition}[10.10.2]
\label{I.10.10.2-fr}
\medskip\noindent
\begin{enumerate}
  \item[\rm{(i)}] $M^\Delta$ est un foncteur exact en $M$, et il existe
    un isomorphisme canonique fonctoriel de $A$-modules
    $\Gamma(\mathfrak{X},M^\Delta)\xrightarrow{\sim}M$.
  \item[\rm{(ii)}] Si $M$ et $N$ sont deux $A$-modules de type fini, il
    existe des isomorphismes canoniques fonctoriels
    \[
      (M\otimes_A N)^\Delta
      \xrightarrow{\sim}
      M^\Delta\otimes_{\mathscr{O}_{\mathfrak{X}}}N^\Delta
      \tag{10.10.2.1}\label{I.10.10.2.1-fr}
    \]
    \[
      (\operatorname{Hom}_A(M,N))^\Delta
      \xrightarrow{\sim}
      \mathscr{H}\!om_{\mathscr{O}_{\mathfrak{X}}}
      (M^\Delta,N^\Delta)
      \tag{10.10.2.2}\label{I.10.10.2.2-fr}
    \]
  \item[\rm{(iii)}] L'application $u\to u^\Delta$ est un isomorphisme
    fonctoriel
    \[
      \operatorname{Hom}_A(M,N)
      \xrightarrow{\sim}
      \operatorname{Hom}_{\mathscr{O}_{\mathfrak{X}}}
      (M^\Delta,N^\Delta)
      \tag{10.10.2.3}\label{I.10.10.2.3-fr}
    \]
\end{enumerate}
\end{proposition}

L'exactitude de $M^\Delta$ résulte de l'exactitude des foncteurs
$\widetilde M$ (\hyperref[I.1.3.5-fr]{1.3.5}) et
$\mathscr{F}_{/\mathfrak{X}}$ (\hyperref[I.10.8.8-fr]{10.8.8}). Par
définition, $\Gamma(\mathfrak{X},M^\Delta)$ est le séparé complété du
$A$-module $\Gamma(X,\widetilde M)=M$ pour la topologie
$\mathfrak{J}$-préadique~; mais comme $A$ est complet et $M$ de type
fini, on sait (\hyperref[0.7.3.6-fr]{0, 7.3.6}) que $M$ est séparé et
complet, ce qui achève de prouver (i). L'isomorphisme
(\hyperref[I.10.10.2.1-fr]{10.10.2.1}) (resp.
\hyperref[I.10.10.2.2-fr]{10.10.2.2}) provient de la composition des
isomorphismes (\hyperref[I.1.3.12-fr]{1.3.12, (i)}) et
(\hyperref[I.10.8.10.1-fr]{10.8.10.1}) (resp.
(\hyperref[I.1.3.12-fr]{1.3.12, (ii)}) et
(\hyperref[I.10.8.10.2-fr]{10.8.10.2})). Enfin, comme
$\operatorname{Hom}_A(M,N)$ est un $A$-module de type fini, on peut lui
appliquer (i), qui identifie
$\Gamma(\mathfrak{X},(\operatorname{Hom}_A(M,N))^\Delta)$ à
$\operatorname{Hom}_A(M,N)$, et utiliser
(\hyperref[I.10.10.2.2-fr]{10.10.2.2}), qui prouve que
l'homomorphisme (\hyperref[I.10.10.2.3-fr]{10.10.2.3}) est un
isomorphisme.

On déduit de (\hyperref[I.10.10.2-fr]{10.10.2}) toute une série de
conséquences analogues à celles déduites de
(\hyperref[I.1.3.7-fr]{1.3.7}) et
(\hyperref[I.1.3.12-fr]{1.3.12}), que nous laissons au lecteur le soin
de formuler.
\oldpage[I]{202}

Notons que la propriété d'exactitude de $M^\Delta$, appliqué à la suite
exacte
$0\to\mathfrak{J}\to A\to A/\mathfrak{J}\to0$ montre que le faisceau
d'idéaux de $\mathscr{O}_{\mathfrak{X}}$ désigné ici par
$\mathfrak{J}^\Delta$ coïncide avec celui qui avait été noté de la même
façon en (\hyperref[I.10.3.1-fr]{10.3.1}), en vertu de
(\hyperref[I.10.3.2-fr]{10.3.2}).

\begin{proposition}[10.10.3]
\label{I.10.10.3-fr}
Sous les hypothèses de (\hyperref[I.10.10.1-fr]{10.10.1}),
$\mathscr{O}_{\mathfrak{X}}$ est un faisceau cohérent d'anneaux.
\end{proposition}

Si $f\in A$, on sait que $A_{\{f\}}$ est un anneau adique noethérien
(\hyperref[0.7.6.11-fr]{0, 7.6.11}) et comme la question est locale, on
est ramené (\hyperref[I.10.1.4-fr]{10.1.4}) à prouver que le noyau d'un
homomorphisme
$v:\mathscr{O}_{\mathfrak{X}}^n\to\mathscr{O}_{\mathfrak{X}}$ est un
$\mathscr{O}_{\mathfrak{X}}$-Module de type fini. On a alors
$v=u^\Delta$, où $u$ est un $A$-homomorphisme $A^n\to A$
(\hyperref[I.10.10.2-fr]{10.10.2})~; comme $A$ est noethérien, le noyau
de $u$ est de type fini, autrement dit on a un homomorphisme
$A^m\xrightarrow{w}A^n$ tel que la suite
$A^m\xrightarrow{w}A^n\xrightarrow{u}A$ soit exacte. On en conclut
(\hyperref[I.10.10.2-fr]{10.10.2}) que la suite
$\mathscr{O}_{\mathfrak{X}}^m\xrightarrow{w^\Delta}
\mathscr{O}_{\mathfrak{X}}^n\xrightarrow{v}
\mathscr{O}_{\mathfrak{X}}$ est exacte, ce qui prouve que le noyau de
$v$ est de type fini.

\begin{env}[10.10.4]
\label{I.10.10.4-fr}
Avec les notations précédentes, posons
$A_n=A/\mathfrak{J}^{n+1}$, et soit $X_n$ le schéma affine
$\operatorname{Spec}(A_n)=(\mathfrak{X},
\mathscr{O}_{\mathfrak{X}}/\mathscr{J}^{n+1})$,
$\mathscr{J}=\mathfrak{J}^\Delta$ étant le faisceau d'idéaux de
définition de $\mathscr{O}_{\mathfrak{X}}$ correspondant à l'idéal
$\mathfrak{J}$. Soit $u_{mn}$ le morphisme de préschémas
$X_m\to X_n$ correspondant à l'homomorphisme canonique
$A_n\to A_m$ pour $m\leq n$~; le schéma formel $\mathfrak{X}$ est
limite inductive des $X_n$ pour les $u_{mn}$
(\hyperref[I.10.6.3-fr]{10.6.3}).
\end{env}

\begin{proposition}[10.10.5]
\label{I.10.10.5-fr}
Sous les hypothèses de (\hyperref[I.10.10.1-fr]{10.10.1}), soit
$\mathscr{F}$ un $\mathscr{O}_{\mathfrak{X}}$-Module. Les conditions
suivantes sont équivalentes~:
\begin{enumerate}
  \item[a)] $\mathscr{F}$ est un $\mathscr{O}_{\mathfrak{X}}$-Module
    cohérent.
  \item[b)] $\mathscr{F}$ est isomorphe à la limite projective
    (\hyperref[I.10.6.6-fr]{10.6.6}) d'une suite $(\mathscr{F}_n)$ de
    $\mathscr{O}_{X_n}$-Modules cohérents tels que
    $u_{mn}^*(\mathscr{F}_n)=\mathscr{F}_m$.
  \item[c)] Il existe un $A$-module de type fini $M$ (déterminé à un
    isomorphisme canonique près par
    (\hyperref[I.10.10.2-fr]{10.10.2, (i)})) tel que $\mathscr{F}$ soit
    isomorphe à $M^\Delta$.
\end{enumerate}
\end{proposition}

Montrons d'abord que b) implique c). On a
$\mathscr{F}_n=\widetilde{M_n}$, où $M_n$ est un $A_n$-module de type
fini, et l'hypothèse entraîne que
$M_m=M_n\otimes_{A_n}A_m$ pour $m\leq n$
(\hyperref[I.1.6.5-fr]{1.6.5})~; les $M_n$ forment donc un système
projectif pour les di-homomorphismes canoniques $M_n\to M_m$
($m\leq n$), et il résulte aussitôt de la définition des $A_n$ que ce
système projectif vérifie les conditions de
(\hyperref[0.7.2.9-fr]{0, 7.2.9})~; sa limite projective $M$ est par
suite un $A$-module de type fini tel que
$M_n=M\otimes_A A_n$ pour tout $n$. On en déduit que $\mathscr{F}_n$
est induit sur $X_n$ par
$\widetilde M\otimes_{\mathscr{O}_X}
(\mathscr{O}_X/\widetilde{\mathfrak{J}}^{n+1})$, donc
$\mathscr{F}=M^\Delta$ par définition
(\hyperref[I.10.8.4-fr]{10.8.4}).

Inversement, c) entraîne b)~; en effet, si $u_n$ est le morphisme
d'immersion $X_n\to X$,
$u_n^*(\widetilde M)=(M\otimes_A A_n)^{\sim}$ est induit sur $X_n$ par
$\widetilde M\otimes_{\mathscr{O}_X}
(\mathscr{O}_X/\widetilde{\mathfrak{J}}^{n+1})$, et
$M^\Delta=\varprojlim u_n^*(\widetilde M)$ par définition
(\hyperref[I.10.8.4-fr]{10.8.4})~; comme
$u_m=u_n\circ u_{mn}$ pour $m\leq n$, les
$\mathscr{F}_n=u_n^*(\widetilde M)$ vérifient les conditions de b),
d'où notre assertion.

Montrons maintenant que c) implique a)~: en effet, on a par définition
$\mathscr{O}_{\mathfrak{X}}=A^\Delta$~; $M$ étant conoyau d'un
homomorphisme $A^m\to A^n$, il résulte de
(\hyperref[I.10.10.2-fr]{10.10.2}) que $M^\Delta$ est conoyau d'un
homomorphisme
$\mathscr{O}_{\mathfrak{X}}^m\to\mathscr{O}_{\mathfrak{X}}^n$, et
comme le faisceau d'anneaux $\mathscr{O}_{\mathfrak{X}}$ est cohérent
(\hyperref[I.10.10.3-fr]{10.10.3}), il en est de même de $M^\Delta$
(\hyperref[0.5.3.4-fr]{0, 5.3.4}).
\oldpage[I]{203}

Enfin, a) entraîne b). Considéré comme
$\mathscr{O}_{\mathfrak{X}}$-Module, on a
$\mathscr{O}_{X_n}=\mathscr{O}_{\mathfrak{X}}/
\mathscr{J}^{n+1}=A_n^\Delta$~; $\mathscr{F}_n=\mathscr{F}
\otimes_{\mathscr{O}_{\mathfrak{X}}}\mathscr{O}_{X_n}$ est un
$\mathscr{O}_{\mathfrak{X}}$-Module cohérent
(\hyperref[0.5.3.5-fr]{0, 5.3.5}), et comme c'est aussi un
$\mathscr{O}_{X_n}$-Module et que $\mathscr{J}^{n+1}$ est cohérent,
on en conclut que $\mathscr{F}_n$ est un
$\mathscr{O}_{X_n}$-Module cohérent
(\hyperref[0.5.3.10-fr]{0, 5.3.10}), et il est immédiat que
$u_{mn}^*(\mathscr{F}_n)=\mathscr{F}_m$ pour $m\leq n$ (en se
souvenant que l'application continue $X_m\to X_n$ des espaces
sous-jacents est l'identité de $\mathfrak{X}$). Le faisceau
$\mathscr{G}=\varprojlim\mathscr{F}_n$ est donc un
$\mathscr{O}_{\mathfrak{X}}$-Module cohérent, puisqu'on a vu que b)
entraîne a). Les homomorphismes canoniques
$\mathscr{F}\to\mathscr{F}_n$ forment un système projectif, qui par
passage à la limite donne un homomorphisme canonique
$w:\mathscr{F}\to\mathscr{G}$, et tout revient à démontrer que $w$ est
bijectif. La question étant maintenant locale, on peut se borner au cas
où $\mathscr{F}$ est conoyau d'un homomorphisme
$\mathscr{O}_{\mathfrak{X}}^p\to\mathscr{O}_{\mathfrak{X}}^q$~; cet
homomorphisme étant de la forme $v^\Delta$, où $v$ est un homomorphisme
$A^m\to A^n$ (\hyperref[I.10.10.2-fr]{10.10.2}), $\mathscr{F}$ est
isomorphe à $M^\Delta$, où $M=\operatorname{Coker}v$
(\hyperref[I.10.10.2-fr]{10.10.2}). On a alors, en vertu de
(\hyperref[I.10.10.2-fr]{10.10.2}),
$\mathscr{F}_n=M^\Delta\otimes_{\mathscr{O}_{\mathfrak{X}}}A_n^\Delta
=(M\otimes_A A_n)^\Delta$, et comme la topologie $\mathfrak{J}$-adique
sur $M\otimes_A A_n$ est discrète, on a
$(M\otimes_A A_n)^\Delta=(M\otimes_A A_n)^{\sim}$ (en tant que
$\mathscr{O}_{X_n}$-Module)~; on a vu plus haut que
$M^\Delta=\varprojlim\mathscr{F}_n$ et $w$ est donc bien dans ce cas
l'identité.
\par\hfill C.Q.F.D.\par

\begin{corollary}[10.10.6]
\label{I.10.10.6-fr}
Si $\mathscr{F}$ vérifie la condition b) de
(\hyperref[I.10.10.5-fr]{10.10.5}), le système projectif
$(\mathscr{F}_n)$ est isomorphe au système des
$\mathscr{F}\otimes_{\mathscr{O}_{\mathfrak{X}}}\mathscr{O}_{X_n}$.
\end{corollary}

\begin{env}[10.10.7]
\label{I.10.10.7-fr}
Soient maintenant $A$, $B$ deux anneaux adiques noethériens,
$\varphi:B\to A$ un homomorphisme continu~; on désignera par
$\mathfrak{J}$ (resp. $\mathfrak{K}$) un idéal de définition de $A$
(resp. $B$), tels que $\varphi(\mathfrak{K})\subset\mathfrak{J}$, et
on posera $X=\operatorname{Spec}(A)$, $Y=\operatorname{Spec}(B)$,
$\mathfrak{X}=\operatorname{Spf}(A)$,
$\mathfrak{Y}=\operatorname{Spf}(B)$. Soient $f:X\to Y$ le morphisme
de préschémas correspondant à $\varphi$
(\hyperref[I.1.6.1-fr]{1.6.1}), $\widehat f:\mathfrak{X}\to
\mathfrak{Y}$ son prolongement aux complétés
(\hyperref[I.10.9.1-fr]{10.9.1}), qui est aussi le morphisme de
préschémas formels correspondant à $\varphi$
(\hyperref[I.10.2.2-fr]{10.2.2}).
\end{env}

\begin{proposition}[10.10.8]
\label{I.10.10.8-fr}
Pour tout $B$-module $N$ de type fini, il existe un isomorphisme
canonique fonctoriel de $\mathscr{O}_{\mathfrak{X}}$-Modules
\[
  \widehat f^*(N^\Delta)
  \xrightarrow{\sim}
  (N\otimes_B A)^\Delta.
\]
\end{proposition}

En effet, en désignant par $i_X:\mathfrak{X}\to X$ et
$i_Y:\mathfrak{Y}\to Y$ les morphismes canoniques, on a
(\hyperref[I.10.8.8-fr]{10.8.8}), à des isomorphismes canoniques
fonctoriels près, $N^\Delta=i_Y^*(\widetilde N)$ et
\[
  (N\otimes_B A)^\Delta
  =i_X^*((N\otimes_B A)^{\sim})
  =i_X^*(f^*(\widetilde N))
\]
(\hyperref[I.1.6.5-fr]{1.6.5})~; la proposition résulte donc de la
commutativité du diagramme (\hyperref[I.10.9.2-fr]{10.9.2}).

\begin{corollary}[10.10.9]
\label{I.10.10.9-fr}
Pour tout idéal $\mathfrak b$ de $B$, on a
$\widehat f^*(\mathfrak b^\Delta)\mathscr{O}_{\mathfrak{X}}
=(\mathfrak bA)^\Delta$.
\end{corollary}

En effet, soit $j$ l'injection canonique $\mathfrak b\to B$, à laquelle
il correspond l'injection canonique
$j^\Delta:\mathfrak b^\Delta\to\mathscr{O}_{\mathfrak{Y}}$ de
faisceaux de $\mathscr{O}_{\mathfrak{Y}}$-Modules~; par définition,
$\widehat f^*(\mathfrak b^\Delta)\mathscr{O}_{\mathfrak{X}}$ est
l'image de l'homomorphisme
$\widehat f^*(j^\Delta):\widehat f^*(\mathfrak b^\Delta)\to
\mathscr{O}_{\mathfrak{X}}=\widehat f^*(\mathscr{O}_{\mathfrak{Y}})$~;
mais cet homomorphisme s'identifie à
$(j\otimes1)^\Delta:(\mathfrak b\otimes_B A)^\Delta\to
\mathscr{O}_{\mathfrak{X}}=(B\otimes_B A)^\Delta$ d'après
(\hyperref[I.10.10.8-fr]{10.10.8}). Comme l'image de $j\otimes1$ est
l'idéal $\mathfrak bA$ de $A$, l'image de $(j\otimes1)^\Delta$ est
donc $(\mathfrak bA)^\Delta$ en vertu de
(\hyperref[I.10.10.2-fr]{10.10.2}), d'où la conclusion.
\oldpage[I]{204}

\subsection{Faisceaux cohérents sur les préschémas formels.}
\label{subsection:I.10.11-fr}

\begin{proposition}[10.11.1]
\label{I.10.11.1-fr}
Si $\mathfrak{X}$ est un préschéma formel localement noethérien, le
faisceau d'anneaux $\mathscr{O}_{\mathfrak{X}}$ est cohérent et tout
faisceau d'idéaux de définition de $\mathfrak{X}$ est cohérent.
\end{proposition}

La question étant locale, on est ramené au cas d'un schéma formel affine
noethérien, et la proposition résulte donc de
(\hyperref[I.10.10.3-fr]{10.10.3}) et
(\hyperref[I.10.10.5-fr]{10.10.5}).

\begin{env}[10.11.2]
\label{I.10.11.2-fr}
Soient $\mathfrak{X}$ un préschéma formel localement noethérien,
$\mathscr{J}$ un faisceau d'idéaux de définition de $\mathfrak{X}$,
$X_n$ le préschéma (usuel) localement noethérien
$(\mathfrak{X},\mathscr{O}_{\mathfrak{X}}/\mathscr{J}^{n+1})$, de
sorte que $\mathfrak{X}$ est limite inductive de la suite $(X_n)$ pour
les morphismes canoniques $u_{mn}:X_m\to X_n$
(\hyperref[I.10.6.3-fr]{10.6.3}). Avec ces notations~:
\end{env}

\begin{theorem}[10.11.3]
\label{I.10.11.3-fr}
Pour qu'un $\mathscr{O}_{\mathfrak{X}}$-Module $\mathscr{F}$ soit
cohérent, il faut et il suffit qu'il soit isomorphe à une limite
projective d'une suite $(\mathscr{F}_n)$, où $\mathscr{F}_n$ est un
$\mathscr{O}_{X_n}$-Module cohérent, tel que
$u_{mn}^*(\mathscr{F}_n)=\mathscr{F}_m$ pour $m\leq n$
(\hyperref[I.10.6.6-fr]{10.6.6}). Le système projectif
$(\mathscr{F}_n)$ est alors isomorphe au système des
$u_n^*(\mathscr{F})=\mathscr{F}
\otimes_{\mathscr{O}_{\mathfrak{X}}}\mathscr{O}_{X_n}$, $u_n$ étant
le morphisme canonique $X_n\to\mathfrak{X}$.
\end{theorem}

La question étant locale, on est ramené au cas où $\mathfrak{X}$ est
un schéma formel affine noethérien, et le théorème est alors conséquence
de (\hyperref[I.10.10.5-fr]{10.10.5}) et
(\hyperref[I.10.10.6-fr]{10.10.6}).

On peut donc dire que \emph{la donnée d'un
$\mathscr{O}_{\mathfrak{X}}$-Module cohérent équivaut à celle d'un
système projectif $(\mathscr{F}_n)$ de
$\mathscr{O}_{X_n}$-Modules cohérents tels que
$u_{mn}^*(\mathscr{F}_n)=\mathscr{F}_m$ pour $m\leq n$}.

\begin{corollary}[10.11.4]
\label{I.10.11.4-fr}
Si $\mathscr{F}$ et $\mathscr{G}$ sont deux
$\mathscr{O}_{\mathfrak{X}}$-Modules cohérents, on peut (avec les
notations de (\hyperref[I.10.11.3-fr]{10.11.3})) définir un
isomorphisme canonique fonctoriel
\[
  \label{I.10.11.4.1-fr}
  \operatorname{Hom}_{\mathscr{O}_{\mathfrak{X}}}
  (\mathscr{F},\mathscr{G})
  \xrightarrow{\sim}
  \varprojlim_n
  \operatorname{Hom}_{\mathscr{O}_{X_n}}
  (\mathscr{F}_n,\mathscr{G}_n).
  \tag{10.11.4.1}
\]
\end{corollary}

La limite projective du second membre doit s'entendre pour les
applications
$\theta_n\mapsto u_{mn}^*(\theta_n)$ ($m\leq n$) de
$\operatorname{Hom}_{\mathscr{O}_{X_n}}
(\mathscr{F}_n,\mathscr{G}_n)$ dans
$\operatorname{Hom}_{\mathscr{O}_{X_m}}
(\mathscr{F}_m,\mathscr{G}_m)$. L'homomorphisme
(\hyperref[I.10.11.4.1-fr]{10.11.4.1}) fait correspondre à un élément
$\theta\in\operatorname{Hom}_{\mathscr{O}_{\mathfrak{X}}}
(\mathscr{F},\mathscr{G})$ la suite $(u_n^*(\theta))$~; on voit
aussitôt qu'on en définit un homomorphisme réciproque du précédent en
faisant correspondre au système projectif
$(\theta_n)\in\varprojlim_n
\operatorname{Hom}_{\mathscr{O}_{X_n}}
(\mathscr{F}_n,\mathscr{G}_n)$
sa limite projective dans
$\operatorname{Hom}_{\mathscr{O}_{\mathfrak{X}}}
(\mathscr{F},\mathscr{G})$, compte tenu de
(\hyperref[I.10.11.3-fr]{10.11.3}).

\begin{corollary}[10.11.5]
\label{I.10.11.5-fr}
Pour qu'un homomorphisme $\theta:\mathscr{F}\to\mathscr{G}$ soit
surjectif, il faut et il suffit que l'homomorphisme correspondant
$\theta_0=u_0^*(\theta):\mathscr{F}_0\to\mathscr{G}_0$ le soit.
\end{corollary}

La question étant locale, on est ramené au cas où
$\mathfrak{X}=\operatorname{Spf}(A)$, $A$ étant adique noethérien,
$\mathscr{F}=M^\Delta$, $\mathscr{G}=N^\Delta$ et $\theta=u^\Delta$,
où $M$ et $N$ sont des $A$-modules de type fini et $u$ un
homomorphisme $M\to N$~; on a en outre alors
$\theta_0=\widetilde{u_0}$, où $u_0$ est l'homomorphisme
$u\otimes1:M\otimes_A A/\mathfrak{J}\to
N\otimes_A A/\mathfrak{J}$~;
la conclusion résulte de ce que $\theta$ et $u$ (resp. $\theta_0$ et
$u_0$) sont simultanément surjectifs
(\hyperref[I.1.3.9-fr]{1.3.9} et
\hyperref[I.10.10.2-fr]{10.10.2}) et de ce que $u$ et $u_0$ sont
simultanément surjectifs (\hyperref[0.7.1.14-fr]{0, 7.1.14}).

\begin{env}[10.11.6]
\label{I.10.11.6-fr}
Le th. (\hyperref[I.10.11.3-fr]{10.11.3}) montre qu'on peut considérer
tout $\mathscr{O}_{\mathfrak{X}}$-Module cohérent $\mathscr{F}$ comme
un $\mathscr{O}_{\mathfrak{X}}$-Module \emph{topologique}, en le
considérant comme limite projective des faisceaux de groupes
\emph{pseudo-discrets} $\mathscr{F}_n$
(\hyperref[0.3.8.1-fr]{0, 3.8.1}). Il résulte alors de
(\hyperref[I.10.11.4-fr]{10.11.4}) que tout homomorphisme
$u:\mathscr{F}\to\mathscr{G}$ de
$\mathscr{O}_{\mathfrak{X}}$-Modules cohérents est automatiquement
\emph{continu}

\oldpage[I]{205}

(\hyperref[0.3.8.2-fr]{0, 3.8.2}). En outre, si $\mathscr{H}$ est un
sous-$\mathscr{O}_{\mathfrak{X}}$-Module cohérent d'un
$\mathscr{O}_{\mathfrak{X}}$-Module cohérent $\mathscr{F}$, pour tout
ouvert $U\subset\mathfrak{X}$, $\Gamma(U,\mathscr{H})$ est un sous-groupe
\emph{fermé} du groupe topologique $\Gamma(U,\mathscr{F})$ car le foncteur
$\Gamma$ étant exact à gauche, $\Gamma(U,\mathscr{H})$ est le noyau de
l'homomorphisme
$\Gamma(U,\mathscr{F})\to\Gamma(U,\mathscr{F}/\mathscr{H})$, qui est
\emph{continu} d'après ce qui précède, puisque $\mathscr{F}/\mathscr{H}$
est cohérent (\hyperref[0.5.3.4-fr]{0, 5.3.4})~; notre assertion résulte
de ce que $\Gamma(U,\mathscr{F}/\mathscr{H})$ est un groupe topologique
séparé.
\end{env}

\begin{proposition}[10.11.7]
\label{I.10.11.7-fr}
Soient $\mathscr{F}$ et $\mathscr{G}$ deux
$\mathscr{O}_{\mathfrak{X}}$-Modules cohérents. On peut définir (avec les
notations de (\hyperref[I.10.11.3-fr]{10.11.3})) des isomorphismes
canoniques fonctoriels de $\mathscr{O}_{\mathfrak{X}}$-Modules
topologiques (\hyperref[I.10.11.6-fr]{10.11.6})
\[
  \label{I.10.11.7.1-fr}
  \mathscr{F}\otimes_{\mathscr{O}_{\mathfrak{X}}}\mathscr{G}
  \xrightarrow{\sim}
  \varprojlim_n
  (\mathscr{F}_n\otimes_{\mathscr{O}_{X_n}}\mathscr{G}_n)
  \tag{10.11.7.1}
\]
\[
  \label{I.10.11.7.2-fr}
  \mathscr{H}\!om_{\mathscr{O}_{\mathfrak{X}}}
  (\mathscr{F},\mathscr{G})
  \xrightarrow{\sim}
  \varprojlim_n
  \mathscr{H}\!om_{\mathscr{O}_{X_n}}(\mathscr{F}_n,\mathscr{G}_n)
  \tag{10.11.7.2}
\]
\end{proposition}

L'existence de l'isomorphisme
(\hyperref[I.10.11.7.1-fr]{10.11.7.1}) résulte de la formule
\[
  \mathscr{F}_n\otimes_{\mathscr{O}_{X_n}}\mathscr{G}_n
  =
  (\mathscr{F}\otimes_{\mathscr{O}_{\mathfrak{X}}}
  \mathscr{O}_{X_n})
  \otimes_{\mathscr{O}_{X_n}}
  (\mathscr{G}\otimes_{\mathscr{O}_{\mathfrak{X}}}
  \mathscr{O}_{X_n})
  =
  (\mathscr{F}\otimes_{\mathscr{O}_{\mathfrak{X}}}\mathscr{G})
  \otimes_{\mathscr{O}_{\mathfrak{X}}}\mathscr{O}_{X_n}
\]
et de (\hyperref[I.10.11.3-fr]{10.11.3}). L'isomorphisme
(\hyperref[I.10.11.7.2-fr]{10.11.7.2}) où les deux membres sont
considérés comme faisceaux de modules sans topologie, résulte de la
définition des sections de
$\mathscr{H}\!om_{\mathscr{O}_{\mathfrak{X}}}
(\mathscr{F},\mathscr{G})$ et
$\mathscr{H}\!om_{\mathscr{O}_{X_n}}(\mathscr{F}_n,\mathscr{G}_n)$ et de
l'existence de l'isomorphisme
(\hyperref[I.10.11.4.1-fr]{10.11.4.1}), appliqué au préschéma induit sur
un ouvert formel affine noethérien arbitraire de $\mathfrak{X}$. Reste à
prouver que l'isomorphisme
(\hyperref[I.10.11.7.2-fr]{10.11.7.2}) est bicontinu au-dessus d'un
ensemble quasi-compact, et on est donc ramené au cas où
$\mathfrak{X}=\operatorname{Spf}(A)$, $A$ étant adique noethérien, d'où
(\hyperref[I.10.10.5-fr]{10.10.5})
$\mathscr{F}=M^\Delta$, $\mathscr{G}=N^\Delta$, $M$, $N$ étant des
$A$-modules de type fini~; compte tenu de
(\hyperref[I.10.10.2.1-fr]{10.10.2.1}),
(\hyperref[I.10.10.2.3-fr]{10.10.2.3}) et
(\hyperref[I.1.3.12-fr]{1.3.12, (ii)}), on est ramené à montrer que
l'isomorphisme canonique
$\operatorname{Hom}_A(M,N)\xrightarrow{\sim}
\varprojlim_n\operatorname{Hom}_{A_n}(M_n,N_n)$ (avec
$M_n=M\otimes_A A_n$, $N_n=N\otimes_A A_n$) est continu, ce qui a été
prouvé dans (\hyperref[0.7.8.2-fr]{0, 7.8.2}).

\begin{env}[10.11.8]
\label{I.10.11.8-fr}
Comme $\operatorname{Hom}_{\mathscr{O}_{\mathfrak{X}}}
(\mathscr{F},\mathscr{G})$ est le groupe des sections du faisceau de
groupes topologiques
$\mathscr{H}\!om_{\mathscr{O}_{\mathfrak{X}}}
(\mathscr{F},\mathscr{G})$, il est muni d'une topologie de groupe. Si
$\mathfrak{X}$ est \emph{noethérien}, il résulte de
(\hyperref[I.10.11.7.2-fr]{10.11.7.2}) qu'un système fondamental de
voisinages de $0$ dans ce groupe s'obtient en prenant les sous-groupes
$\operatorname{Hom}_{\mathscr{O}_{\mathfrak{X}}}
(\mathscr{F},\mathscr{J}^n\mathscr{G})$ ($n$ arbitraire).
\end{env}

\begin{proposition}[10.11.9]
\label{I.10.11.9-fr}
Soient $\mathfrak{X}$ un préschéma formel noethérien, $\mathscr{F}$ et
$\mathscr{G}$ deux $\mathscr{O}_{\mathfrak{X}}$-Modules cohérents. Dans
le groupe topologique
$\operatorname{Hom}_{\mathscr{O}_{\mathfrak{X}}}
(\mathscr{F},\mathscr{G})$ les homomorphismes surjectifs (resp.
injectifs, bijectifs) forment une partie ouverte.
\end{proposition}

En vertu de (\hyperref[I.10.11.5-fr]{10.11.5}), l'ensemble des
homomorphismes surjectifs dans
$\operatorname{Hom}_{\mathscr{O}_{\mathfrak{X}}}
(\mathscr{F},\mathscr{G})$ est l'image réciproque par l'application
continue
$\operatorname{Hom}_{\mathscr{O}_{\mathfrak{X}}}
(\mathscr{F},\mathscr{G})\to
\operatorname{Hom}_{\mathscr{O}_{X_0}}(\mathscr{F}_0,\mathscr{G}_0)$
d'une partie du groupe discret
$\operatorname{Hom}_{\mathscr{O}_{X_0}}(\mathscr{F}_0,\mathscr{G}_0)$
d'où la première assertion. Pour démontrer la seconde, recouvrons
$\mathfrak{X}$ par un nombre fini d'ouverts formels affines noethériens
$U_i$. Pour que
$\theta\in\operatorname{Hom}_{\mathscr{O}_{\mathfrak{X}}}
(\mathscr{F},\mathscr{G})$ soit injectif, il faut et il suffit que toutes
ses images par les applications (continues) de restriction
$\operatorname{Hom}_{\mathscr{O}_{\mathfrak{X}}}
(\mathscr{F},\mathscr{G})\to
\operatorname{Hom}_{\mathscr{O}_{\mathfrak{X}}|U_i}
(\mathscr{F}|U_i,\mathscr{G}|U_i)$ le soient~; on est donc ramené au cas
affine, et alors cela a déjà été prouvé dans
(\hyperref[0.7.8.3-fr]{0, 7.8.3}).

\oldpage[I]{206}

\subsection{Morphismes adiques de préschémas formels.}
\label{subsection:I.10.12-fr}

\begin{env}[10.12.1]
\label{I.10.12.1-fr}
Soient $\mathfrak{X}$, $\mathfrak{S}$ deux préschémas formels
\emph{localement noethériens}~; nous dirons qu'un morphisme
$f:\mathfrak{X}\to\mathfrak{S}$ est \emph{adique} s'il existe un Idéal
de définition $\mathscr{J}$ de $\mathfrak{S}$ tel que
$\mathscr{K}=f^*(\mathscr{J})\mathscr{O}_{\mathfrak{X}}$ soit un Idéal de
définition de $\mathfrak{X}$~; on dit aussi alors que $\mathfrak{X}$ est
un \emph{$\mathfrak{S}$-préschéma adique} (pour $f$). Lorsqu'il en est
ainsi, pour \emph{tout} Idéal de définition $\mathscr{J}_1$ de
$\mathfrak{S}$,
$\mathscr{K}_1=f^*(\mathscr{J}_1)\mathscr{O}_{\mathfrak{X}}$ est un Idéal
de définition de $\mathfrak{X}$. En effet, la question étant locale, on
peut supposer $\mathfrak{X}$ et $\mathfrak{S}$ affines noethériens~; il
existe donc un entier $n$ tel que $\mathscr{J}^n\subset\mathscr{J}_1$ et
$\mathscr{J}_1^n\subset\mathscr{J}$
(\hyperref[I.10.3.6-fr]{10.3.6} et
\hyperref[0.7.1.4-fr]{0, 7.1.4}), d'où
$\mathscr{K}^n\subset\mathscr{K}_1$ et
$\mathscr{K}_1^n\subset\mathscr{K}$. La première de ces relations prouve
que $\mathscr{K}_1=\mathfrak{K}_1^\Delta$, où $\mathfrak{K}_1$ est un
idéal ouvert de
$A=\Gamma(\mathfrak{X},\mathscr{O}_{\mathfrak{X}})$, et la seconde prouve
que $\mathfrak{K}_1$ est un idéal de définition de $A$
(\hyperref[0.7.1.4-fr]{0, 7.1.4}), d'où notre assertion.
\end{env}

Il résulte aussitôt de ce qui précède que si $\mathfrak{X}$ et
$\mathfrak{Y}$ sont deux $\mathfrak{S}$-préschémas adiques, \emph{tout
$\mathfrak{S}$-morphisme $u:\mathfrak{X}\to\mathfrak{Y}$ est adique}~:
en effet, si $f:\mathfrak{X}\to\mathfrak{S}$,
$g:\mathfrak{Y}\to\mathfrak{S}$ sont les morphismes structuraux, et
$\mathscr{J}$ un Idéal de définition de $\mathfrak{S}$, on a
$f=g\circ u$, donc
$u^*(g^*(\mathscr{J})\mathscr{O}_{\mathfrak{Y}})
\mathscr{O}_{\mathfrak{X}}=f^*(\mathscr{J})
\mathscr{O}_{\mathfrak{X}}$ est un Idéal de définition de
$\mathfrak{X}$, et par hypothèse
$g^*(\mathscr{J})\mathscr{O}_{\mathfrak{Y}}$ est un Idéal de définition
de $\mathfrak{Y}$.

\begin{env}[10.12.2]
\label{I.10.12.2-fr}
Dans ce qui suit, nous supposerons fixé un préschéma formel localement
noethérien $\mathfrak{S}$ et un Idéal de définition $\mathscr{J}$ de
$\mathfrak{S}$~; nous poserons
$S_n=(\mathfrak{S},\mathscr{O}_{\mathfrak{S}}/\mathscr{J}^{n+1})$. Les
$\mathfrak{S}$-préschémas adiques (localement noethériens) forment
évidemment une \emph{catégorie}. Nous dirons qu'un système inductif
$(X_n)$ de $S_n$-préschémas (usuels) localement noethériens est un
\emph{$(S_n)$-système inductif adique} si les morphismes structuraux
$f_n:X_n\to S_n$ sont tels que, pour $m\leq n$, les diagrammes
\[
  \label{I.10.12.2.1-fr}
  \xymatrix{
    X_n \ar[d]_{f_n}
    & X_m \ar[l] \ar[d]^{f_m}\\
    S_n
    & S_m \ar[l]
  }
  \tag{10.12.2.1}
\]
soient commutatifs et \emph{identifient $X_m$ au produit
$X_n\times_{S_n}S_m=(X_n)_{(S_m)}$}. Les systèmes inductifs adiques
forment une \emph{catégorie}~: il suffit en effet de définir un morphisme
$(X_n)\to(Y_n)$ de tels systèmes comme un \emph{système inductif de
$S_n$-morphismes $u_n:X_n\to Y_n$} tel que $u_m$ s'identifie à
$(u_n)_{(S_m)}$ pour $m\leq n$. Cela étant~:
\end{env}

\begin{theorem}[10.12.3]
\label{I.10.12.3-fr}
Il y a équivalence canonique entre la catégorie des
$\mathfrak{S}$-préschémas adiques et la catégorie des
$(S_n)$-systèmes inductifs adiques.
\end{theorem}

L'équivalence en question s'obtient de la façon suivante~: si
$\mathfrak{X}$ est un $\mathfrak{S}$-préschéma adique,
$f:\mathfrak{X}\to\mathfrak{S}$ le morphisme structural,
$\mathscr{K}=f^*(\mathscr{J})\mathscr{O}_{\mathfrak{X}}$ est un Idéal de
définition de $\mathfrak{X}$ et on fait correspondre à $\mathfrak{X}$ le
système inductif des
$X_n=(\mathfrak{X},\mathscr{O}_{\mathfrak{X}}/\mathscr{K}^{n+1})$, le
morphisme structural $f_n:X_n\to S_n$ correspondant à $f$
(\hyperref[I.10.5.6-fr]{10.5.6}). Montrons d'abord que $(X_n)$ est un
\emph{système inductif adique}~: si $f=(\psi,\theta)$, on a
$\psi^*(\mathscr{J})\mathscr{O}_{\mathfrak{X}}=\mathscr{K}$, donc
$\psi^*(\mathscr{J}^n)\mathscr{O}_{\mathfrak{X}}=\mathscr{K}^n$ pour tout
$n$, et (par l'exactitude du foncteur $\psi^*$)
$\mathscr{K}^{m+1}/\mathscr{K}^{n+1}
=\psi^*(\mathscr{J}^{m+1}/\mathscr{J}^{n+1})
(\mathscr{O}_{\mathfrak{X}}/\mathscr{K}^{n+1})$ pour $m\leq n$~; notre
conclusion résulte donc de (\hyperref[I.4.4.5-fr]{4.4.5}). Il est immédiat
en outre de vérifier qu'à un $\mathfrak{S}$-morphisme
$u:\mathfrak{X}\to\mathfrak{Y}$ de $\mathfrak{S}$-préschémas adiques
correspond (avec des notations évidentes)

\oldpage[I]{207}

un système inductif de $S_n$-morphismes $u_n:X_n\to Y_n$ tel que $u_m$
s'identifie à $(u_n)_{(S_m)}$ pour $m\leq n$.

Le fait qu'on a bien défini ainsi une équivalence résultera de la proposition
plus précise suivante~:

\begin{proposition}[10.12.3.1]
\label{I.10.12.3.1-fr}
Soit $(X_n)$ un système inductif de $S_n$-préschémas~; on suppose que les
morphismes structuraux $f_n:X_n\to S_n$ sont tels que les diagrammes
(\hyperref[I.10.12.2.1-fr]{10.12.2.1}) soient commutatifs et identifient
$X_m$ à $X_n\times_{S_n}S_m$ pour $m\leq n$. Alors le système inductif
$(X_n)$ vérifie les conditions b) et c) de
(\hyperref[I.10.6.3-fr]{10.6.3})~; soient $\mathfrak{X}$ sa limite inductive,
$f:\mathfrak{X}\to\mathfrak{S}$ le morphisme limite inductive du système
inductif $(f_n)$. Alors, si $X_0$ est localement noethérien,
$\mathfrak{X}$ est localement noethérien et $f$ est un morphisme adique.
\end{proposition}

Comme le faisceau d'idéaux de $\mathscr{O}_{S_n}$ qui définit le
sous-préschéma $S_m$ de $S_n$ est nilpotent, il en est de même, en vertu de
(\hyperref[I.4.4.5-fr]{4.4.5}) du faisceau d'idéaux de
$\mathscr{O}_{X_n}$ définissant le sous-préschéma $X_m$ de $X_n$, donc les
conditions de (\hyperref[I.10.6.3-fr]{10.6.3}) sont bien vérifiées. La
question est par suite locale sur $\mathfrak{X}$ et $\mathfrak{S}$ et on
peut supposer que $\mathfrak{S}=\operatorname{Spf}(A)$,
$\mathscr{J}=\mathfrak{J}^\Delta$, $A$ étant un anneau $\mathfrak{J}$-adique
noethérien, et $X_n=\operatorname{Spec}(B_n)$~; si
$A_n=A/\mathfrak{J}^{n+1}$, l'hypothèse entraîne que $B_0$ est noethérien et
que si on pose $\mathfrak{J}_n=\mathfrak{J}/\mathfrak{J}^{n+1}$,
$B_m=B_n/\mathfrak{J}_n^{m+1}B_n$. Le noyau de $B_n\to B_0$ est donc
$\mathfrak{K}_n=\mathfrak{J}_nB_n$ et le noyau de $B_n\to B_m$ est
$\mathfrak{K}_n^{m+1}$ pour $m\leq n$~; en outre, comme $A_1$ est
noethérien, $\mathfrak{J}_1$ est de type fini sur $A_1$, donc
$\overline{\mathfrak{K}}_1=\mathfrak{K}_1/\mathfrak{K}_1^2$ est de type fini
sur $B_1$, et \emph{a fortiori} sur
$B_0=B_1/\mathfrak{K}_1$~; le fait que $\mathfrak{X}$ soit noethérien
résulte alors de (\hyperref[I.10.6.4-fr]{10.6.4})~; si
$B=\varprojlim B_n$, on a $\mathfrak{X}=\operatorname{Spf}(B)$, et si
$\mathfrak{K}$ est le noyau de $B\to B_0$,
$B_n=B/\mathfrak{K}^{n+1}$. Si
$\rho_n:A/\mathfrak{J}^{n+1}\to B/\mathfrak{K}^{n+1}$ est l'homomorphisme
correspondant à $f_n$, on a donc
\[
  \mathfrak{K}/\mathfrak{K}^{n+1}
  =(B/\mathfrak{K}^{n+1})\rho_n(\mathfrak{J}/\mathfrak{J}^{n+1})
\]
comme l'homomorphisme $\rho:A\to B$ correspondant à $f$ est égal à
$\varprojlim\rho_n$, l'idéal $\mathfrak{J}B$ de $B$ est dense dans
$\mathfrak{K}$, et comme tout idéal de $B$ est fermé
(\hyperref[0.7.3.5-fr]{0, 7.3.5}), on a $\mathfrak{K}=\mathfrak{J}B$. Si
$\mathscr{K}=\mathfrak{K}^\Delta$, la relation
$f^*(\mathscr{J})\mathscr{O}_{\mathfrak{X}}=\mathscr{K}$ résulte alors de
(\hyperref[I.10.10.9-fr]{10.10.9}) et achève la démonstration.

\begin{env}[10.12.3.2]
\label{I.10.12.3.2-fr}
L'équivalence précédente fournit, pour deux $\mathfrak{S}$-préschémas
adiques $\mathfrak{X}$, $\mathfrak{Y}$, une \emph{bijection canonique}
\[
  \operatorname{Hom}_{\mathfrak{S}}(\mathfrak{X},\mathfrak{Y})
  \xrightarrow{\sim}
  \varprojlim_n\operatorname{Hom}_{S_n}(X_n,Y_n)
\]
la limite projective étant relative aux applications
$u_n\to(u_n)_{(S_m)}$ pour $m\leq n$.
\end{env}

\subsection{Morphismes de type fini.}
\label{subsection:I.10.13-fr}

\begin{proposition}[10.13.1]
\label{I.10.13.1-fr}
Soient $\mathfrak{Y}$ un préschéma formel localement noethérien,
$\mathscr{K}$ un Idéal de définition de $\mathfrak{Y}$,
$f:\mathfrak{X}\to\mathfrak{Y}$ un morphisme de préschémas formels. Les
conditions suivantes sont équivalentes~:
\begin{enumerate}
  \item[a)] $\mathfrak{X}$ est localement noethérien, $f$ est un morphisme
    adique (\hyperref[I.10.12.1-fr]{10.12.1}) et si l'on pose
    $\mathscr{J}=f^*(\mathscr{K})\mathscr{O}_{\mathfrak{X}}$, le morphisme
    $f_0:(\mathfrak{X},\mathscr{O}_{\mathfrak{X}}/\mathscr{J})
    \to(\mathfrak{Y},\mathscr{O}_{\mathfrak{Y}}/\mathscr{K})$ déduit de
    $f$ est de type fini.
  \item[b)] $\mathfrak{X}$ est localement noethérien, et est limite
    inductive d'un $(Y_n)$-système inductif adique $(X_n)$ tel que le
    morphisme $X_0\to Y_0$ soit de type fini.

\oldpage[I]{208}

  \item[c)] Tout point de $\mathfrak{Y}$ possède un voisinage ouvert formel
    affine noethérien $V$ ayant la propriété suivante~:
    \begin{enumerate}
      \item[(\textbf{Q})] $f^{-1}(V)$ est réunion d'une famille finie
        d'ouverts formels affines noethériens $U_i$ tels que l'anneau adique
        noethérien $\Gamma(U_i,\mathscr{O}_{\mathfrak{X}})$ soit
        topologiquement isomorphe au quotient d'une algèbre de séries
        formelles restreintes (\hyperref[0.7.5.1-fr]{0, 7.5.1}) sur
        $\Gamma(V,\mathscr{O}_{\mathfrak{Y}})$, par un idéal (nécessairement
        fermé).
    \end{enumerate}
\end{enumerate}
\end{proposition}

Il est immédiat que a) entraîne b) en vertu de
(\hyperref[I.10.12.3-fr]{10.12.3}). Pour montrer que b) entraîne c), on
peut, puisque la question est locale sur $\mathfrak{Y}$, supposer que
$\mathfrak{Y}=\operatorname{Spf}(B)$, où $B$ est adique noethérien~; soit
$\mathscr{K}=\mathfrak{K}^\Delta$, $\mathfrak{K}$ étant un idéal de
définition de $B$. Comme par hypothèse, $X_0$ est de type fini sur $Y_0$,
$X_0$ est réunion finie d'ouverts affines $U_i$ tels que l'anneau $A_{i0}$
du schéma affine induit par $X_0$ sur $U_i$ soit une algèbre de type fini sur
l'anneau $B/\mathfrak{K}$ de $Y_0$
(\hyperref[I.6.3.2-fr]{6.3.2}). En vertu de
(\hyperref[I.5.1.9-fr]{5.1.9}), $U_i$ est aussi un ouvert affine dans chacun
des préschémas noethériens $X_n$, et si $A_{in}$ est l'anneau du schéma
affine induit par $X_n$ sur $U_i$, l'hypothèse b) entraîne que pour
$m\leq n$, $A_{im}$ est isomorphe à
$A_{in}/\mathfrak{K}^{m+1}A_{in}$. Par suite, le préschéma formel induit sur
$U_i$ par $\mathfrak{X}$ est isomorphe à $\operatorname{Spf}(A_i)$, où
$A_i=\varprojlim_n A_{in}$ (\hyperref[I.10.6.4-fr]{10.6.4})~; $A_i$ est un
anneau $\mathfrak{K}A_i$-adique, et $A_i/\mathfrak{K}A_i$, isomorphe à
$A_{i0}$, est une algèbre de type fini sur $B/\mathfrak{K}$. On en conclut
(\hyperref[0.7.5.5-fr]{0, 7.5.5}) que $A_i$ est topologiquement isomorphe à
un quotient d'une algèbre de séries formelles restreintes sur $B$ (par un
idéal nécessairement fermé, puisqu'une telle algèbre est noethérienne
(\hyperref[0.7.5.4-fr]{0, 7.5.4})).

Pour démontrer que c) entraîne a), on peut se limiter au cas où
$\mathfrak{X}=\operatorname{Spf}(A)$ est aussi affine, $A$ étant un anneau
adique noethérien, isomorphe au quotient d'une algèbre de séries formelles
restreintes sur $B$ par un idéal fermé. Alors
(\hyperref[0.7.5.5-fr]{0, 7.5.5}), $A/\mathfrak{K}A$ est une algèbre de type
fini sur $B/\mathfrak{K}$, et $\mathfrak{K}A=\mathfrak{J}$ est un idéal de
définition de $A$, donc, en vertu de
(\hyperref[I.10.10.9-fr]{10.10.9}), les conditions de a) sont satisfaites.

On notera que si les conditions de la prop.
(\hyperref[I.10.13.1-fr]{10.13.1}) sont remplies, la propriété a) est
valable pour \emph{tout Idéal de définition $\mathscr{K}$} de
$\mathfrak{Y}$ (en vertu de c)), et par suite, dans la propriété b),
\emph{tous les $f_n$ sont des morphismes de type fini}.

\begin{corollary}[10.13.2]
\label{I.10.13.2-fr}
Si les conditions de (\hyperref[I.10.13.1-fr]{10.13.1}) sont vérifiées,
tout ouvert formel affine noethérien $V$ de $\mathfrak{Y}$ possède la
propriété (\textbf{Q}) et si $\mathfrak{Y}$ est noethérien, il en est de
même de $\mathfrak{X}$.
\end{corollary}

Cela résulte aussitôt de (\hyperref[I.10.13.1-fr]{10.13.1}) et de
(\hyperref[I.6.3.2-fr]{6.3.2}).

\begin{definition}[10.13.3]
\label{I.10.13.3-fr}
Lorsque les propriétés équivalentes a), b), c) de
(\hyperref[I.10.13.1-fr]{10.13.1}) sont vérifiées, on dit que le morphisme
$f$ est \emph{de type fini}, ou que $\mathfrak{X}$ est un
\emph{$\mathfrak{Y}$-préschéma formel de type fini}, ou un
\emph{préschéma formel de type fini au-dessus de $\mathfrak{Y}$}.
\end{definition}

\begin{corollary}[10.13.4]
\label{I.10.13.4-fr}
Soient $\mathfrak{X}=\operatorname{Spf}(A)$,
$\mathfrak{Y}=\operatorname{Spf}(B)$ deux schémas formels affines
noethériens~; pour que $\mathfrak{X}$ soit de type fini sur
$\mathfrak{Y}$, il faut et il suffit que l'anneau adique noethérien $A$ soit
isomorphe au quotient d'une algèbre de séries formelles restreintes sur $B$
par un idéal fermé.
\end{corollary}

En effet, avec les notations de (\hyperref[I.10.13.1-fr]{10.13.1}), si
$\mathfrak{X}$ est de type fini sur $\mathfrak{Y}$,
$A/\mathfrak{K}A$ est alors une $(B/\mathfrak{K})$-algèbre de type fini en
vertu de (\hyperref[I.6.3.3-fr]{6.3.3}) et $\mathfrak{K}A$ est un idéal de
définition de $A$ (\hyperref[I.10.10.9-fr]{10.10.9}). On conclut donc par
(\hyperref[0.7.5.5-fr]{0, 7.5.5}).

\begin{proposition}[10.13.5]
\label{I.10.13.5-fr}
\medskip\noindent
\begin{enumerate}
  \item[\rm{(i)}] Le composé de deux morphismes de préschémas formels qui
    sont de type fini est de type fini.

\oldpage[I]{209}

  \item[\rm{(ii)}] Soient $\mathfrak{X}$, $\mathfrak{S}$,
    $\mathfrak{S}'$ trois préschémas formels localement noethériens (resp.
    noethériens), $f:\mathfrak{X}\to\mathfrak{S}$,
    $g:\mathfrak{S}'\to\mathfrak{S}$ deux morphismes. Si $f$ est de type
    fini, $\mathfrak{X}\times_{\mathfrak{S}}\mathfrak{S}'$ est localement
    noethérien (resp. noethérien) et est de type fini sur $\mathfrak{S}'$.
  \item[\rm{(iii)}] Soient $\mathfrak{S}$ un préschéma formel localement
    noethérien, $\mathfrak{X}'$, $\mathfrak{Y}'$ deux
    $\mathfrak{S}$-préschémas formels localement noethériens tels que
    $\mathfrak{X}'\times_{\mathfrak{S}}\mathfrak{Y}'$ soit localement
    noethérien. Si $\mathfrak{X}$, $\mathfrak{Y}$ sont des
    $\mathfrak{S}$-préschémas formels localement noethériens,
    $f:\mathfrak{X}\to\mathfrak{X}'$,
    $g:\mathfrak{Y}\to\mathfrak{Y}'$ deux $\mathfrak{S}$-morphismes de
    type fini, $\mathfrak{X}\times_{\mathfrak{S}}\mathfrak{Y}$ est
    localement noethérien et $f\times_{\mathfrak{S}}g$ est un
    $\mathfrak{S}$-morphisme de type fini.
\end{enumerate}
\end{proposition}

(iii) se déduit de (i) et (ii) par le raisonnement formel de
(\hyperref[I.3.5.1-fr]{3.5.1}) et il suffit donc de prouver (i) et (ii).

Soient $\mathfrak{X}$, $\mathfrak{Y}$, $\mathfrak{Z}$ trois préschémas
formels localement noethériens, $f:\mathfrak{X}\to\mathfrak{Y}$,
$g:\mathfrak{Y}\to\mathfrak{Z}$ deux morphismes de type fini. Si
$\mathscr{L}$ est un Idéal de définition de $\mathfrak{Z}$,
$\mathscr{K}=g^*(\mathscr{L})\mathscr{O}_{\mathfrak{Y}}$ en est un pour
$\mathfrak{Y}$ et
$\mathscr{J}=f^*(g^*(\mathscr{L}))\mathscr{O}_{\mathfrak{X}}$ en est un
pour $\mathfrak{X}$. Posons
$X_0=(\mathfrak{X},\mathscr{O}_{\mathfrak{X}}/\mathscr{J})$,
$Y_0=(\mathfrak{Y},\mathscr{O}_{\mathfrak{Y}}/\mathscr{K})$,
$Z_0=(\mathfrak{Z},\mathscr{O}_{\mathfrak{Z}}/\mathscr{L})$ et soient
$f_0:X_0\to Y_0$, $g_0:Y_0\to Z_0$ les morphismes correspondant à $f$ et
$g$. Comme par hypothèse $f_0$ et $g_0$ sont de type fini, il en est de même
de $g_0\circ f_0$ (\hyperref[I.6.3.4-fr]{6.3.4}) qui correspond à
$g\circ f$~; donc $g\circ f$ est de type fini par
(\hyperref[I.10.13.1-fr]{10.13.1}).

Sous les conditions de (ii), $\mathfrak{S}$ (resp. $\mathfrak{X}$,
$\mathfrak{S}'$) est limite inductive d'une suite $(S_n)$ (resp. $(X_n)$,
$(S'_n)$) de préschémas localement noethériens et on peut supposer
(\hyperref[I.10.13.1-fr]{10.13.1}) que
$X_m=X_n\times_{S_n}S_m$ pour $m\leq n$. Le préschéma formel
$\mathfrak{X}\times_{\mathfrak{S}}\mathfrak{S}'$ est alors limite inductive
des préschémas $X_n\times_{S_n}S'_n$
(\hyperref[I.10.7.4-fr]{10.7.4}), et on a
\[
  X_m\times_{S_m}S'_m
  =(X_n\times_{S_n}S_m)\times_{S_m}S'_m
  =(X_n\times_{S_n}S'_n)\times_{S'_n}S'_m.
\]
En outre, $X_0\times_{S_0}S'_0$ est localement noethérien puisque $X_0$
est de type fini sur $S_0$ (\hyperref[I.6.3.8-fr]{6.3.8}). On en conclut
d'abord (\hyperref[I.10.12.3.1-fr]{10.12.3.1}) que
$\mathfrak{X}\times_{\mathfrak{S}}\mathfrak{S}'$ est localement
noethérien~; en outre, comme $X_0\times_{S_0}S'_0$ est de type fini sur
$S'_0$ (\hyperref[I.6.3.8-fr]{6.3.8}), il résulte de
(\hyperref[I.10.12.3.1-fr]{10.12.3.1}) et de
(\hyperref[I.10.13.1-fr]{10.13.1}) que
$\mathfrak{X}\times_{\mathfrak{S}}\mathfrak{S}'$ est de type fini sur
$\mathfrak{S}'$, ce qui achève de prouver (ii) (l'assertion relative aux
préschémas noethériens étant conséquence immédiate de
(\hyperref[I.6.3.8-fr]{6.3.8})).

\begin{corollary}[10.13.6]
\label{I.10.13.6-fr}
Sous les hypothèses de (\hyperref[I.10.9.9-fr]{10.9.9}), si $f$ est un
morphisme de type fini, il en est de même de son prolongement $\widehat f$
aux complétés.
\end{corollary}

\subsection{Sous-préschémas fermés des préschémas formels.}
\label{subsection:I.10.14-fr}

\begin{proposition}[10.14.1]
\label{I.10.14.1-fr}
Soient $\mathfrak{X}$ un préschéma formel localement noethérien,
$\mathscr{A}$ un faisceau d'idéaux cohérent de
$\mathscr{O}_{\mathfrak{X}}$. Si $\mathfrak{Y}$ est le support (fermé) de
$\mathscr{O}_{\mathfrak{X}}/\mathscr{A}$, l'espace topologiquement annelé
$(\mathfrak{Y},(\mathscr{O}_{\mathfrak{X}}/\mathscr{A})|\mathfrak{Y})$ est
un préschéma formel localement noethérien, qui est noethérien si
$\mathfrak{X}$ l'est.
\end{proposition}

Notons que $\mathscr{O}_{\mathfrak{X}}/\mathscr{A}$ est cohérent en vertu de
(\hyperref[I.10.10.3-fr]{10.10.3}) et
(\hyperref[0.5.3.4-fr]{0, 5.3.4}), donc son support $\mathfrak{Y}$ est fermé
(\hyperref[0.5.2.2-fr]{0, 5.2.2}). Soit $\mathscr{J}$ un Idéal de définition
de $\mathfrak{X}$, et soit
$X_n=(\mathfrak{X},\mathscr{O}_{\mathfrak{X}}/\mathscr{J}^{n+1})$~; le
faisceau d'anneaux $\mathscr{O}_{\mathfrak{X}}/\mathscr{A}$ est limite
projective des faisceaux
$\mathscr{O}_{\mathfrak{X}}/(\mathscr{A}+\mathscr{J}^{n+1})
=(\mathscr{O}_{\mathfrak{X}}/\mathscr{A})
\otimes_{\mathscr{O}_{\mathfrak{X}}}
(\mathscr{O}_{\mathfrak{X}}/\mathscr{J}^{n+1})$
(\hyperref[I.10.11.3-fr]{10.11.3}), qui ont tous pour support
$\mathfrak{Y}$. Le faisceau
$(\mathscr{A}+\mathscr{J}^{n+1})/\mathscr{J}^{n+1}$ est un
$\mathscr{O}_{\mathfrak{X}}$-Module cohérent, puisque
$\mathscr{J}^{n+1}$ est cohérent, donc
$(\mathscr{A}+\mathscr{J}^{n+1})/\mathscr{J}^{n+1}$ est aussi un
$(\mathscr{O}_{\mathfrak{X}}/\mathscr{J}^{n+1})$-Module cohérent
(\hyperref[0.5.3.10-fr]{0, 5.3.10})~; si $Y_n$ est le sous-préschéma fermé
de $X_n$ défini par ce faisceau d'idéaux, il est immédiat que
$(\mathfrak{Y},(\mathscr{O}_{\mathfrak{X}}/\mathscr{A})|\mathfrak{Y})$

\oldpage[I]{210}
est le préschéma formel limite inductive des $Y_n$, et comme les conditions de
(\hyperref[I.10.6.4-fr]{10.6.4}) sont satisfaites, cela prouve que ce préschéma
formel est localement noethérien, et noethérien si $\mathfrak{X}$ l'est (puisque
alors $Y_0$ l'est en vertu de (\hyperref[I.6.1.4-fr]{6.1.4})).

\begin{definition}[10.14.2]
\label{I.10.14.2-fr}
On appelle \emph{sous-préschéma fermé} d'un préschéma formel $\mathfrak{X}$ tout
préschéma formel
$(\mathfrak{Y},(\mathscr{O}_{\mathfrak{X}}/\mathscr{A})|\mathfrak{Y})$ où
$\mathscr{A}$ est un $\mathscr{O}_{\mathfrak{X}}$-Module cohérent~; on dit que ce
préschéma est le sous-préschéma fermé défini par $\mathscr{A}$.
\end{definition}

Il est clair que la correspondance ainsi définie entre
$\mathscr{O}_{\mathfrak{X}}$-Modules cohérents et sous-préschémas fermés de
$\mathfrak{X}$ est biunivoque.

Le morphisme d'espaces topologiquement annelés
$j=(\psi,\theta):\mathfrak{Y}\to\mathfrak{X}$, où $\psi$ est l'injection
$\mathfrak{Y}\to\mathfrak{X}$ et $\theta^\sharp$ l'homomorphisme canonique
$\mathscr{O}_{\mathfrak{X}}\to\mathscr{O}_{\mathfrak{X}}/\mathscr{A}$, est
évidemment (\hyperref[I.10.4.5-fr]{10.4.5}) un morphisme de préschémas formels,
qu'on appelle l'\emph{injection canonique} de $\mathfrak{Y}$ dans $\mathfrak{X}$.
On notera que si $\mathfrak{X}=\operatorname{Spf}(A)$, où $A$ est adique
noethérien, on a $\mathscr{A}=\mathfrak{a}^\Delta$, où $\mathfrak{a}$ est un
idéal de $A$ (\hyperref[I.10.10.5-fr]{10.10.5}), et il résulte aussitôt de ce
qui précède que l'on a alors
$\mathfrak{Y}=\operatorname{Spf}(A/\mathfrak{a})$ à un isomorphisme près, et que
$j$ correspond (\hyperref[I.10.2.2-fr]{10.2.2}) à l'homomorphisme canonique
$A\to A/\mathfrak{a}$.

On dit qu'un morphisme $f:\mathfrak{Z}\to\mathfrak{X}$ de préschémas formels
localement noethériens est une \emph{immersion fermée} si elle se factorise en
$\mathfrak{Z}\xrightarrow{g}\mathfrak{Y}\xrightarrow{j}\mathfrak{X}$, où $g$
est un isomorphisme de $\mathfrak{Z}$ sur un sous-préschéma fermé
$\mathfrak{Y}$ de $\mathfrak{X}$ et $j$ l'injection canonique. Comme $j$ est un
monomorphisme d'espaces annelés, $g$ et $\mathfrak{Y}$ sont nécessairement
\emph{uniques}.

\begin{proposition}[10.14.3]
\label{I.10.14.3-fr}
Une immersion fermée est un morphisme de type fini.
\end{proposition}

On se ramène aussitôt au cas où $\mathfrak{X}$ est un schéma formel affine
$\operatorname{Spf}(A)$ et $\mathfrak{Y}=\operatorname{Spf}(A/\mathfrak{a})$~;
la proposition résulte de (\hyperref[I.10.13.1-fr]{10.13.1, c)}).

\begin{lemma}[10.14.4]
\label{I.10.14.4-fr}
Soit $f:\mathfrak{Y}\to\mathfrak{X}$ un morphisme de préschémas formels
localement noethériens, et soit $(U_\alpha)$ un recouvrement de
$f(\mathfrak{Y})$ par des ouverts formels affines noethériens de
$\mathfrak{X}$, tels que les $f^{-1}(U_\alpha)$ soient des ouverts formels
affines noethériens de $\mathfrak{Y}$. Pour que $f$ soit une immersion fermée,
il faut et il suffit que $f(\mathfrak{Y})$ soit une partie fermée de
$\mathfrak{X}$ et que, pour tout $\alpha$, la restriction de $f$ à
$f^{-1}(U_\alpha)$ corresponde (\hyperref[I.10.4.6-fr]{10.4.6}) à un
homomorphisme surjectif
$\Gamma(U_\alpha,\mathscr{O}_{\mathfrak{X}})\to
\Gamma(f^{-1}(U_\alpha),\mathscr{O}_{\mathfrak{Y}})$.
\end{lemma}

Les conditions sont évidemment nécessaires. Inversement, si elles sont
remplies, et si on désigne par $\mathfrak{a}_\alpha$ le noyau de
$\Gamma(U_\alpha,\mathscr{O}_{\mathfrak{X}})\to
\Gamma(f^{-1}(U_\alpha),\mathscr{O}_{\mathfrak{Y}})$, on définit un faisceau
cohérent d'idéaux $\mathscr{A}$ de $\mathscr{O}_{\mathfrak{X}}$ en prenant
$\mathscr{A}|U_\alpha=\mathfrak{a}_\alpha^\Delta$, et en prenant $\mathscr{A}$
nul dans le complémentaire de la réunion des $U_\alpha$. En effet, puisque
$f(\mathfrak{Y})$ est fermé et que le support de
$\mathfrak{a}_\alpha^\Delta$ est $U_\alpha\cap f(\mathfrak{Y})$, tout revient
à vérifier que $\mathfrak{a}_\alpha^\Delta$ et
$\mathfrak{a}_\beta^\Delta$ induisent le même faisceau sur un ouvert formel
affine noethérien $V\subset U_\alpha\cap U_\beta$. Or, la restriction de $f$
à $f^{-1}(U_\alpha)$ étant une immersion fermée de ce préschéma formel dans
$U_\alpha$, $f^{-1}(V)$ est un ouvert formel affine noethérien dans
$f^{-1}(U_\alpha)$ et la restriction de $f$ à $f^{-1}(V)$ est une immersion
fermée~; si $\mathfrak{b}$ est le noyau de l'homomorphisme surjectif
$\Gamma(V,\mathscr{O}_{\mathfrak{X}})\to
\Gamma(f^{-1}(V),\mathscr{O}_{\mathfrak{Y}})$ correspondant à cette
restriction, il est immédiat (\hyperref[I.10.10.2-fr]{10.10.2}) que
$\mathfrak{a}_\alpha^\Delta$ induit $\mathfrak{b}^\Delta$ sur $V$. Le faisceau
d'idéaux $\mathscr{A}$ étant ainsi défini, il est alors clair que
$f=g\circ j$, où $j:\mathfrak{Z}\to\mathfrak{X}$ est l'injection canonique du
sous-préschéma fermé $\mathfrak{Z}$ de $\mathfrak{X}$ défini par
$\mathscr{A}$, et $g$ un isomorphisme de $\mathfrak{Y}$ sur $\mathfrak{Z}$.

\begin{proposition}[10.14.5]
\label{I.10.14.5-fr}
\begin{enumerate}
  \item[\rm{(i)}] Si $f:\mathfrak{Z}\to\mathfrak{Y}$,
  $g:\mathfrak{Y}\to\mathfrak{X}$ sont des immersions fermées de préschémas
  formels localement noethériens, $g\circ f$ est une immersion fermée.

\oldpage[I]{211}
  \item[\rm{(ii)}] Soient $\mathfrak{X}$, $\mathfrak{Y}$, $\mathfrak{S}$ trois
  préschémas formels localement noethériens,
  $f:\mathfrak{X}\to\mathfrak{S}$ une immersion fermée,
  $g:\mathfrak{Y}\to\mathfrak{S}$ un morphisme. Alors le morphisme
  $\mathfrak{X}\times_{\mathfrak{S}}\mathfrak{Y}\to\mathfrak{Y}$ est une
  immersion fermée.
  \item[\rm{(iii)}] Soient $\mathfrak{S}$ un préschéma formel localement
  noethérien, $\mathfrak{X}'$, $\mathfrak{Y}'$ deux
  $\mathfrak{S}$-préschémas formels localement noethériens tels que
  $\mathfrak{X}'\times_{\mathfrak{S}}\mathfrak{Y}'$ soit localement
  noethérien. Si $\mathfrak{X}$, $\mathfrak{Y}$ sont des
  $\mathfrak{S}$-préschémas formels localement noethériens,
  $f:\mathfrak{X}\to\mathfrak{X}'$, $g:\mathfrak{Y}\to\mathfrak{Y}'$ deux
  $\mathfrak{S}$-morphismes qui sont des immersions fermées, alors
  $f\times_{\mathfrak{S}}g$ est une immersion fermée.
\end{enumerate}
\end{proposition}

En vertu de (\hyperref[I.3.5.1-fr]{3.5.1}), il suffit encore de prouver (i)
et (ii).

Pour démontrer (i), on peut supposer que $\mathfrak{Y}$ (resp.
$\mathfrak{Z}$) est un sous-préschéma fermé de $\mathfrak{X}$ (resp.
$\mathfrak{Y}$) défini par un faisceau cohérent $\mathscr{J}$ (resp.
$\mathscr{K}$) d'idéaux de $\mathscr{O}_{\mathfrak{X}}$ (resp.
$\mathscr{O}_{\mathfrak{Y}}$)~; si $\psi$ est l'injection
$\mathfrak{Y}\to\mathfrak{X}$ des espaces sous-jacents,
$\psi_*(\mathscr{K})$ est un faisceau cohérent d'idéaux de
$\psi_*(\mathscr{O}_{\mathfrak{Y}})
=\mathscr{O}_{\mathfrak{X}}/\mathscr{J}$
(\hyperref[0.5.3.12-fr]{0, 5.3.12}), donc aussi un
$\mathscr{O}_{\mathfrak{X}}$-Module cohérent
(\hyperref[0.5.3.10-fr]{0, 5.3.10})~; le noyau $\mathscr{K}_1$ de
$\mathscr{O}_{\mathfrak{X}}\to
(\mathscr{O}_{\mathfrak{X}}/\mathscr{J})/\psi_*(\mathscr{K})$ est donc un
faisceau cohérent d'idéaux de $\mathscr{O}_{\mathfrak{X}}$
(\hyperref[0.5.3.4-fr]{0, 5.3.4}), et
$\mathscr{O}_{\mathfrak{X}}/\mathscr{K}_1$ est isomorphe à
$\psi_*(\mathscr{O}_{\mathfrak{Y}}/\mathscr{K})$, ce qui prouve que
$\mathfrak{Z}$ est isomorphe à un sous-préschéma fermé de $\mathfrak{X}$.

Pour démontrer (ii), il est immédiat qu'on peut se limiter au cas où
$\mathfrak{S}=\operatorname{Spf}(A)$,
$\mathfrak{X}=\operatorname{Spf}(B)$,
$\mathfrak{Y}=\operatorname{Spf}(C)$, $A$ étant un anneau
$\mathfrak{J}$-adique noethérien, $B=A/\mathfrak{a}$, où $\mathfrak{a}$ est
un idéal de $A$, $C$ une $A$-algèbre topologique adique et noethérienne. Tout
revient à prouver que l'homomorphisme
$C\to C\widehat{\otimes}_A(A/\mathfrak{a})$ est \emph{surjectif}~: or,
$A/\mathfrak{a}$ est un $A$-module de type fini, et sa topologie est la
topologie $\mathfrak{J}$-adique~; il résulte alors de
(\hyperref[0.7.7.8-fr]{0, 7.7.8}) que
$C\widehat{\otimes}_A(A/\mathfrak{a})$ s'identifie à
$C\otimes_A(A/\mathfrak{a})=C/\mathfrak{a}C$, d'où notre assertion.

\begin{corollary}[10.14.6]
\label{I.10.14.6-fr}
Sous les hypothèses de (\hyperref[I.10.14.5-fr]{10.14.5, (ii)}), soient
$p:\mathfrak{X}\times_{\mathfrak{S}}\mathfrak{Y}\to\mathfrak{X}$,
$q:\mathfrak{X}\times_{\mathfrak{S}}\mathfrak{Y}\to\mathfrak{Y}$ les
projections, de sorte que le diagramme
\[
  \xymatrix{
    \mathfrak{X} \ar[d]_{f}
    & \mathfrak{X}\times_{\mathfrak{S}}\mathfrak{Y}
      \ar[l]_{p} \ar[d]^{q}\\
    \mathfrak{S}
    & \mathfrak{Y} \ar[l]^{g}
  }
\]
est commutatif. Pour tout $\mathscr{O}_{\mathfrak{X}}$-Module cohérent
$\mathscr{F}$, on a alors un isomorphisme canonique de
$\mathscr{O}_{\mathfrak{Y}}$-Modules
\[
  \label{I.10.14.6.1-fr}
  u:g^*(f_*(\mathscr{F}))
  \xrightarrow{\sim}
  q_*(p^*(\mathscr{F})).
  \tag{10.14.6.1}
\]
\end{corollary}

Pour définir un homomorphisme
$g^*(f_*(\mathscr{F}))\to q_*(p^*(\mathscr{F}))$, on sait qu'il revient au
même de définir un homomorphisme
$f_*(\mathscr{F})\to g_*(q_*(p^*(\mathscr{F})))
=f_*(p_*(p^*(\mathscr{F})))$
(\hyperref[0.4.4.3-fr]{0, 4.4.3})~: nous prendrons $u=f_*(\rho)$, où $\rho$
est l'homomorphisme canonique
$\mathscr{F}\to p_*(p^*(\mathscr{F}))$
(\hyperref[0.4.4.3-fr]{0, 4.4.3}). Pour voir que $u$ est un isomorphisme, on
se ramène aussitôt au cas où $\mathfrak{S}$, $\mathfrak{X}$,
$\mathfrak{Y}$ sont des spectres formels d'anneaux adiques noethériens $A$,
$B$, $C$, avec les conditions vues ci-dessus dans
(\hyperref[I.10.14.5-fr]{10.14.5, (ii)})~; on a alors
$\mathscr{F}=M^\Delta$, où $M$ est un $(A/\mathfrak{a})$-module de type fini
(\hyperref[I.10.10.5-fr]{10.10.5}), et les deux membres de
(\hyperref[I.10.14.6.1-fr]{10.14.6.1}) s'identifient respectivement, en vertu
de (\hyperref[I.10.10.8-fr]{10.10.8}), à
$(C\otimes_A M)^\Delta$ et à
$((C/\mathfrak{a}C)\otimes_{A/\mathfrak{a}}M)^\Delta$, d'où le corollaire,
puisque
$(C/\mathfrak{a}C)\otimes_{A/\mathfrak{a}}M
=(C\otimes_A(A/\mathfrak{a}))\otimes_{A/\mathfrak{a}}M$ s'identifie
canoniquement à $C\otimes_A M$.

\oldpage[I]{212}
\begin{corollary}[10.14.7]
\label{I.10.14.7-fr}
Soient $X$ un préschéma usuel localement noethérien, $Y$ un sous-préschéma
fermé de $X$, $j$ l'injection canonique $Y\to X$, $X'$ une partie fermée de
$X$ et $Y'=Y\cap X'$~; alors
$\widehat{j}:Y_{/Y'}\to X_{/X'}$ est une immersion fermée, et pour tout
$\mathscr{O}_Y$-Module cohérent $\mathscr{F}$, on a
\[
  \widehat{j}_*(\mathscr{F}_{/Y'})=(j_*(\mathscr{F}))_{/X'}.
\]
\end{corollary}

Comme $Y'=j^{-1}(X')$, il suffit d'utiliser
(\hyperref[I.10.9.9-fr]{10.9.9}) et d'appliquer
(\hyperref[I.10.14.5-fr]{10.14.5}) et
(\hyperref[I.10.14.6-fr]{10.14.6}).

\subsection{Préschémas formels séparés.}
\label{subsection:I.10.15-fr}

\begin{definition}[10.15.1]
\label{I.10.15.1-fr}
Soient $\mathfrak{S}$ un préschéma formel, $\mathfrak{X}$ un
$\mathfrak{S}$-préschéma formel,
$f:\mathfrak{X}\to\mathfrak{S}$ le morphisme structural. On appelle
morphisme diagonal
$\Delta_{\mathfrak{X}|\mathfrak{S}}:\mathfrak{X}\to
\mathfrak{X}\times_{\mathfrak{S}}\mathfrak{X}$ (noté aussi
$\Delta_{\mathfrak{X}}$) le morphisme
$(1_{\mathfrak{X}},1_{\mathfrak{X}})_{\mathfrak{S}}$. On dit que
$\mathfrak{X}$ est séparé sur $\mathfrak{S}$, ou un
$\mathfrak{S}$-schéma formel, ou que $f$ est un morphisme séparé, si l'image
par $\Delta_{\mathfrak{X}}$ de l'espace sous-jacent $\mathfrak{X}$ est une
partie fermée de l'espace sous-jacent à
$\mathfrak{X}\times_{\mathfrak{S}}\mathfrak{X}$. On dit qu'un préschéma
formel $\mathfrak{X}$ est séparé, ou est un schéma formel, s'il est séparé
sur $\mathbf{Z}$.
\end{definition}

\begin{proposition}[10.15.2]
\label{I.10.15.2-fr}
Supposons que les préschémas formels $\mathfrak{S}$, $\mathfrak{X}$ soient
respectivement limites inductives des suites $(S_n)$, $(X_n)$ de
préschémas usuels, et que le morphisme
$f:\mathfrak{X}\to\mathfrak{S}$ soit limite inductive d'une suite de
morphismes $f_n:X_n\to S_n$. Pour que $f$ soit séparé, il faut et il suffit
que le morphisme $f_0:X_0\to S_0$ le soit.
\end{proposition}

En effet, $\Delta_{\mathfrak{X}|\mathfrak{S}}$ est alors limite inductive de
la suite de morphismes $\Delta_{X_n|S_n}$
(\hyperref[I.10.7.4-fr]{10.7.4}), et l'image par
$\Delta_{\mathfrak{X}|\mathfrak{S}}$ de l'espace sous-jacent
$\mathfrak{X}$ (resp. l'espace sous-jacent
$\mathfrak{X}\times_{\mathfrak{S}}\mathfrak{X}$) est identique à l'image
par $\Delta_{X_0|S_0}$ de l'espace sous-jacent $X_0$ (resp. à l'espace
sous-jacent $X_0\times_{S_0}X_0$)~; d'où la conclusion.

\begin{proposition}[10.15.3]
\label{I.10.15.3-fr}
On suppose dans ce qui suit que tous les préschémas formels (resp.
morphismes de préschémas formels) considérés soient limites inductives de
suites de préschémas usuels (resp. de morphismes de préschémas usuels).
\begin{enumerate}
  \item[\rm{(i)}] Le composé de deux morphismes séparés est séparé.
  \item[\rm{(ii)}] Si $f:\mathfrak{X}\to\mathfrak{X}'$,
  $g:\mathfrak{Y}\to\mathfrak{Y}'$ sont deux $\mathfrak{S}$-morphismes
  séparés, $f\times_{\mathfrak{S}}g$ est séparé.
  \item[\rm{(iii)}] Si $f:\mathfrak{X}\to\mathfrak{Y}$ est un
  $\mathfrak{S}$-morphisme séparé, le $\mathfrak{S}'$-morphisme
  $f_{(\mathfrak{S}')}$ est séparé pour toute extension
  $\mathfrak{S}'\to\mathfrak{S}$ du préschéma formel de base.
  \item[\rm{(iv)}] Si le composé $g\circ f$ de deux morphismes est séparé,
  $f$ est séparé.
\end{enumerate}
\emph{(On sous-entend dans cet énoncé que si un même préschéma formel
$\mathfrak{Z}$ intervient plusieurs fois dans une même proposition, on le
considère comme limite inductive de la \emph{même} suite $(Z_n)$ de
préschémas usuels partout où il figure, et les morphismes de $\mathfrak{Z}$
dans un préschéma formel (resp. d'un préschéma formel dans $\mathfrak{Z}$)
comme limites inductives de morphismes des $Z_n$ dans des préschémas usuels
(resp. de préschémas usuels dans les $Z_n$)).}
\end{proposition}

Avec les notations de (\hyperref[I.10.15.2-fr]{10.15.2}), on a en effet
$(g\circ f)_0=g_0\circ f_0$, et
$(f\times_{\mathfrak{S}}g)_0=f_0\times_{S_0}g_0$~; les assertions de
(\hyperref[I.10.15.3-fr]{10.15.3}) sont alors conséquences immédiates de
(\hyperref[I.10.15.2-fr]{10.15.2}) et des assertions correspondantes de
(\hyperref[I.5.5.1-fr]{5.5.1}) pour les préschémas usuels.

Nous laissons au lecteur le soin d'énoncer pour la même espèce de
préschémas formels et de morphismes que dans
(\hyperref[I.10.15.3-fr]{10.15.3}) les propositions correspondant à
(\hyperref[I.5.5.5-fr]{5.5.5}),

\oldpage[I]{213}
(\hyperref[I.5.5.9-fr]{5.5.9}) et
(\hyperref[I.5.5.10-fr]{5.5.10}) (en y remplaçant «~ouvert affine~» par
«~ouvert formel affine vérifiant la condition b) de
(\hyperref[I.10.6.3-fr]{10.6.3})~»).

Un raisonnement analogue montre aussi que tout schéma formel affine
\emph{noethérien} est séparé, ce qui justifie la terminologie.

\begin{proposition}[10.15.4]
\label{I.10.15.4-fr}
Soient $\mathfrak{S}$ un préschéma formel localement noethérien,
$\mathfrak{X}$, $\mathfrak{Y}$ deux $\mathfrak{S}$-préschémas formels
localement noethériens, tels que $\mathfrak{X}$ ou $\mathfrak{Y}$ soit de
type fini sur $\mathfrak{S}$ (de sorte que
$\mathfrak{X}\times_{\mathfrak{S}}\mathfrak{Y}$ est localement noethérien)
et que $\mathfrak{Y}$ soit séparé sur $\mathfrak{S}$. Soit
$f:\mathfrak{X}\to\mathfrak{Y}$ un $\mathfrak{S}$-morphisme~; alors le
morphisme graphe
$\Gamma_f=(1_{\mathfrak{X}},f)_{\mathfrak{S}}:\mathfrak{X}\to
\mathfrak{X}\times_{\mathfrak{S}}\mathfrak{Y}$ est une immersion fermée.
\end{proposition}

On peut supposer que $\mathfrak{S}$ est limite inductive d'une suite
$(S_n)$ de préschémas localement noethériens, $\mathfrak{X}$ (resp.
$\mathfrak{Y}$) limite inductive d'une suite $(X_n)$ (resp. $(Y_n)$) de
$S_n$-préschémas, $f$ limite inductive d'une suite de $S_n$-morphismes
$f_n:X_n\to Y_n$~; alors
$\mathfrak{X}\times_{\mathfrak{S}}\mathfrak{Y}$ est limite inductive de la
suite $(X_n\times_{S_n}Y_n)$ et $\Gamma_f$ de la suite
$(\Gamma_{f_n})$ (\hyperref[I.10.7.4-fr]{10.7.4})~; par hypothèse, $Y_0$ est
séparé sur $S_0$ (\hyperref[I.10.15.2-fr]{10.15.2}), donc l'espace
$\Gamma_{f_0}(X_0)$ est un sous-espace fermé de
$X_0\times_{S_0}Y_0$~; comme les espaces sous-jacents de
$\mathfrak{X}\times_{\mathfrak{S}}\mathfrak{Y}$ (resp.
$\Gamma_f(\mathfrak{X})$) et $X_0\times_{S_0}Y_0$ (resp.
$\Gamma_{f_0}(X_0)$) sont les mêmes, on voit déjà que
$\Gamma_f(\mathfrak{X})$ est un sous-espace \emph{fermé} de
$\mathfrak{X}\times_{\mathfrak{S}}\mathfrak{Y}$. Remarquons maintenant que
lorsque $(U,V)$ parcourt l'ensemble des couples formés d'un ouvert formel
affine noethérien $U$ (resp. $V$) de $\mathfrak{X}$ (resp.
$\mathfrak{Y}$) tels que $f(U)\subset V$, les ouverts
$U\times_{\mathfrak{S}}V$ forment un recouvrement de
$\Gamma_f(\mathfrak{X})$ dans
$\mathfrak{X}\times_{\mathfrak{S}}\mathfrak{Y}$, et si
$f_U:U\to V$ est la restriction de $f$ à $U$,
$\Gamma_{f_U}:U\to U\times_{\mathfrak{S}}V$ est la restriction de
$\Gamma_f$ à $U$. Si nous montrons que $\Gamma_{f_U}$ est une immersion
fermée, il en sera donc de même de $\Gamma_f$
(\hyperref[I.10.14.4-fr]{10.14.4}), autrement dit, on est ramené au cas où
$\mathfrak{S}=\operatorname{Spf}(A)$,
$\mathfrak{X}=\operatorname{Spf}(B)$,
$\mathfrak{Y}=\operatorname{Spf}(C)$ sont affines ($A$, $B$, $C$ adiques
noethériens), $f$ correspondant à un $A$-homomorphisme continu
$\varphi:C\to B$~; $\Gamma_f$ correspond alors à l'unique homomorphisme
continu $\omega:B\widehat{\otimes}_A C\to B$ qui, composé avec les
homomorphismes canoniques $B\to B\widehat{\otimes}_A C$ et
$C\to B\widehat{\otimes}_A C$, donne respectivement l'identité et
$\varphi$. Or, il est clair que $\omega$ est \emph{surjectif}, d'où notre
assertion.

\begin{corollary}[10.15.5]
\label{I.10.15.5-fr}
Soient $\mathfrak{S}$ un préschéma formel localement noethérien,
$\mathfrak{X}$ un $\mathfrak{S}$-préschéma de type fini~; pour que
$\mathfrak{X}$ soit séparé sur $\mathfrak{S}$, il faut et il suffit que le
morphisme diagonal
$\mathfrak{X}\to\mathfrak{X}\times_{\mathfrak{S}}\mathfrak{X}$ soit une
immersion fermée.
\end{corollary}

\begin{proposition}[10.15.6]
\label{I.10.15.6-fr}
Une immersion fermée $j:\mathfrak{Y}\to\mathfrak{X}$ de préschémas formels
localement noethériens est un morphisme séparé.
\end{proposition}

Avec les notations de (\hyperref[I.10.14.2-fr]{10.14.2}),
$j_0:Y_0\to X_0$ est une immersion fermée, donc un morphisme séparé, et il
suffit d'appliquer (\hyperref[I.10.15.2-fr]{10.15.2}).

\begin{proposition}[10.15.7]
\label{I.10.15.7-fr}
Soient $X$ un préschéma (usuel) localement noethérien, $X'$ une partie
fermée de $X$ et $\widehat{X}=X_{/X'}$. Pour que $\widehat{X}$ soit séparé,
il faut et il suffit que $\widehat{X}_{\mathrm{red}}$ le soit, et il suffit
que $X$ le soit.
\end{proposition}

En effet, avec les notations de (\hyperref[I.10.8.5-fr]{10.8.5}), pour que
$\widehat{X}$ soit séparé, il faut et il suffit que $X'_0$ le soit
(\hyperref[I.10.15.2-fr]{10.15.2}), et comme
$\widehat{X}_{\mathrm{red}}=(X'_0)_{\mathrm{red}}$, il est équivalent de
dire que $\widehat{X}_{\mathrm{red}}$ l'est
(\hyperref[I.5.5.1-fr]{5.5.1, (vi)}).
